% Section file for Chapter: Twisted Holography and Quantum Gravity
% Result (G7): Modular Bootstrap as MC Recursion
%
% Cross-references:
% thm:mc2-bar-intrinsic, thm:recursive-existence,
% thm:convolution-dg-lie-structure, const:vol1-genus-spectral-sequence,
% def:modular-convolution-dg-lie

% =====================================
% Local macros (provided only if absent from the ambient manuscript)
% =====================================
\providecommand{\cA}{\mathcal{A}}
\providecommand{\cH}{\mathcal{H}}
\providecommand{\cM}{\mathcal{M}}
\providecommand{\cC}{\mathcal{C}}
\providecommand{\fg}{\mathfrak{g}}
\providecommand{\fC}{\mathfrak{C}}
\providecommand{\gAmod}{\mathfrak{g}_{\cA}^{\mathrm{mod}}}
\providecommand{\dzero}{d_0}
\providecommand{\dfib}{d_{\mathrm{fib}}}
\providecommand{\MC}{\text{MC}}
\providecommand{\barB}{\overline{B}}
\providecommand{\Defcyc}{\operatorname{Def}_{\mathrm{cyc}}}
\providecommand{\Definfmod}{\operatorname{Def}_{\infty}^{\mathrm{mod}}}
\providecommand{\Convstr}{\operatorname{Conv}_{\mathrm{str}}}
\providecommand{\Convinf}{\operatorname{Conv}_{\infty}}
\providecommand{\Gmod}{\mathcal{G}^{\mathrm{mod}}}
\providecommand{\Tr}{\operatorname{Tr}}
\providecommand{\Hom}{\operatorname{Hom}}
\providecommand{\Res}{\operatorname{Res}}
\providecommand{\id}{\mathrm{id}}
\providecommand{\ad}{\operatorname{ad}}
\providecommand{\gr}{\operatorname{gr}}
\providecommand{\wt}{\operatorname{wt}}
\providecommand{\Aut}{\operatorname{Aut}}
\providecommand{\Det}{\operatorname{Det}}
\providecommand{\Sym}{\operatorname{Sym}}
\providecommand{\Omegach}{\Omega^{\text{ch}}}

% Bootstrap-specific macros
\providecommand{\Ob}{\mathrm{Ob}}
\providecommand{\Shad}{\operatorname{Sh}}
\providecommand{\hCont}{\mathsf{h}}
\providecommand{\BVop}{\Delta_{\mathrm{BV}}}
\providecommand{\dsew}{d_{\mathrm{sew}}}
\providecommand{\dpf}{d_{\mathrm{pf}}}

\section{Modular bootstrap as MC recursion}
% label removed: sec:thqg-modular-bootstrap
\index{modular bootstrap|textbf}
\index{genus spectral sequence!bootstrap}
\index{Maurer--Cartan equation!genus expansion}
\index{bootstrap!modular, from MC recursion}

The genus-$g$ partition function is determined inductively from
genus-$0$ data by the Maurer--Cartan equation in the modular
convolution algebra~$\gAmod$
(Definition~\ref*{V1-def:modular-convolution-dg-lie}):

\begin{center}
\emph{The MC equation $D_\cA\,\Theta_\cA
+ \tfrac{1}{2}[\Theta_\cA,\Theta_\cA] = 0$
is a modular bootstrap equation. Expanded genus by genus,
it determines $\Theta^{(g)}_\cA$ from $\Theta^{(<g)}_\cA$
by solving a linear equation with an explicitly computable
inhomogeneity, the genus-$g$ obstruction class $\Ob_g$.}
\end{center}

The conformal bootstrap proceeds by consistency conditions
(crossing symmetry, modular invariance) that constrain but
do not uniquely determine the correlators. The MC recursion
is \emph{algebraically complete}: it determines the partition
function at each genus from the OPE data at genus~$0$, with
no ambiguity beyond gauge equivalence. The mechanism is the
contracting homotopy of the tree-twisted differential
$d_{\Theta^{(0)}} = \dzero + [\Theta^{(0)}, -]$ on the
bar complex.

% ======================================================================
%
% PART 1: THE MC EQUATION EXPANDED BY GENUS
%
% ======================================================================

\subsection{The MC equation expanded by genus}
% label removed: subsec:thqg-VII-mc-genus-expansion
\index{Maurer--Cartan equation!genus expansion|textbf}

\subsubsection{The genus filtration on the convolution algebra}
% label removed: subsubsec:thqg-VII-genus-filtration

We recall from Definition~\ref*{V1-def:modular-convolution-dg-lie}
that the modular convolution dg~Lie algebra
\begin{equation}% label removed: eq:thqg-VII-gAmod-def
\gAmod
\;=\;
\prod_{g \geq 0}\;
\prod_{n \geq 1}\;
\Hom_{S_n}\!\bigl(
C_*(\overline{\cM}_{g,n}),\;
\cA^{\otimes n}
\bigr)^{S_n}[-1]
\end{equation}
carries a descending genus filtration
\begin{equation}% label removed: eq:thqg-VII-genus-filtration
F^g\,\gAmod
\;:=\;
\prod_{g' \geq g}\;
\prod_{n \geq 1}\;
\Hom_{S_n}\!\bigl(
C_*(\overline{\cM}_{g',n}),\;
\cA^{\otimes n}
\bigr)^{S_n}[-1].
\end{equation}
This filtration is complete and Hausdorff:
$\gAmod = F^0\,\gAmod$ and
$\bigcap_{g \geq 0} F^g\,\gAmod = 0$.
The associated graded is
\begin{equation}% label removed: eq:thqg-VII-assoc-graded
\gr^g_F\,\gAmod
\;=\;
F^g / F^{g+1}
\;\cong\;
\prod_{n \geq 1}\;
\Hom_{S_n}\!\bigl(
C_*(\overline{\cM}_{g,n}),\;
\cA^{\otimes n}
\bigr)^{S_n}[-1],
\end{equation}
the space of genus-$g$ operations.

\begin{remark}[Completeness and $\hbar$-adic topology]
\label{rem:thqg-VII-hbar-topology}
\index{hbar-adic topology@$\hbar$-adic topology}
The genus filtration is equivalent to the $\hbar$-adic filtration
when one introduces the formal genus-counting parameter~$\hbar$
and writes $\gAmod[\![\hbar]\!]$ with $\hbar^g$ multiplying the
genus-$g$ sector. The $\hbar$-adic topology on
$\gAmod[\![\hbar]\!]$ is the linear topology defined by the
ideals $\hbar^g \cdot \gAmod[\![\hbar]\!]$, and this coincides
with the topology defined by~\eqref{V1-eq:thqg-VII-genus-filtration}.
Completeness holds because the product in~%
\eqref{V1-eq:thqg-VII-gAmod-def} is an infinite product, not
a direct sum.
\end{remark}

\subsubsection{Decomposition of the differential}
% label removed: subsubsec:thqg-VII-differential-decomp

The differential $D$ on $\gAmod$ decomposes as the sum of
five components
(Theorem~\ref{V1-thm:convolution-dg-lie-structure}):
\begin{equation}% label removed: eq:thqg-VII-D-components
D \;=\; d_{\mathrm{Ch}_\infty}
 + d_{\check{C}}
 + d_{\mathrm{coll}}
 + d_{\mathrm{loop}}
 + d_{\mathrm{pf}},
\end{equation}
where:
\begin{enumerate}[label=(\roman*)]
\item $d_{\mathrm{Ch}_\infty}$: the homotopy chiral algebra
 structure maps (internal to each vertex of the graph);
\item $d_{\check{C}}$: the \v{C}ech differential (from the
 boundary stratification of $\overline{\cM}_{g,n}$, encoding
 separating degenerations);
\item $d_{\mathrm{coll}}$: the collision differential (residues
 along the collision divisors in the Fulton--MacPherson
 compactification);
\item $d_{\mathrm{loop}}$: the non-separating clutching operator
 (the BV operator $\Delta$ that raises genus by one, dual to
 the non-separating node of the boundary);
\item $d_{\mathrm{pf}}$: the planted-forest corrections
 (from the log-FM boundary strata of
 Mok~\cite{Mok25}).
\end{enumerate}
With respect to the genus filtration, these have the following
behaviour:

\begin{lemma}[Genus shifts of the differential components]
% label removed: lem:thqg-VII-genus-shifts
\ClaimStatusProvedElsewhere
\begin{enumerate}[label=\textup{(\roman*)}]
\item $d_{\mathrm{Ch}_\infty}$, $d_{\check{C}}$, and
 $d_{\mathrm{coll}}$ \emph{preserve} genus:
 $d_{\mathrm{Ch}_\infty}(F^g) \subset F^g$,
 and similarly for $d_{\check{C}}$ and $d_{\mathrm{coll}}$.
\item $d_{\mathrm{loop}}$ \emph{raises} genus by one:
 $d_{\mathrm{loop}}(F^g) \subset F^{g+1}$.
\item $d_{\mathrm{pf}}$ \emph{raises} genus by at least one:
 $d_{\mathrm{pf}}(F^g) \subset F^{g+1}$.
\end{enumerate}
\end{lemma}

\begin{proof}
The non-separating clutching map
$\overline{\cM}_{g-1,n+2} \to \partial\overline{\cM}_{g,n}$
identifies two marked points on a genus-$(g{-}1)$ surface to
produce a genus-$g$ surface. The dual operation on the
convolution algebra maps genus-$g$ elements to genus-$(g{+}1)$
elements. The components $d_{\mathrm{Ch}_\infty}$,
$d_{\check{C}}$, and $d_{\mathrm{coll}}$ all act fibrewise over
a fixed moduli space $\overline{\cM}_{g,n}$ and hence preserve
the genus. The planted-forest component $d_{\mathrm{pf}}$
involves boundary strata indexed by planted forests of positive
loop genus (each internal edge raising the genus), hence also
raises genus.

The genus filtration satisfies the multiplicative axiom
$F^g \cdot F^{g'} \subset F^{g+g'}$: the Lie bracket
$[\varphi, \psi]$ for $\varphi \in F^g$ and $\psi \in F^{g'}$
involves stable graphs of genus $\geq g + g'$, because the
bracket glues a genus-$g$ and a genus-$g'$ graph along a
single edge, producing a graph of genus $\geq g + g'$.
This multiplicative property ensures that the genus filtration
is compatible with the dg~Lie algebra structure.
\end{proof}

\begin{definition}[Genus-preserving and genus-raising parts]
% label removed: def:thqg-VII-D-split
\index{differential!genus-preserving part}
\index{differential!genus-raising part}
Define the genus-preserving part of $D$:
\begin{equation}% label removed: eq:thqg-VII-D0
D_0
\;:=\;
d_{\mathrm{Ch}_\infty}
+ d_{\check{C}}
+ d_{\mathrm{coll}},
\end{equation}
and the genus-raising part:
\begin{equation}% label removed: eq:thqg-VII-D1
D_1
\;:=\;
d_{\mathrm{loop}}
+ d_{\mathrm{pf}}.
\end{equation}
Then $D = D_0 + D_1$ with $D_0^2 = 0$,
$D_0 D_1 + D_1 D_0 = 0$, $D_1^2 = 0$
(these follow from $D^2 = 0$ by genus bookkeeping).
\end{definition}

\begin{remark}[The three identities]
% label removed: rem:thqg-VII-three-identities
The factorization $D^2 = 0$ into three homogeneous pieces is the
structural core of the genus spectral sequence. Explicitly:
\begin{align}
D^2 &= (D_0 + D_1)^2 \notag\\
&= D_0^2 + (D_0 D_1 + D_1 D_0) + D_1^2 \notag\\
&= 0,% label removed: eq:thqg-VII-D-squared-decomp
\end{align}
and each of the three terms vanishes separately because they
have different genus shifts: $D_0^2$ preserves genus,
$D_0 D_1 + D_1 D_0$ shifts by one, $D_1^2$ shifts by two.
Since these map into disjoint genus sectors of $\gAmod$,
each must vanish independently.
\end{remark}

\subsubsection{The genus expansion of $\Theta_\cA$}
% label removed: subsubsec:thqg-VII-genus-expansion

The universal MC element $\Theta_\cA \in \MC(\gAmod)$
(Theorem~\ref*{V1-thm:mc2-bar-intrinsic}) admits a genus expansion
\begin{equation}% label removed: eq:thqg-VII-theta-genus-exp
\Theta_\cA
\;=\;
\sum_{g \geq 0} \hbar^g\,\Theta^{(g)}_\cA,
\end{equation}
where $\Theta^{(g)}_\cA \in \gr^g_F\,\gAmod$ is the genus-$g$
component. The MC equation reads
\begin{equation}% label removed: eq:thqg-VII-mc-full
D\,\Theta_\cA
+ \tfrac{1}{2}[\Theta_\cA,\, \Theta_\cA]
\;=\; 0.
\end{equation}
We expand both sides in powers of~$\hbar$.

\begin{proposition}[Genus-$g$ MC equation; \ClaimStatusProvedHere]
% label removed: prop:thqg-VII-genus-g-mc
\index{Maurer--Cartan equation!genus-$g$ component}
The coefficient of $\hbar^g$ in the MC
equation~\eqref{V1-eq:thqg-VII-mc-full} is:
\begin{equation}% label removed: eq:thqg-VII-genus-g-mc-eq
\boxed{
d_{\Theta^{(0)}}\bigl(\Theta^{(g)}\bigr)
\;=\;
-\Ob_g\bigl(\Theta^{(<g)}\bigr),
\qquad g \geq 1,
}
\end{equation}
where the tree-twisted differential is
\begin{equation}% label removed: eq:thqg-VII-d-twisted
d_{\Theta^{(0)}}
\;:=\;
D_0 + \ad_{\Theta^{(0)}}
\;=\;
D_0 + [\Theta^{(0)},\, -],
\end{equation}
and the genus-$g$ obstruction class is
\begin{equation}% label removed: eq:thqg-VII-obs-class
\Ob_g\bigl(\Theta^{(<g)}\bigr)
\;:=\;
\underbrace{
\sum_{\substack{g_1 + g_2 = g \\ g_1, g_2 \geq 1}}
\tfrac{1}{2}\,[\Theta^{(g_1)},\, \Theta^{(g_2)}]
}_{\text{separating clutching}}
\;+\;
\underbrace{
\Delta\bigl(\Theta^{(g-1)}\bigr)
}_{\text{non-separating BV}}
\;+\;
\underbrace{
\dsew^{(g)} + \dpf^{(g)}
}_{\text{boundary strata}}.
\end{equation}
\end{proposition}

\begin{proof}
Write $\Theta = \sum_{g \geq 0} \hbar^g\,\Theta^{(g)}$ and
expand the MC equation~\eqref{V1-eq:thqg-VII-mc-full}.

\medskip\noindent\emph{Step 1: Expansion of $D\,\Theta$.}
The differential $D = D_0 + D_1$ acts on $\Theta^{(g)}$ as
\begin{align}
(D\,\Theta)_g
&= D_0(\Theta^{(g)}) + D_1(\Theta^{(g-1)})\notag\\
&= D_0(\Theta^{(g)})
 + d_{\mathrm{loop}}(\Theta^{(g-1)})
 + d_{\mathrm{pf}}(\Theta^{(g-1)}),
% label removed: eq:thqg-VII-D-theta-expand
\end{align}
since $D_0$ preserves genus while $D_1$ raises it by one.
The component $D_1(\Theta^{(g-1)})$ contributes to genus~$g$
by the non-separating clutching $\Delta(\Theta^{(g-1)})$ and
the planted-forest correction $\dpf^{(g)}$.

\medskip\noindent\emph{Step 2: Expansion of $[\Theta,\Theta]$.}
The bracket $[\Theta,\Theta]$ expanded in $\hbar$ gives
\begin{equation}% label removed: eq:thqg-VII-bracket-expand
\bigl[\Theta,\Theta\bigr]_g
\;=\;
\sum_{g_1 + g_2 = g}
[\Theta^{(g_1)},\, \Theta^{(g_2)}]
\;=\;
2\,[\Theta^{(0)},\, \Theta^{(g)}]
\;+\;
\sum_{\substack{g_1 + g_2 = g \\ g_1, g_2 \geq 1}}
[\Theta^{(g_1)},\, \Theta^{(g_2)}],
\end{equation}
where we have isolated the two terms with $g_1 = 0$ or
$g_2 = 0$ (which contribute $[\Theta^{(0)}, \Theta^{(g)}]$
each, and by the graded antisymmetry of the Lie bracket,
$[\Theta^{(0)}, \Theta^{(g)}] = -(-1)^{|\Theta^{(0)}|
|\Theta^{(g)}|}[\Theta^{(g)}, \Theta^{(0)}]$; since both
are degree~$1$, this gives $[\Theta^{(0)}, \Theta^{(g)}]
= -[\Theta^{(g)}, \Theta^{(0)}]$ and the two terms with
$g_1 = 0$ or $g_2 = 0$ contribute
$2[\Theta^{(0)}, \Theta^{(g)}]$).

\medskip\noindent\emph{Step 3: Assembly.}
Collecting the genus-$g$ coefficient of the MC equation
$D\,\Theta + \tfrac{1}{2}[\Theta,\Theta] = 0$:
\begin{align}
0 &= D_0(\Theta^{(g)})
 + \Delta(\Theta^{(g-1)})
 + \dsew^{(g)} + \dpf^{(g)}
 \notag\\
 &\qquad
 + [\Theta^{(0)},\, \Theta^{(g)}]
 + \tfrac{1}{2}
 \sum_{\substack{g_1 + g_2 = g \\ g_1, g_2 \geq 1}}
 [\Theta^{(g_1)},\, \Theta^{(g_2)}].
% label removed: eq:thqg-VII-mc-g-assembled
\end{align}
Rearranging, with $d_{\Theta^{(0)}} = D_0 +
[\Theta^{(0)}, -]$:
\begin{equation}% label removed: eq:thqg-VII-mc-g-rearranged
d_{\Theta^{(0)}}(\Theta^{(g)})
\;=\;
-\Delta(\Theta^{(g-1)})
- \dsew^{(g)} - \dpf^{(g)}
- \tfrac{1}{2}
 \sum_{\substack{g_1 + g_2 = g \\ g_1, g_2 \geq 1}}
 [\Theta^{(g_1)},\, \Theta^{(g_2)}],
\end{equation}
which is exactly $d_{\Theta^{(0)}}(\Theta^{(g)}) = -\Ob_g(\Theta^{(<g)})$.

\medskip\noindent\emph{Step 4: The tree-twisted differential squares to zero.}
We verify $d_{\Theta^{(0)}}^2 = 0$. The genus-$0$ MC equation gives
\begin{equation}% label removed: eq:thqg-VII-genus0-mc
D_0(\Theta^{(0)}) + \tfrac{1}{2}[\Theta^{(0)},\Theta^{(0)}] = 0.
\end{equation}
Write $d_{\Theta^{(0)}}(x) = D_0(x) + [\Theta^{(0)}, x]$.
By the Leibniz rule for a degree-$1$ derivation acting on the bracket:
\begin{equation}\label{eq:thqg-VII-leibniz-explicit}
D_0[\Theta^{(0)}, x]
\;=\;
[D_0\Theta^{(0)}, x]
+ (-1)^{|\Theta^{(0)}|}[\Theta^{(0)}, D_0(x)]
\;=\;
[D_0\Theta^{(0)}, x]
- [\Theta^{(0)}, D_0(x)],
\end{equation}
since $|\Theta^{(0)}| = 1$ (odd degree).
Then:
\begin{align}
d_{\Theta^{(0)}}^2(x)
&= (D_0 + \ad_{\Theta^{(0)}})^2(x)\notag\\
&= D_0^2(x) + D_0[\Theta^{(0)},x]
 + [\Theta^{(0)}, D_0(x)]
 + [\Theta^{(0)},[\Theta^{(0)},x]]\notag\\
&= 0
 + \bigl([D_0\Theta^{(0)}, x]
 - [\Theta^{(0)}, D_0(x)]\bigr)
 + [\Theta^{(0)}, D_0(x)]
 + [\Theta^{(0)},[\Theta^{(0)},x]]\notag\\
&= [D_0\Theta^{(0)}, x]
 + [\Theta^{(0)},[\Theta^{(0)},x]]\notag\\
&= [D_0\Theta^{(0)}, x]
 + \tfrac{1}{2}[[\Theta^{(0)},\Theta^{(0)}], x]\notag\\
&= [D_0\Theta^{(0)}
 + \tfrac{1}{2}[\Theta^{(0)},\Theta^{(0)}], x]\notag\\
&= 0,
% label removed: eq:thqg-VII-d-twisted-squared
\end{align}
where we used: $D_0^2 = 0$ in the first term; the Leibniz
rule~\eqref{eq:thqg-VII-leibniz-explicit} expanding
$D_0[\Theta^{(0)},x]$, after which the
$-[\Theta^{(0)}, D_0(x)]$ cancels the $+[\Theta^{(0)}, D_0(x)]$
from the cross term; the graded Jacobi identity
$[\Theta^{(0)},[\Theta^{(0)},x]] = \tfrac{1}{2}
[[\Theta^{(0)},\Theta^{(0)}], x]$
(valid because $|\Theta^{(0)}| = 1$ is odd); and
\eqref{V1-eq:thqg-VII-genus0-mc} for the final vanishing.
\end{proof}

\begin{remark}[Disjointness of the three terms]
% label removed: rem:thqg-VII-three-term-disjointness
\index{obstruction class!three-term decomposition}
The three terms in~\eqref{V1-eq:thqg-VII-obs-class} are disjoint
because they correspond to different codimension-$1$ boundary strata
of $\overline{\cM}_{g,n}$: separating nodes (clutching two
lower-genus surfaces, contributing
$\sum_{g_1+g_2=g} \tfrac{1}{2}[\Theta^{(g_1)},\Theta^{(g_2)}]$),
non-separating nodes (pinching a non-separating cycle on a single
surface, contributing $\Delta(\Theta^{(g-1)})$), and boundary
corrections from the log-FM strata (no topology change in the
surface but corrections from the planted-forest compactification,
contributing $\dsew^{(g)} + \dpf^{(g)}$). These are distinct
codimension-$1$ strata with no overlap: a degeneration is either
separating, non-separating, or a log-FM correction, never more
than one simultaneously.
\end{remark}

\begin{remark}[The obstruction class is a cocycle]
% label removed: rem:thqg-VII-obs-cocycle
\index{obstruction class!cocycle property}
The right side of~\eqref{V1-eq:thqg-VII-genus-g-mc-eq} is
$d_{\Theta^{(0)}}$-closed. Indeed, applying
$d_{\Theta^{(0)}}$ to~\eqref{V1-eq:thqg-VII-genus-g-mc-eq}:
\begin{equation}% label removed: eq:thqg-VII-obs-closed
d_{\Theta^{(0)}}\bigl(\Ob_g(\Theta^{(<g)})\bigr)
\;=\;
d_{\Theta^{(0)}}^2(\Theta^{(g)})
\;=\; 0,
\end{equation}
so $\Ob_g(\Theta^{(<g)})$ defines a cohomology class
\begin{equation}% label removed: eq:thqg-VII-obs-class-coh
[\Ob_g]
\;\in\;
H^2\bigl(\gAmod[\text{genus }g],\; d_{\Theta^{(0)}}\bigr).
\end{equation}
The solvability of the recursion at genus~$g$ is equivalent to
the vanishing of this class. This vanishing is \emph{automatic}
by the bar-intrinsic construction
(Theorem~\ref*{V1-thm:mc2-bar-intrinsic}): the global element
$\Theta_\cA = D_\cA - \dzero$ exists by virtue of $D_\cA^2 = 0$,
hence $\Ob_g$ is exact for every~$g$.
\end{remark}

\subsubsection{Explicit low-genus MC equations}
% label removed: subsubsec:thqg-VII-low-genus

The MC equation~\eqref{V1-eq:thqg-VII-genus-g-mc-eq} at each genus
from $g = 0$ through $g = 4$:

\begin{proposition}[Genus-$0$ MC equation; \ClaimStatusProvedHere]
% label removed: prop:thqg-VII-genus0
At genus $g = 0$, the MC equation reduces to the $\mathrm{Ch}_\infty$-algebra
structure equation on $\cA$:
\begin{equation}% label removed: eq:thqg-VII-mc-g0
D_0(\Theta^{(0)})
+ \tfrac{1}{2}[\Theta^{(0)},\, \Theta^{(0)}]
\;=\; 0.
\end{equation}
This is the statement that $\Theta^{(0)}$, the collection of
tree-level OPE structure maps, defines a homotopy chiral algebra
structure on~$\cA$.
\end{proposition}

\begin{proof}
At genus~$0$, the obstruction class $\Ob_0$ is empty (there are no
lower-genus terms), and the MC equation is simply the first equation
in the genus tower. The identity $D_0^2 = 0$ ensures that the
Arnold relation is satisfied, giving $d^2 = 0$ on the genus-$0$
bar complex.
\end{proof}

\begin{proposition}[Genus-$1$ MC equation; \ClaimStatusProvedHere]
\label{prop:thqg-VII-genus1}
\index{genus-1 MC equation}
At genus $g = 1$:
\begin{equation}% label removed: eq:thqg-VII-mc-g1
d_{\Theta^{(0)}}(\Theta^{(1)})
\;=\;
-\Delta(\Theta^{(0)})
- \dsew^{(1)},
\end{equation}
where $\Delta(\Theta^{(0)})$ is the non-separating BV operator
applied to the tree-level data, and $\dsew^{(1)}$ encodes the
genus-$1$ separating clutching (the correction from the boundary
divisor in $\overline{\cM}_{1,n}$ corresponding to a separating
node producing a genus-$0$ and a genus-$1$ component).

There are no bracket terms on the right
(since $g_1, g_2 \geq 1$ with $g_1 + g_2 = 1$ has no solutions)
and no planted-forest correction (the planted-forest
strata do not contribute at genus~$1$).

The separating clutching at genus~$1$ acts explicitly as
\begin{equation}% label removed: eq:thqg-VII-dsew1-explicit
\dsew^{(1)}(\varphi)
\;=\;
\sum_i \Delta_i(\varphi),
\end{equation}
where $\Delta_i$ contracts the $i$-th pair of inputs using
the invariant pairing and the clutching map
$\xi \colon \overline{\cM}_{0,n+2} \to \overline{\cM}_{1,n}$.
For the Heisenberg algebra at genus~$1$ with no external legs:
$\dsew^{(1)} = \kappa \cdot \Tr(\id) \cdot [\overline{\cM}_{1,1}]$,
where $\Tr$ is the supertrace. This term vanishes for the Heisenberg
because the separating boundary divisor
$\delta_0 \in \partial\overline{\cM}_{1,1}$ factorizes through two
genus-$0$ components with scalar data, producing a contribution
that cancels against the non-separating term in the one-channel case.
\end{proposition}

\begin{proof}
Specialize~\eqref{V1-eq:thqg-VII-genus-g-mc-eq} to $g = 1$. The
sum $\sum_{g_1+g_2=1, g_i \geq 1}$ is empty.
The planted-forest correction $\dpf^{(1)}$ vanishes because
the log-FM planted-forest boundary strata of
$\overline{\cM}_{1,n}$ are empty: every planted forest of
positive loop genus has first Betti number at least~$2$,
which exceeds genus~$1$. The only boundary stratum is the
non-separating node
$\delta_{\mathrm{irr}} \in \partial\overline{\cM}_{1,n}$.
\end{proof}

\begin{proposition}[Genus-$2$ MC equation; \ClaimStatusProvedHere]
\label{prop:thqg-VII-genus2}
\index{genus-2 MC equation}
At genus $g = 2$:
\begin{equation}% label removed: eq:thqg-VII-mc-g2
d_{\Theta^{(0)}}(\Theta^{(2)})
\;=\;
-\tfrac{1}{2}[\Theta^{(1)},\, \Theta^{(1)}]
- \Delta(\Theta^{(1)})
- \dsew^{(2)}
- \dpf^{(2)}.
\end{equation}
The right side decomposes into three shells, present at
\emph{every} degree~$k$:
\begin{enumerate}[label=\textup{(\roman*)}]
\item the \emph{loop-squared shell}
 $\Theta^{(2)}_{\mathrm{loop}^2} :=
 -\tfrac{1}{2}[\Theta^{(1)},\Theta^{(1)}]$,
 corresponding to the separating clutching
 $\sum_{g_1+g_2=2}[\Theta^{(g_1)},\Theta^{(g_2)}]
 = [\Theta^{(1)},\Theta^{(1)}]$
 (a genus-$0$ surface with two non-separating nodes);
\item the \emph{separating-loop shell}
 $\Theta^{(2)}_{\mathrm{sep}\circ\mathrm{loop}} :=
 -\Delta(\Theta^{(1)}) - \dsew^{(2)}$,
 corresponding to the non-separating BV operator
 $\Delta(\Theta^{(1)})$ (pinching a non-separating cycle)
 and the separating boundary strata $\dsew^{(2)}$;
\item the \emph{planted-forest shell}
 $\Theta^{(2)}_{\mathrm{pf}} :=
 -\dpf^{(2)}$,
 corresponding to the log-FM boundary strata indexed by
 depth-$1$ planted forests at genus~$2$.
\end{enumerate}
At genus~$2$ with degree~$k$, the source term $\Ob_2$ decomposes
as: (a)~separating clutching
$[\Theta^{(1)},\Theta^{(1)}]$, (b)~non-separating BV
$\Delta(\Theta^{(1)})$, and (c)~boundary strata contributions
from $\overline{\cM}_{2,k}$. These three shells are present
at every degree~$k$, not only at $(2,0)$: the boundary
stratification of $\overline{\cM}_{2,k}$ has the same three
types of codimension-$1$ strata (separating node, non-separating
node, log-FM correction) for all~$k$.
\end{proposition}

\begin{proof}
Specialize~\eqref{V1-eq:thqg-VII-genus-g-mc-eq} to $g = 2$. The
sum $\sum_{g_1+g_2=2, g_i \geq 1}$ has one term: $g_1 = g_2 = 1$,
contributing $\tfrac{1}{2}[\Theta^{(1)},\Theta^{(1)}]$. The
three-shell decomposition follows from the source of each
contribution: the bracket term from the separating node
connecting two genus-$1$ components, the BV term from the
non-separating node applied to genus-$1$ data, and the
planted-forest term from the Mok boundary strata.
\end{proof}

\begin{proposition}[Genus-$3$ MC equation; \ClaimStatusProvedHere]
\label{prop:thqg-VII-genus3}
\index{genus-3 MC equation}
At genus $g = 3$:
\begin{equation}% label removed: eq:thqg-VII-mc-g3
d_{\Theta^{(0)}}(\Theta^{(3)})
\;=\;
-[\Theta^{(1)},\, \Theta^{(2)}]
- \Delta(\Theta^{(2)})
- \dsew^{(3)}
- \dpf^{(3)}.
\end{equation}
The bracket sum $\sum_{g_1+g_2=3, g_i \geq 1}$ has two terms:
$(g_1,g_2) = (1,2)$ and $(2,1)$, which contribute
$[\Theta^{(1)},\Theta^{(2)}]$ by graded antisymmetry
(the factor $\tfrac{1}{2}$ is cancelled by the two orderings).
\end{proposition}

\begin{proof}
Specialize~\eqref{V1-eq:thqg-VII-genus-g-mc-eq} to $g = 3$. The
sum has $(g_1,g_2) \in \{(1,2),(2,1)\}$:
$\tfrac{1}{2}\bigl([\Theta^{(1)},\Theta^{(2)}]
+ [\Theta^{(2)},\Theta^{(1)}]\bigr)
= \tfrac{1}{2}\bigl([\Theta^{(1)},\Theta^{(2)}]
- [\Theta^{(1)},\Theta^{(2)}]\bigr)
+ [\Theta^{(1)},\Theta^{(2)}]
= [\Theta^{(1)},\Theta^{(2)}]$,
since $[\Theta^{(2)},\Theta^{(1)}] = -[\Theta^{(1)},
\Theta^{(2)}]$ (both elements are odd degree).
\end{proof}

\begin{proposition}[Genus-$4$ MC equation; \ClaimStatusProvedHere]
\label{prop:thqg-VII-genus4}
\index{genus-4 MC equation}
At genus $g = 4$:
\begin{equation}% label removed: eq:thqg-VII-mc-g4
d_{\Theta^{(0)}}(\Theta^{(4)})
\;=\;
-[\Theta^{(1)},\, \Theta^{(3)}]
- \tfrac{1}{2}[\Theta^{(2)},\, \Theta^{(2)}]
- \Delta(\Theta^{(3)})
- \dsew^{(4)}
- \dpf^{(4)}.
\end{equation}
The bracket sum has three terms:
$(g_1,g_2) \in \{(1,3),(2,2),(3,1)\}$, yielding
$[\Theta^{(1)},\Theta^{(3)}]
+ \tfrac{1}{2}[\Theta^{(2)},\Theta^{(2)}]$.
\end{proposition}

\begin{proof}
The computation parallels the previous cases. The three terms
in the bracket sum are:
\begin{align*}
&\tfrac{1}{2}\bigl(
[\Theta^{(1)},\Theta^{(3)}]
+ [\Theta^{(2)},\Theta^{(2)}]
+ [\Theta^{(3)},\Theta^{(1)}]
\bigr)\\
&= \tfrac{1}{2}\bigl(
2[\Theta^{(1)},\Theta^{(3)}]
+ [\Theta^{(2)},\Theta^{(2)}]
\bigr)\\
&= [\Theta^{(1)},\Theta^{(3)}]
+ \tfrac{1}{2}[\Theta^{(2)},\Theta^{(2)}],
\end{align*}
using $[\Theta^{(3)},\Theta^{(1)}]
= -[\Theta^{(1)},\Theta^{(3)}]$ and
$[\Theta^{(2)},\Theta^{(2)}] \neq 0$ in general
(it vanishes only in the Gaussian case).
\end{proof}

\begin{remark}[General pattern]
% label removed: rem:thqg-VII-general-pattern
\index{genus-$g$ MC equation!general structure}
At genus $g \geq 1$, the obstruction class
$\Ob_g(\Theta^{(<g)})$ has $\lfloor g/2 \rfloor + 2$ distinct
contributions:
\begin{enumerate}[label=(\alph*)]
\item $\lfloor g/2 \rfloor$ bracket terms
 $[\Theta^{(g_1)},\Theta^{(g_2)}]$ for $1 \leq g_1 \leq g_2$,
 $g_1 + g_2 = g$;
\item one BV term $\Delta(\Theta^{(g-1)})$;
\item boundary-stratum corrections $\dsew^{(g)} + \dpf^{(g)}$.
\end{enumerate}
The bracket terms record the \emph{separating degenerations}
(pinching a cycle that separates the surface into two parts of
genera $g_1$ and $g_2$). The BV term records the
\emph{non-separating degeneration} (pinching a non-separating
cycle, which lowers the genus by one). The boundary terms encode
higher-codimension strata. The structure is:
\begin{equation}% label removed: eq:thqg-VII-obs-general
\Ob_g
\;=\;
\underbrace{
\sum_{\substack{g_1+g_2=g \\ g_1,g_2 \geq 1}}
\tfrac{1}{2}[\Theta^{(g_1)},\Theta^{(g_2)}]
}_{\text{separating}}
\;+\;
\underbrace{
\Delta(\Theta^{(g-1)})
}_{\text{non-separating}}
\;+\;
\underbrace{
\dsew^{(g)} + \dpf^{(g)}
}_{\text{boundary strata}}.
\end{equation}
This is the genus analogue of the Postnikov tower in
obstruction theory: the extension from genus~$<g$ to genus~$g$
is controlled by a class in $H^2$ of the tree-twisted complex.
\end{remark}

\subsubsection{The bracket and BV terms in graph-sum language}
% label removed: subsubsec:thqg-VII-graph-sum

The bracket $[\Theta^{(g_1)},\Theta^{(g_2)}]$ in the
convolution algebra $\gAmod$ has a graph-theoretic interpretation.
A genus-$g_1$ graph $\Gamma_1$ and a genus-$g_2$ graph $\Gamma_2$
are joined by a single edge to produce a genus-$(g_1{+}g_2)$
graph:
\begin{equation}% label removed: eq:thqg-VII-bracket-graph
[\Theta^{(g_1)},\Theta^{(g_2)}]
\;=\;
\sum_{\substack{\Gamma_1 \in \mathsf{Gr}_{g_1}^{\mathrm{st}}\\
 \Gamma_2 \in \mathsf{Gr}_{g_2}^{\mathrm{st}}}}
\frac{1}{|\Aut(\Gamma_1)||\Aut(\Gamma_2)|}
\;\ell_{\Gamma_1 \cup_e \Gamma_2},
\end{equation}
where $\ell_\Gamma$ is the graph amplitude
(vertex labels from the transferred cyclic minimal model, edge
contractions by the complementarity propagator
$P_\cA = H_\cA^{-1}$), and $\Gamma_1 \cup_e \Gamma_2$ denotes
the graph obtained by joining $\Gamma_1$ and $\Gamma_2$ with a
single new edge.

The BV operator $\Delta(\Theta^{(g-1)})$ has the interpretation
of closing a single edge within a genus-$(g{-}1)$ graph to
produce a genus-$g$ graph:
\begin{equation}% label removed: eq:thqg-VII-bv-graph
\Delta(\Theta^{(g-1)})
\;=\;
\sum_{\Gamma \in \mathsf{Gr}_{g-1}^{\mathrm{st}}}
\frac{1}{|\Aut(\Gamma)|}
\sum_{v \in V(\Gamma)}
\ell_{\Gamma^v},
\end{equation}
where $\Gamma^v$ denotes the graph obtained by adding a loop edge
at vertex~$v$ of~$\Gamma$ (joining two half-edges at $v$ with
the propagator, increasing the loop number by one).

\begin{remark}[Feynman diagram interpretation]
% label removed: rem:thqg-VII-feynman-interpretation
\index{Feynman diagrams!MC recursion}
Equations~\eqref{V1-eq:thqg-VII-bracket-graph}
and~\eqref{V1-eq:thqg-VII-bv-graph} are the standard Feynman
diagram operations: the bracket is the tree-level connecting
propagator (joining two subdiagrams with one line), and the BV
operator is the loop closure (contracting one internal line into
a loop). The MC recursion is thus equivalent to the statement
that the genus-$g$ amplitude is determined by lower-genus
amplitudes through these two operations plus boundary corrections.
This is the chiral-algebraic version of the cut-and-sew relations
in string theory.
\end{remark}

% ======================================================================
%
% PART 2: THE MC RECURSION
%
% ======================================================================

\subsection{The MC recursion: unique determination from lower genera}
% label removed: subsec:thqg-VII-mc-recursion
\index{MC recursion|textbf}
\index{contracting homotopy!MC recursion}

\subsubsection{The contracting homotopy}
% label removed: subsubsec:thqg-VII-contracting-homotopy

The solvability of the recursion~\eqref{V1-eq:thqg-VII-genus-g-mc-eq}
rests on two facts:
\begin{enumerate}[label=(\alph*)]
\item the obstruction class $\Ob_g$ is
 $d_{\Theta^{(0)}}$-exact (hence the equation has a
 solution), and
\item the solution is unique up to gauge equivalence
 (hence the genus-$g$ partition function is determined).
\end{enumerate}
Both facts follow from the existence of a contracting homotopy
for the tree-twisted complex.

\begin{theorem}[Contracting homotopy; \ClaimStatusProvedHere]
\label{thm:thqg-VII-contracting-homotopy}
\index{contracting homotopy|textbf}
Let $\cA$ be a modular Koszul chiral algebra with bar-intrinsic
MC element $\Theta_\cA = D_\cA - \dzero$
\textup{(}Theorem~\textup{\ref*{V1-thm:mc2-bar-intrinsic})}. Denote
the tree-level data $\Theta^{(0)} \in \gr^0_F\,\gAmod$. Then:

\begin{enumerate}[label=\textup{(\roman*)}]
\item The tree-twisted complex $(\gAmod, d_{\Theta^{(0)}})$
 admits a contracting homotopy
 $\hCont \colon {\gAmod}^{\bullet}[\mathrm{genus\,}g]
 \to {\gAmod}^{\bullet - 1}[\mathrm{genus\,}g]$
 satisfying
 \begin{equation}% label removed: eq:thqg-VII-homotopy
 d_{\Theta^{(0)}} \circ \hCont
 + \hCont \circ d_{\Theta^{(0)}}
 \;=\; \id - \pi_H,
 \end{equation}
 where $\pi_H$ is the projection onto the cohomology
 $H^*(\gAmod[\mathrm{genus\,}g], d_{\Theta^{(0)}})$.

\item The cohomology $H^*(\gAmod[\mathrm{genus\,}g],
 d_{\Theta^{(0)}})$ vanishes in degree~$2$ for $g \geq 1$:
 \begin{equation}% label removed: eq:thqg-VII-h2-vanishing
 H^2\bigl(\gAmod[\mathrm{genus\,}g],\; d_{\Theta^{(0)}}\bigr)
 \;=\; 0, \qquad g \geq 1.
 \end{equation}
 In particular, every $d_{\Theta^{(0)}}$-cocycle of
 degree~$2$ is exact, and the obstruction class $[\Ob_g] = 0$.
\end{enumerate}
\end{theorem}

\begin{proof}
\emph{Part (i).} The contracting homotopy is constructed by
homological perturbation. Start with the untwisted complex
$(\gAmod, D_0)$ at genus~$g$. This is a chain complex of
cochain-valued functions on moduli spaces, and its cohomology
at genus~$g$ is computed by the $E_1$~page of the genus
spectral sequence
(Construction~\ref{V1-const:vol1-genus-spectral-sequence}).
The untwisted complex $(\gAmod[\text{genus }g], D_0)$
admits a standard contracting homotopy $\hCont_0$ from the
Hodge-type decomposition (choosing a metric on $\overline{\cM}_{g,n}$
gives a Hodge decomposition, and the Green's operator provides
$\hCont_0$).

The twist by $\Theta^{(0)}$ replaces $D_0$ with $d_{\Theta^{(0)}}
= D_0 + \ad_{\Theta^{(0)}}$. By the homological perturbation
lemma (HPL), the contracting homotopy for the twisted complex is
\begin{equation}% label removed: eq:thqg-VII-hpl-homotopy
\hCont
\;=\;
\hCont_0 \circ
\sum_{n \geq 0}
\bigl(-\ad_{\Theta^{(0)}} \circ \hCont_0\bigr)^n
\;=\;
\hCont_0 \circ
\bigl(\id + \ad_{\Theta^{(0)}} \circ \hCont_0\bigr)^{-1},
\end{equation}
which converges because the degree filtration makes
$\ad_{\Theta^{(0)}} \circ \hCont_0$ nilpotent on each
finite-degree component (the degree-cutoff of
Lemma~\ref{V1-lem:degree-cutoff} ensures that the series terminates
at each finite weight).

\medskip\noindent
\emph{Part (ii).} The vanishing of $H^2$ at positive genus
follows from the Koszul property. At genus~$0$, the Koszul
assumption gives diagonal Ext vanishing for $\cA$, which
translates to acyclicity of the tree-twisted bar complex in
the relevant degrees. At genus $g \geq 1$, the tree-twisted
differential $d_{\Theta^{(0)}}$ on the genus-$g$ sector of
$\gAmod$ is a deformation of the genus-$0$ differential
$D_0$; the deformation is controlled by the degree filtration,
and the acyclicity propagates by induction on degree
(each degree component is a finite-dimensional perturbation
of the acyclic genus-$0$ complex).

More precisely, the global solution $\Theta_\cA = D_\cA - \dzero$
(Theorem~\ref*{V1-thm:mc2-bar-intrinsic}) exists at all genera.
Its existence implies that every $\Ob_g$ is exact, hence
$H^2 = 0$ at every positive genus. This is not circular: the
bar-intrinsic construction produces $\Theta_\cA$ directly
from $D_\cA^2 = 0$, without solving the recursion step by step.
The recursion is a \emph{decomposition} of the already-existing
element into genus components, not a \emph{construction} that
might fail.
\end{proof}

\subsubsection{The recursion algorithm}
% label removed: subsubsec:thqg-VII-recursion-algorithm

\begin{construction}[MC recursion; \ClaimStatusProvedHere]
\label{constr:thqg-VII-mc-recursion}
\index{MC recursion!algorithm}
Given the genus-$0$ data $\Theta^{(0)}$ and the contracting
homotopy $\hCont$ of
Theorem~\textup{\ref{thm:thqg-VII-contracting-homotopy}},
the genus-$g$ MC element is computed recursively:

\medskip
\noindent
\textbf{Input:} $\Theta^{(0)}, \Theta^{(1)}, \ldots,
\Theta^{(g-1)}$.

\medskip
\noindent
\textbf{Step 1.} Compute the obstruction class:
\begin{equation}% label removed: eq:thqg-VII-recursion-obs
\Ob_g
\;:=\;
\sum_{\substack{g_1+g_2=g \\ g_1,g_2 \geq 1}}
\tfrac{1}{2}[\Theta^{(g_1)},\Theta^{(g_2)}]
\;+\;
\Delta(\Theta^{(g-1)})
\;+\;
\dsew^{(g)} + \dpf^{(g)}.
\end{equation}

\textbf{Step 2.} Verify the cocycle condition:
\begin{equation}% label removed: eq:thqg-VII-recursion-cocycle
d_{\Theta^{(0)}}(\Ob_g) = 0.
\end{equation}
This is a consistency check; it must hold by the inductive
hypothesis.

\textbf{Step 3.} Apply the contracting homotopy:
\begin{equation}% label removed: eq:thqg-VII-recursion-solve
\Theta^{(g)}
\;:=\;
-\hCont\bigl(\Ob_g\bigr).
\end{equation}

\textbf{Step 4.} Verify:
$d_{\Theta^{(0)}}(\Theta^{(g)})
= -(\id - \pi_H)(\Ob_g)
= -\Ob_g$
(since $\pi_H(\Ob_g) = 0$ by the vanishing of $H^2$).

\medskip
\noindent
\textbf{Output:} $\Theta^{(g)}$, uniquely determined up to
gauge equivalence.
\end{construction}

\begin{theorem}[Uniqueness up to gauge; \ClaimStatusProvedHere]
% label removed: thm:thqg-VII-uniqueness
\index{MC recursion!uniqueness}
The genus-$g$ MC element $\Theta^{(g)}$ is unique up to the
action of the gauge group
$\exp\bigl({\gAmod}^0[\mathrm{genus\,}g]\bigr)$.
Two solutions $\Theta^{(g)}$ and $\widetilde{\Theta}^{(g)}$
of~\eqref{V1-eq:thqg-VII-genus-g-mc-eq} with the same
lower-genus data $\Theta^{(<g)}$ satisfy
\begin{equation}% label removed: eq:thqg-VII-gauge-diff
\Theta^{(g)} - \widetilde{\Theta}^{(g)}
\;\in\;
\ker(d_{\Theta^{(0)}})
\;=\;
\im(d_{\Theta^{(0)}})
\;\oplus\;
H^1(\gAmod[\text{genus }g], d_{\Theta^{(0)}}),
\end{equation}
and the gauge group acts transitively on the fiber over any
fixed obstruction class.
\end{theorem}

\begin{proof}
Let $\Theta^{(g)}$ and $\widetilde{\Theta}^{(g)}$ be two
solutions with the same obstruction class. Their difference
$\xi := \Theta^{(g)} - \widetilde{\Theta}^{(g)}$ satisfies
$d_{\Theta^{(0)}}(\xi) = 0$. The gauge group
$G := \exp({\gAmod}^0[\text{genus }g])$ is pro-unipotent
(the weight filtration provides the lower central series), hence
acts freely on degree-$1$ cocycles. The gauge orbit of~$\xi$
is exactly the set of degree-$1$ cocycles cohomologous to~$\xi$.
Since $H^2 = 0$
(Theorem~\ref{thm:thqg-VII-contracting-homotopy}(ii)),
the ambiguity lies in $H^1$ and is absorbed by the gauge action.

The bar-intrinsic element $\Theta^{(g)} = (D_\cA - \dzero)^{(g)}$
is the canonical representative: it is the unique element in
the gauge orbit selected by the property $D_\cA^2 = 0$ (which
fixes the contracting homotopy to be the one induced by the
bar differential).
\end{proof}

\subsubsection{Base case: genus $g = 1$}
% label removed: subsubsec:thqg-VII-base-case

\begin{proposition}[Genus-$1$ solution; \ClaimStatusProvedHere]
\label{prop:thqg-VII-genus1-solution}
\index{genus-1 MC element}
The genus-$1$ MC element is
\begin{equation}% label removed: eq:thqg-VII-theta1
\Theta^{(1)}
\;=\;
-\hCont\bigl(
\Delta(\Theta^{(0)}) + \dsew^{(1)}
\bigr).
\end{equation}
For a one-channel algebra (e.g.\ the Heisenberg, the affine
Kac--Moody at generic level), this reduces to
\begin{equation}% label removed: eq:thqg-VII-theta1-scalar
\Theta^{(1)}
\;=\;
\kappa(\cA) \cdot [\delta_{\mathrm{irr}}] \otimes \Delta,
\end{equation}
where $[\delta_{\mathrm{irr}}] \in H^1(\overline{\cM}_{1,1})$
is the class of the irreducible boundary divisor and $\Delta$ is
the BV operator (the propagator inserted at the non-separating
node). The scalar $\kappa(\cA)$ is the modular characteristic
\textup{(}Theorem~D\textup{)}.
\end{proposition}

\begin{proof}
The obstruction class at genus~$1$ is
$\Ob_1 = \Delta(\Theta^{(0)}) + \dsew^{(1)}$
(Proposition~\ref{prop:thqg-VII-genus1}). The solution is
$\Theta^{(1)} = -\hCont(\Ob_1)$ by
Construction~\ref{constr:thqg-VII-mc-recursion}.

For a one-channel algebra, the BV operator $\Delta$ applied to
the genus-$0$ data $\Theta^{(0)}$ produces a scalar:
$\Delta(\Theta^{(0)}) = \kappa(\cA) \cdot \omega_1$, where
$\omega_1$ is the genus-$1$ modular form (the holomorphic
$1$-form on $\overline{\cM}_{1,1}$). The contracting homotopy
inverts the tree-twisted differential, and the result is the
scalar expression~\eqref{V1-eq:thqg-VII-theta1-scalar}.
The separating clutching $\dsew^{(1)}$ vanishes for genus~$1$
of the one-channel algebra because the separating boundary
$\delta_0 \in \partial\overline{\cM}_{1,1}$ decomposes the
surface into genus-$0$ and genus-$1$ components, and the
genus-$1$ component at the first step of the recursion
contributes only through $\Delta$.
\end{proof}

\subsubsection{Worked example: genus $g = 2$}
% label removed: subsubsec:thqg-VII-genus2-worked

\begin{computation}[Genus-$2$ MC recursion; \ClaimStatusProvedHere]
% label removed: comp:thqg-VII-genus2
\index{genus-2 MC element!worked computation}
We carry out the recursion for genus $g = 2$ in full detail.

\medskip\noindent\emph{Step 1: Compute the obstruction.}
From Proposition~\ref{prop:thqg-VII-genus2}:
\begin{equation}% label removed: eq:thqg-VII-obs2-full
\Ob_2
\;=\;
\tfrac{1}{2}[\Theta^{(1)},\Theta^{(1)}]
+ \Delta(\Theta^{(1)})
+ \dsew^{(2)} + \dpf^{(2)}.
\end{equation}

\medskip\noindent\emph{Step 2: Evaluate each term.}

\emph{(a) The bracket term.} Using
$\Theta^{(1)} = \kappa \cdot [\delta_{\mathrm{irr}}]
\otimes \Delta$ (for a one-channel algebra):
\begin{align}
\tfrac{1}{2}[\Theta^{(1)},\Theta^{(1)}]
&= \tfrac{1}{2}\,\kappa^2\,
 [[\delta_{\mathrm{irr}}] \otimes \Delta,\,
 [\delta_{\mathrm{irr}}] \otimes \Delta]
\notag\\
&= \tfrac{1}{2}\,\kappa^2\,
 [\delta_{\mathrm{irr}} \times \delta_{\mathrm{irr}}]
 \otimes \Delta \circ \Delta,
% label removed: eq:thqg-VII-bracket-term-g2
\end{align}
where $\delta_{\mathrm{irr}} \times \delta_{\mathrm{irr}}$
is the class in $H^*(\overline{\cM}_{2,n})$ corresponding to the
boundary stratum with two non-separating nodes, and
$\Delta \circ \Delta$ is the iterated BV operator (double
trace on the propagator).

\emph{(b) The BV term.} The BV operator applied to the genus-$1$
element:
\begin{equation}% label removed: eq:thqg-VII-bv-term-g2
\Delta(\Theta^{(1)})
= \kappa \cdot \Delta([\delta_{\mathrm{irr}}] \otimes \Delta)
= \kappa \cdot [\Delta(\delta_{\mathrm{irr}})]
 \otimes \Delta^2,
\end{equation}
where $\Delta(\delta_{\mathrm{irr}})$ is the pushforward of
the irreducible boundary class through the non-separating
clutching.

\emph{(c) The boundary terms.} The separating clutching
$\dsew^{(2)}$ and the planted-forest correction $\dpf^{(2)}$
encode the codimension-$1$ and codimension-$2$ boundary strata
of $\overline{\cM}_{2,n}$.

\medskip\noindent\emph{Step 3: Apply the homotopy.}
\begin{equation}% label removed: eq:thqg-VII-theta2-solution
\Theta^{(2)}
= -\hCont\bigl(\Ob_2\bigr)
= -\hCont\bigl(
 \tfrac{1}{2}\kappa^2\,
 [\delta_{\mathrm{irr}}^2] \otimes \Delta^2
 + \kappa\,[\Delta(\delta_{\mathrm{irr}})]
 \otimes \Delta^2
 + \dsew^{(2)} + \dpf^{(2)}
 \bigr).
\end{equation}

\medskip\noindent\emph{Step 4: Shell decomposition.} The
three shells (Proposition~\ref{prop:thqg-VII-genus2}) are:
\begin{align}
\Theta^{(2)}_{\mathrm{loop}^2}
&= -\hCont\bigl(\tfrac{1}{2}[\Theta^{(1)},\Theta^{(1)}]\bigr)
\neq 0 \text{ for all $\kappa \neq 0$},
% label removed: eq:thqg-VII-shell-loop2\\
\Theta^{(2)}_{\mathrm{sep}\circ\mathrm{loop}}
&= -\hCont\bigl(\Delta(\Theta^{(1)}) + \dsew^{(2)}\bigr)
\neq 0 \text{ iff } \fC(\cA) \neq 0,
% label removed: eq:thqg-VII-shell-sep\\
\Theta^{(2)}_{\mathrm{pf}}
&= -\hCont\bigl(\dpf^{(2)}\bigr)
\neq 0 \text{ iff } S_3 \neq 0
\;\text{(cubic shadow nonzero, from genus-$2$ planted-forest graphs)}.
% label removed: eq:thqg-VII-shell-pf
\end{align}
By the genus-$2$ shell activation theorem
(Theorem~\ref{V1-rem:genus2-shell-activation}), the separating shell~$\sigma$
distinguishes $\{G,C\}$ from $\{L,M\}$; the planted-forest shell
depends on $S_3$ (not just $S_4$), so classes $L$ and~$M$ share the
same genus-$2$ pattern. The full four-class diagnostic requires the
critical discriminant $\Delta = 8\kappa S_4$:
\begin{center}
\begin{tabular}{ccccc}
$\sigma$ & $\pi$ & $\Delta$ & Class & Example \\
\hline
$0$ & $0$ & $0$ & G (Gaussian) & $\cH_k$ \\
$1$ & $1$ & $0$ & L (Lie/tree) & $V_k(\fg)$ \\
$0$ & $0$ & $\neq 0$ & C (contact) & $\beta\gamma$ \\
$1$ & $1$ & $\neq 0$ & M (mixed) & $\mathrm{Vir}_c$
\end{tabular}
\end{center}
\end{computation}

\begin{remark}[Genus-$2$ as the first diagnostic]
% label removed: rem:thqg-VII-genus2-diagnostic
\index{genus-2 MC element!diagnostic role}
Genus~$2$ is the first level at which the MC recursion
distinguishes between the four shadow-depth classes. At genus~$1$,
the MC element $\Theta^{(1)}$ is controlled entirely by the scalar
$\kappa$; all four classes look identical at genus~$1$. At genus~$2$,
the three shells probe the cubic and quartic shadow data, and the
vanishing pattern separates the classes. This is a concrete
manifestation of the principle that the modular bootstrap becomes
nontrivial at genus~$2$.
\end{remark}

% ======================================================================
%
% PART 3: THE GENUS SPECTRAL SEQUENCE
%
% ======================================================================

\subsection{The genus spectral sequence}
% label removed: subsec:thqg-VII-genus-spectral-sequence
\index{genus spectral sequence|textbf}
\index{spectral sequence!genus filtration!bootstrap interpretation}

The MC recursion is organized by a spectral sequence that arises
from the genus filtration on $\gAmod$. This spectral sequence was
constructed in
Construction~\ref{V1-const:vol1-genus-spectral-sequence}; here we
develop its bootstrap interpretation and prove convergence.

\subsubsection{Construction}
% label removed: subsubsec:thqg-VII-ss-construction

\begin{theorem}[Genus spectral sequence for the bootstrap;
\ClaimStatusProvedHere]
\label{thm:thqg-VII-genus-ss}
\index{genus spectral sequence!construction}
The genus filtration~\eqref{V1-eq:thqg-VII-genus-filtration}
on the modular convolution algebra $\gAmod$ induces a
spectral sequence
\begin{equation}% label removed: eq:thqg-VII-genus-ss
E_1^{p,q}
\;\Longrightarrow\;
H^{p+q}\bigl(\gAmod,\; D\bigr),
\end{equation}
with the following structure:

\begin{enumerate}[label=\textup{(\roman*)}]
\item \emph{$E_0$~page:}
 $E_0^{p,q} = \gr^p_F\,{\gAmod}^{p+q}$,
 the genus-$p$ component in total degree $p + q$.
 The differential is $d_0 = D_0$ restricted to
 $\gr^p_F$, which is the genus-preserving part of~$D$.

\item \emph{$E_1$~page:}
 \begin{equation}% label removed: eq:thqg-VII-E1
 E_1^{p,q}
 \;=\;
 H^{p+q}\bigl(\gr^p_F\,\gAmod,\; D_0\bigr),
 \end{equation}
 the cohomology of the genus-$p$ complex with
 genus-preserving differential. The three columns are:
 \begin{itemize}
 \item $p = 0$: \emph{tree-level cohomology}
 $H^q(\gAmod[\text{genus }0], D_0)$, which is the
 bar cohomology $H^*(\barB(\cA))$. For Koszul~$\cA$,
 this is concentrated on the diagonal.
 \item $p = 1$: \emph{one-loop cohomology}
 $H^{1+q}(\gAmod[\text{genus }1], D_0)$, encoding
 the genus-$1$ bar complex modulo the genus-$0$
 differential.
 \item $p \geq 2$: \emph{higher-genus shells},
 computed by the moduli-space cohomology of
 $\overline{\cM}_{p,n}$ tensored with the cyclic
 bar data.
 \end{itemize}

\item \emph{$d_1$ differential:}
 $d_1 \colon E_1^{p,q} \to E_1^{p+1,q}$.
 This is the obstruction map
 $\Ob_{p+1} \colon E_1^{p,*} \to E_1^{p+1,*-1}$,
 which sends the genus-$p$ data to the genus-$(p{+}1)$
 obstruction class.
 \begin{equation}% label removed: eq:thqg-VII-d1
 d_1
 \;=\;
 [\Ob_{p+1}]
 \;\colon\;
 E_1^{p,q} \to E_1^{p+1,q}.
 \end{equation}
 The $d_1$ differential at $p = 0$ is the genus-$1$
 obstruction $\Ob_1 = \Delta + \dsew^{(1)}$:
 it sends tree-level data to the one-loop correction.

\item \emph{Higher differentials:}
 $d_r \colon E_r^{p,q} \to E_r^{p+r,q-r+1}$.
 The differential $d_r$ is the genus-$r$ obstruction map
 that sends genus-$p$ data through the MC recursion to the
 genus-$(p{+}r)$ sector. These encode the multi-loop
 corrections.
\end{enumerate}
\end{theorem}

\begin{proof}
The genus filtration is a complete descending filtration on a
cochain complex. The standard construction (cf.\
\cite[Ch.~5]{Wei94}) produces a spectral sequence with the
stated $E_0$~page and $E_1$~page.

\medskip\noindent\emph{The $d_0$ differential.}
On the associated graded $\gr^p_F$, the differential induced by
$D = D_0 + D_1$ is $D_0$ (the genus-preserving part), since
$D_1$ raises genus and hence maps into the next filtration step.
Thus $E_0^{p,q} = \gr^p_F\,{\gAmod}^{p+q}$ with $d_0 = D_0$.

\medskip\noindent\emph{The $d_1$ differential.}
The $d_1$ differential on the $E_1$~page is induced by $D_1$
on the $D_0$-cohomology. Explicitly: for a $D_0$-cocycle
$\alpha \in \gr^p_F$, choose a lift $\widetilde{\alpha}$ to
$F^p\,\gAmod$. Then $D(\widetilde{\alpha}) \in F^{p+1}$ and
$D_0(D(\widetilde{\alpha})) \in F^{p+1}$; the class of
$D(\widetilde{\alpha})$ modulo $D_0$-exact terms is the $d_1$
differential. By the decomposition $D = D_0 + D_1$:
\[
D(\widetilde{\alpha})
= D_0(\widetilde{\alpha}) + D_1(\widetilde{\alpha}).
\]
Since $D_0(\widetilde{\alpha}) = 0$ (as $\alpha$ is a cocycle),
we get $d_1([\alpha]) = [D_1(\widetilde{\alpha})]$ in
$E_1^{p+1,*}$. The map $D_1$ is exactly the BV operator
$\Delta$ plus the planted-forest corrections $\dpf$, which is
the genus-raising obstruction. Thus $d_1 = [\Ob_{p+1}]$.

\medskip\noindent\emph{Higher differentials.}
The $d_r$ differential is constructed by the standard
extension procedure: starting with a class in $E_r^{p,q}$
(which survives to the $E_r$~page because $d_1, \ldots,
d_{r-1}$ vanish on it), lift to an element of $F^p$, apply
$D$, and project to the genus-$(p{+}r)$ component modulo
earlier terms. The result is the genus-$(p{+}r)$ obstruction
accumulated from $r$ iterations of the genus-raising operations.
\end{proof}

\subsubsection{The $E_1$~page in detail}
% label removed: subsubsec:thqg-VII-E1-detail

\begin{proposition}[$E_1$~page for a Koszul algebra;
\ClaimStatusProvedHere]
% label removed: prop:thqg-VII-E1-koszul
\index{genus spectral sequence!$E_1$ page}
Let $\cA$ be a modular Koszul chiral algebra. Then:
\begin{enumerate}[label=\textup{(\roman*)}]
\item $E_1^{0,q} = H^q(\barB(\cA))$, the bar cohomology of
 $\cA$. By the Koszul property, this is concentrated on the
 diagonal: $E_1^{0,q} = 0$ for $q \neq 0$ (after appropriate
 regrading).
\item $E_1^{1,q}$ is the genus-$1$ bar cohomology, computed
 from the elliptic bar complex
 (Theorem~\ref{V1-thm:elliptic-bar}). For the Heisenberg
 algebra, $E_1^{1,q} = \mathbb{C}$ for $q = 0$ (generated by
 $\kappa \cdot \omega_1$) and vanishes for $q \neq 0$.
\item $E_1^{p,q}$ for $p \geq 2$ is the genus-$p$ bar
 cohomology, which is computed from the moduli-space
 cohomology $H^*(\overline{\cM}_{p,n})$ via the Mumford
 classes. For the Heisenberg algebra, each $E_1^{p,0}$ is
 one-dimensional, generated by $\kappa^p \cdot \lambda_p$.
\end{enumerate}
\end{proposition}

\begin{proof}
Part (i) is the definition of the Koszul property: the bar
cohomology $H^*(\barB(\cA))$ is the Koszul dual $\cA^!$,
concentrated in degree~$0$.

For parts (ii) and (iii) in the Heisenberg case, the genus-$p$
bar complex has differential $D_0 = \dzero$ (no collision
corrections, since the OPE is quadratic and abelian). The
cohomology is determined by the moduli-space integration:
at genus~$p$, the only nontrivial class is the top Hodge class
$\lambda_p = c_p(\mathbb{E})$, paired with the scalar modular
characteristic $\kappa^p$. This gives a one-dimensional
$E_1^{p,0} = \mathbb{C} \cdot \kappa^p \lambda_p$.
\end{proof}

\begin{remark}[The $E_1$~page as the ``classical bootstrap'']
% label removed: rem:thqg-VII-E1-classical
The $E_1$~page isolates the data at each genus before the
inter-genus corrections are turned on. One may think of
$E_1^{p,*}$ as the ``classical'' (i.e.\ tree-level with respect
to the genus filtration) contribution at genus~$p$. The
higher differentials $d_r$ encode the quantum corrections that
propagate from lower genera to higher genera. The spectral
sequence converges to the full MC moduli, which incorporates
all inter-genus interactions.
\end{remark}

\subsubsection{Convergence}
% label removed: subsubsec:thqg-VII-convergence

\begin{theorem}[Convergence of the genus spectral sequence;
\ClaimStatusProvedHere]
\label{thm:thqg-VII-ss-convergence}
\index{genus spectral sequence!convergence}
The spectral
sequence~\eqref{V1-eq:thqg-VII-genus-ss} converges:
\begin{equation}% label removed: eq:thqg-VII-convergence
E_\infty^{p,q}
\;\cong\;
\gr^p_F\, H^{p+q}\bigl(\gAmod,\; D\bigr).
\end{equation}
\end{theorem}

\begin{proof}
The genus filtration is complete
($\gAmod = \varprojlim_g\, \gAmod / F^g$) and
exhaustive ($F^0 = \gAmod$). The Eilenberg--Moore
convergence theorem (cf.\ \cite[Theorem~5.5.1]{Wei94})
applies to complete, exhaustive filtrations on cochain
complexes: the spectral sequence converges to the graded
pieces of the cohomology of the total complex.

The completeness condition is exactly the $\hbar$-adic
completeness of
Remark~\ref{rem:thqg-VII-hbar-topology}: the
$\hbar$-adic topology is the linear topology defined by
the genus filtration, and the convolution algebra is complete
(being defined as an infinite product over genera).

The pronilpotence of $F^1\,\gAmod$ ensures that the spectral
sequence is strongly convergent: for each total degree $n$,
only finitely many nonzero differentials land in a given
bidegree, so the inverse system of quotients stabilizes.
\end{proof}

\begin{corollary}[Degeneration for Gaussian algebras;
\ClaimStatusProvedHere]
\label{cor:thqg-VII-gaussian-degeneration}
\index{genus spectral sequence!degeneration}
\index{Gaussian archetype!spectral sequence}
For a Gaussian algebra $\cA$ \textup{(}shadow depth $r_{\max} = 2$,
e.g.\ the Heisenberg algebra\textup{)}, the genus spectral
sequence degenerates at $E_1$:
\begin{equation}% label removed: eq:thqg-VII-gaussian-degeneration
E_1^{p,q} \;=\; E_\infty^{p,q}.
\end{equation}
\end{corollary}

\begin{proof}
For a Gaussian algebra, the MC element $\Theta_\cA$ is
scalar: $\Theta_\cA = \kappa \cdot \eta \otimes \Lambda$,
with no higher-degree corrections. All inter-genus
brackets $[\Theta^{(g_1)},\Theta^{(g_2)}]$ and BV applications
$\Delta(\Theta^{(g)})$ produce scalar multiples of
$\kappa^{g+1} \cdot \lambda_{g+1}$; the differentials $d_r$
for $r \geq 1$ land on pre-existing classes (they are
multiplied by $\kappa$ at each step but do not create new
cohomological directions). The tree-twisted differential
$d_{\Theta^{(0)}}$ is exact on positive-degree classes at
every genus (the Koszul property), so $d_1 = 0$ on the
$E_1$~page: every obstruction class is automatically solved
by the scalar homotopy, and the spectral sequence collapses.

Concretely, the one-dimensional space $E_1^{p,0}$ survives
to $E_\infty$ because $d_r(E_1^{p,0}) = 0$ for all $r \geq 1$
(the scalar obstruction at each genus is solved by the scalar
homotopy, leaving no residual class on the $E_1$ page).
\end{proof}

\begin{proposition}[Non-degeneration for mixed-type algebras;
\ClaimStatusProvedHere]
\label{prop:thqg-VII-mixed-nondegeneration}
\index{genus spectral sequence!non-degeneration}
\index{mixed archetype!spectral sequence}
For a mixed-type algebra $\cA$ with $r_{\max} = \infty$
\textup{(}e.g.\ the Virasoro algebra\textup{)}, the genus
spectral sequence does not degenerate at any finite page:
for every $r \geq 1$, there exists a bidegree $(p,q)$ with
$d_r \neq 0$ on $E_r^{p,q}$.
\end{proposition}

\begin{proof}
The mixed archetype has $\fC(\cA) \neq 0$ and
$\mathfrak{Q}(\cA) \neq 0$, hence both the separating-loop
shell and the planted-forest shell are activated at genus~$2$
(Theorem~\ref{V1-rem:genus2-shell-activation}). The $d_1$
differential carries the genus-$0$ cubic shadow data into the
genus-$1$ sector, and this propagates: the genus-$r$
obstruction class receives contributions from the
$r$-fold iterated application of the BV operator and the
$r$-fold nested brackets, both of which are nontrivial when
the shadow obstruction tower is infinite.

More precisely, the $d_r$ differential at the $(0,*)$ column
maps tree-level data through the $r$-loop recursion. Since
the shadow obstruction tower $\Theta^{\leq r}$ is strictly larger than
$\Theta^{\leq r-1}$ for every~$r$ (the obstruction
$o_{r+1}(\cA) \neq 0$ by the infinite-depth condition),
the $d_r$ differential is nontrivial for each~$r$.
\end{proof}

\subsubsection{The spectral sequence and the shadow obstruction tower}
% label removed: subsubsec:thqg-VII-ss-shadow

\begin{theorem}[Spectral sequence differentials as shadow
obstructions; \ClaimStatusProvedHere]
\label{thm:thqg-VII-ss-shadow-obstruction}
\index{genus spectral sequence!shadow tower correspondence}
The differentials of the genus spectral sequence encode the
shadow obstruction tower:
\begin{enumerate}[label=\textup{(\roman*)}]
\item $d_1$ encodes the modular characteristic $\kappa(\cA)$:
 it is the scalar genus-raising map
 $\Theta^{(0)} \mapsto \kappa \cdot \omega_1$.
\item $d_2$ encodes the cubic shadow $\fC(\cA)$:
 it measures the failure of the genus-$2$ element to be
 determined solely by the scalar data.
\item $d_3$ encodes the quartic shadow $\mathfrak{Q}(\cA)$:
 it is the genus-$3$ obstruction class coming from the
 quartic contact invariant.
\item $d_r$ for $r \geq 4$ encodes the degree-$(r{+}1)$
 shadow: the $r$-loop obstruction from the shadow
 obstruction tower at weight~$r+1$.
\end{enumerate}
In the Gaussian archetype, all $d_r = 0$ for $r \geq 1$.
In the Lie/tree archetype, $d_1 \neq 0$ but $d_r = 0$ for
$r \geq 2$. In the contact archetype, $d_1 = d_2 = 0$,
$d_3 \neq 0$. In the mixed archetype, infinitely many
$d_r \neq 0$.
\end{theorem}

\begin{proof}
The identification follows from the explicit formula for the
differentials. At the $E_r$~page, the differential $d_r$
is induced by $r$ applications of the genus-raising part $D_1$,
composed with the tree-level contracting homotopy. The
$r$-fold application of $D_1$ produces an expression that
involves all brackets and BV operators of total genus shift~$r$,
which is exactly the degree-$(r{+}1)$ obstruction class
$o_{r+1}(\cA)$ projected through the shadow obstruction tower.

The vanishing pattern follows from the shadow depth
classification: if $\cA$ has shadow depth $r_{\max}$, then
$o_{r+1}(\cA) = 0$ for $r \geq r_{\max}$, hence $d_r = 0$
for $r \geq r_{\max}$ (as the obstruction classes that feed
the higher differentials vanish).
\end{proof}

\begin{remark}[Comparison with the PBW spectral sequence]
% label removed: rem:thqg-VII-pbw-vs-genus
\index{PBW spectral sequence!vs genus spectral sequence}
The genus spectral sequence is distinct from the PBW spectral
sequence. The PBW spectral sequence uses the conformal-weight
filtration \emph{within each fixed genus} to compute bar
cohomology; it is a computational tool for determining
$E_1^{p,*}$. The genus spectral sequence uses the genus
filtration \emph{across genera} to lift the MC element; it is
a structural tool for the bootstrap. The two spectral sequences
are complementary: the PBW spectral sequence computes the input
(the $E_1$~page), and the genus spectral sequence computes the
output (the full MC element $\Theta_\cA$).
\end{remark}

% ======================================================================
%
% PART 4: COMPARISON WITH THE CONFORMAL BOOTSTRAP
%
% ======================================================================

\subsection{Comparison with the conformal bootstrap}
% label removed: subsec:thqg-VII-bootstrap-comparison
\index{conformal bootstrap!comparison with MC recursion|textbf}

The conformal bootstrap programme (see \cite{BPRS15} for a review)
constrains correlation functions and partition functions by
imposing:
\begin{enumerate}[label=(\alph*)]
\item \emph{crossing symmetry} of the four-point function on the
 sphere;
\item \emph{modular invariance} of the genus-$1$ partition
 function under $\mathrm{SL}_2(\mathbb{Z})$;
\item \emph{unitarity} bounds on OPE coefficients;
\item \emph{causality/Regge-limit} constraints on the OPE
 convergence.
\end{enumerate}
These are \emph{necessary conditions} on the partition function,
and the programme proceeds by eliminating regions of parameter
space that violate them. The remaining allowed region is the
``bootstrap island'' in the space of CFT data.

The MC recursion is qualitatively different.

\begin{theorem}[MC recursion vs.\ conformal bootstrap;
\ClaimStatusProvedHere]
% label removed: thm:thqg-VII-mc-vs-bootstrap
\index{MC recursion!vs conformal bootstrap}
Let $\cA$ be a modular Koszul chiral algebra with boundary OPE
data $\Theta^{(0)}$. Then:

\begin{enumerate}[label=\textup{(\roman*)}]
\item \emph{Completeness.} The MC recursion
 \textup{(}Construction~\textup{\ref{constr:thqg-VII-mc-recursion})}
 determines $\Theta^{(g)}$ for every $g \geq 1$, uniquely
 up to gauge equivalence. The conformal bootstrap provides
 constraints but not unique determination.

\item \emph{Determinism.} The conformal bootstrap is a
 set of inequalities and consistency conditions; it does not
 specify the genus-$g$ partition function as a function of
 the genus-$0$ OPE data. The MC recursion does: the
 formula $\Theta^{(g)} = -\hCont(\Ob_g)$ is an explicit
 function of $\Theta^{(<g)}$.

\item \emph{Systematicity.} The genus spectral sequence
 organizes the MC recursion into a controlled sequence of
 algebraic extensions. There is no analogous systematic
 organization for the conformal bootstrap beyond genus~$1$.

\item \emph{Structural input.} The conformal bootstrap
 requires unitarity, which is a metric condition (positive
 definiteness of the Hilbert space inner product). The MC
 recursion requires the Koszul property, which is a homological
 condition (diagonal Ext vanishing). These are logically
 independent.

\item \emph{Scope.} The conformal bootstrap applies to
 unitary CFTs. The MC recursion applies to all modular
 Koszul chiral algebras, including non-unitary ones
 \textup{(}e.g.\ the Virasoro algebra at $c < 0$\textup{)}.
\end{enumerate}
\end{theorem}

\begin{proof}
(i) follows from the $H^2$-vanishing
(Theorem~\ref{thm:thqg-VII-contracting-homotopy}(ii))
and the existence of the bar-intrinsic element
(Theorem~\ref*{V1-thm:mc2-bar-intrinsic}).

(ii) follows from the explicit formula
$\Theta^{(g)} = -\hCont(\Ob_g)$, which is an algebraic
function of the lower-genus data.

(iii) follows from the convergent spectral sequence
(Theorem~\ref{thm:thqg-VII-ss-convergence}).

(iv) is a direct comparison of the structural inputs: the
Koszul property (Ext vanishing on the bar complex) versus
unitarity (positivity of the inner product). Neither
implies the other: there exist Koszul algebras that are not
unitary (e.g.\ $\mathrm{Vir}_c$ at $c < 0$), and there exist
unitary CFTs whose chiral algebras have non-formal Swiss-cheese
structure (e.g.\ the minimal models at exceptional values where the
transferred $\Ainf$ operations are nonvanishing).

(v) follows from (iv): the Koszul property is defined purely
in terms of the bar complex, without reference to an inner
product.
\end{proof}

\begin{remark}[Koszulness as bootstrap closure]
\label{rem:vol2-koszulness-bootstrap}
\index{conformal bootstrap!Koszulness correspondence}
The correspondence between Koszulness and bootstrap structure
is made precise in Volume~I
\textup{(}Theorem~\textup{\ref*{V1-thm:koszulness-bootstrap})}: for
a chirally Koszul algebra, the MC moduli is discrete
\textup{(}$\Theta_\cA$ is the unique solution\textup{)}, the shadow
depth $r_{\max}$ equals the bootstrap closure order, and the
$G$/$L$/$C$/$M$ classification corresponds to bootstrap complexity.
This sharpens the comparison above: where the conformal bootstrap
provides necessary constraints, the MC recursion on the Koszul locus
provides unique determination with a finite complexity measure.
\end{remark}

\begin{theorem}[Crossing symmetry from the MC equation;
\ClaimStatusProvedHere]
% label removed: thm:thqg-VII-crossing-from-mc
% label removed: rem:thqg-VII-crossing-as-mc
\index{crossing symmetry!as MC equation|textbf}
\index{conformal bootstrap!as MC projection|textbf}
Let $\cA$ be a modular Koszul chiral algebra with MC element
$\Theta_\cA$. The projection of the MC equation to genus~$0$,
degree~$2$ \textup{(}four external points with binary internal
pairing\textup{)} is equivalent to the conformal bootstrap
crossing-symmetry equation. Precisely: the MC equation
$D_0(\Theta^{(0)}) + \frac{1}{2}[\Theta^{(0)},\Theta^{(0)}] = 0$
evaluated at $n = 4$ external points gives
\begin{equation}% label removed: eq:thqg-VII-crossing-symmetry
\sum_p C_{12}^p\, C_{p34}\, G_p^{(s)}(z, \bar{z})
\;=\;
\sum_p C_{14}^p\, C_{p23}\, G_p^{(t)}(z, \bar{z}),
\end{equation}
where $C_{ij}^p$ are the OPE structure constants,
$G_p^{(s)}$ and $G_p^{(t)}$ are the $s$- and $t$-channel
conformal blocks, and
$z = z_{12}z_{34}/(z_{13}z_{24})$ is the cross-ratio.
\end{theorem}

\begin{proof}
The genus-$0$ MC equation at $n = 4$ reads
\[
D_0(\Theta^{(0)}_{n=4})
+ \tfrac{1}{2}[\Theta^{(0)}_{n=2},\, \Theta^{(0)}_{n=2}]
= 0.
\]
The bracket $[\Theta^{(0)}_{n=2}, \Theta^{(0)}_{n=2}]$
sums over the three channels in which two pairs of the four
external points collide:
\begin{itemize}
\item \emph{$s$-channel}: points $1,2$ collide, then $3,4$.
 The coderivation pairing contracts
 $\Theta^{(0)}_{n=2}(z_1, z_2)$ with
 $\Theta^{(0)}_{n=2}(z_3, z_4)$ at an internal point,
 producing $\sum_p C_{12}^p C_{p34} G_p^{(s)}$.
\item \emph{$t$-channel}: points $1,4$ collide, then $2,3$.
 This produces $\sum_p C_{14}^p C_{p23} G_p^{(t)}$.
\item \emph{$u$-channel}: points $1,3$ collide, then $2,4$.
 This produces $\sum_p C_{13}^p C_{p24} G_p^{(u)}$.
\end{itemize}
The MC equation requires the sum of all three channels to
vanish: this is the associativity of the OPE.

The three channels correspond to the three boundary divisors of
$\overline{\cM}_{0,4} \cong \mathbb{P}^1$:
$z = 0$ ($s$-channel), $z = \infty$ ($t$-channel),
$z = 1$ ($u$-channel). The boundary relation
$[z{=}0] + [z{=}1] + [z{=}\infty] = 0$ in
$H_0(\partial\overline{\cM}_{0,4})$ is the topological identity
underlying crossing symmetry; the MC equation at $(0,4)$ is its
transport through the convolution algebra.

The equivalence between OPE associativity and crossing symmetry
is standard:
$G_p^{(u)}(z, \bar{z}) = G_p^{(t)}(1-z, 1-\bar{z})$, so
the three-channel vanishing reduces to the
$s$--$t$ crossing equation~\eqref{V1-eq:thqg-VII-crossing-symmetry}.
\end{proof}

\begin{corollary}[Genus-$g$ bootstrap from MC;
\ClaimStatusProvedHere]
% label removed: cor:thqg-VII-genus-g-bootstrap
\index{conformal bootstrap!genus $g$}
The MC equation at $(g, n)$ is the genus-$g$, $n$-point
conformal bootstrap equation. At $n = 0$, it becomes the
genus-$g$ modular bootstrap:
\begin{equation}% label removed: eq:thqg-VII-genus-g-bootstrap
\Ob_g(\Theta^{(<g)}) \;\in\; \mathrm{Im}(d_{\Theta^{(0)}}).
\end{equation}
The exactness of $\Ob_g$ is the statement that the genus-$g$
partition function is uniquely determined \textup{(}up to gauge
equivalence\textup{)} by the lower-genus data. This is the
algebraic content of the modular bootstrap: modular invariance
is not an additional axiom but a consequence of the MC equation.
\end{corollary}

\begin{proof}
The MC equation at $(g, 0)$ is
$d_{\Theta^{(0)}}(\Theta^{(g)}) = -\Ob_g(\Theta^{(<g)})$.
The bar-intrinsic construction
(Theorem~\ref*{V1-thm:mc2-bar-intrinsic}) guarantees that
$\Theta_\cA$ exists, hence $\Ob_g$ is exact for
every~$g$. Uniqueness follows from acyclicity of the
tree-twisted complex (Koszul property).
\end{proof}

\begin{remark}[Higher-degree crossing]
% label removed: rem:thqg-VII-higher-degree-crossing
\index{crossing symmetry!higher degree}
The MC equation at $(0, n)$ for $n \geq 5$ gives higher-degree
crossing-symmetry constraints. The boundary of
$\overline{\cM}_{0,n}$ (the Stasheff associahedron $K_n$) has
codimension-$1$ faces indexed by $2$-level trees, and
$\partial^2 K_n = 0$ gives the $A_\infty$-associativity
equations. The MC equation at $(0, n)$ is the $n$-point
conformal bootstrap equation, encoding consistency of all
possible OPE orderings.
\end{remark}

\begin{remark}[Modular invariance as genus-$1$ MC]
% label removed: rem:thqg-VII-modular-as-mc
\index{modular invariance!as MC equation}
The conformal bootstrap's modular invariance condition is
a consequence of the genus-$1$ MC equation. The mapping class
group $\mathrm{MCG}(\Sigma_1) \cong \mathrm{SL}_2(\mathbb{Z})$
acts on the genus-$1$ sector of $\gAmod$; the MC equation at
genus~$1$, combined with the equivariance of the bar-intrinsic
construction under this action
(Proposition~\ref{V1-prop:mcg-equivariance-tower}), forces the
genus-$1$ partition function to be modular invariant.
Thus modular invariance is not an additional constraint
imposed from outside, but an automatic consequence of the
MC equation.
\end{remark}

\begin{remark}[Why the MC approach is ``stronger'']
% label removed: rem:thqg-VII-mc-stronger
The conformal bootstrap carves out allowed regions by
constraints (crossing, modular invariance, unitarity). The MC
recursion \emph{parameterizes} the moduli space
$\MC(\gAmod)/\text{gauge}$: given the genus-$0$ data, the
higher-genus data are uniquely computed via the contracting
homotopy $\hCont$, which exists because the Koszul property
gives acyclicity. The genus filtration supplies a coordinate
system in which each genus-$g$ coordinate is an explicit
function of the lower-genus coordinates.

The trade-off: the MC recursion requires more algebraic
structure (the full bar complex and Koszul property) but yields
unique determination. The conformal bootstrap requires less
structure (unitarity) but yields only constraints.
\end{remark}

\begin{proposition}[The bootstrap gap]
% label removed: prop:thqg-VII-bootstrap-gap
\index{bootstrap gap}
\ClaimStatusProvedHere
The conformal bootstrap's genus-$1$ modular constraint is
equivalent to the $d_1$ differential of the genus spectral
sequence:
\begin{equation}% label removed: eq:thqg-VII-bootstrap-gap
d_1 \colon E_1^{0,*} \to E_1^{1,*}
\quad\Longleftrightarrow\quad
\text{modular invariance of } Z_1(\cA).
\end{equation}
The ``bootstrap gap'' (the difference between what the
conformal bootstrap constrains and what the MC recursion
determines) corresponds to the higher differentials
$d_r$ for $r \geq 2$: these carry information about genus
$g \geq 2$ that is not accessible from genus-$0$ and genus-$1$
constraints alone. In the MC recursion, this information
is \emph{determined} by the higher differentials;
in the conformal bootstrap literature, no general mechanism
for accessing genus $g \geq 2$ constraints is currently known
(though specific examples, such as the genus-$2$ partition
function of the monster CFT, are fixed by other means).
\end{proposition}

\begin{proof}
The genus-$1$ modular constraint is the requirement that
$\Theta^{(1)}$ be a solution of the genus-$1$ MC equation
\eqref{V1-eq:thqg-VII-mc-g1}, which is precisely the condition
that the $d_1$ differential maps the tree-level data into a
coboundary in the genus-$1$ cohomology. The higher differentials
$d_r$ ($r \geq 2$) encode the genus-$(r{+}1)$ constraints, which
are uniquely determined by the MC recursion.
\end{proof}

% ======================================================================
%
% PART 5: EXPLICIT COMPUTATIONS FOR THE HEISENBERG ALGEBRA
%
% ======================================================================

\subsection{Heisenberg genus expansion through genus~$4$}
% label removed: subsec:thqg-VII-heisenberg-genus4
\index{Heisenberg algebra!genus expansion through genus 4|textbf}

The Heisenberg algebra $\cH_k$ is the Gaussian archetype:
the shadow obstruction tower terminates at weight~$2$ and
$\Theta_{\cH_k} = k \cdot \eta \otimes \Lambda$. The MC
recursion for $\cH_k$ is entirely scalar. We compute
$\Theta^{(g)}_{\cH_k}$ for $g = 0, 1, 2, 3, 4$ and verify
agreement with the Faber--Pandharipande $\lambda_g$-formula
(Theorem~\ref{V1-thm:dmvv-agreement}).

\subsubsection{Setup and notation}
% label removed: subsubsec:thqg-VII-heis-setup

The rank-$1$ Heisenberg algebra $\cH_k$ has a single generator~$J$
with OPE
\begin{equation}% label removed: eq:thqg-VII-heis-ope
J(z)\,J(w) \;\sim\; \frac{k}{(z-w)^2},
\end{equation}
where $k$ is the level. The modular characteristic is
$\kappa(\cH_k) = k$. The bar complex $\barB(\cH_k)$ is
generated by the logarithmic forms
$\eta_{ij} = d\log(z_i - z_j)$, and the bar differential satisfies
$d^2 = 0$ by the Arnold relation. The Koszul dual is the
free fermion $\cH_k^! = \mathrm{ff}_{k^{-1}}$ (not the
Heisenberg at level $k^{-1}$: the duality exchanges bosonic
and fermionic statistics).

The genus expansion has the form
\begin{equation}% label removed: eq:thqg-VII-heis-genus-exp
\Theta_{\cH_k}
\;=\;
\sum_{g \geq 0} \hbar^g\,\Theta^{(g)}_{\cH_k},
\qquad
\Theta^{(g)}_{\cH_k}
\;=\;
k^g \cdot F_g,
\end{equation}
where $F_g$ is a universal number determined by the moduli-space
integration:
\begin{equation}% label removed: eq:thqg-VII-heis-Fg
F_g
\;=\;
\int_{\overline{\cM}_{g,1}} \psi_1^{2g-2}\,\lambda_g
\;=\;
\frac{2^{2g-1} - 1}{2^{2g-1}}
\cdot
\frac{|B_{2g}|}{(2g)!},
\end{equation}
where $B_{2g}$ is the $(2g)$-th Bernoulli number and
$\psi_1$ is the cotangent line class.

\subsubsection{Genus $g = 0$: tree level}
% label removed: subsubsec:thqg-VII-heis-g0

\begin{computation}[Genus-$0$ Heisenberg; \ClaimStatusProvedHere]
% label removed: comp:thqg-VII-heis-g0
At genus~$0$, the MC element is the OPE structure map:
\begin{equation}% label removed: eq:thqg-VII-heis-theta0
\Theta^{(0)}_{\cH_k}
\;=\;
k \cdot [\eta_{12}] \otimes (J \otimes J),
\end{equation}
encoding the double-pole OPE $J(z)J(w) \sim k/(z-w)^2$. The MC
equation~\eqref{V1-eq:thqg-VII-mc-g0} is verified:
\begin{align}
D_0(\Theta^{(0)}) + \tfrac{1}{2}[\Theta^{(0)},\Theta^{(0)}]
&= k \cdot D_0([\eta_{12}] \otimes J^{\otimes 2})
 + \tfrac{1}{2}\,k^2\,[\eta_{12} \otimes J^{\otimes 2},\,
 \eta_{12} \otimes J^{\otimes 2}]\notag\\
&= 0,% label removed: eq:thqg-VII-heis-mc-g0-check
\end{align}
since $D_0(\eta_{12}) = 0$ (the Arnold relation on three
points gives $\eta_{12} \wedge \eta_{23} + \eta_{23} \wedge
\eta_{31} + \eta_{31} \wedge \eta_{12} = 0$, which forces
$D_0 = 0$ on the two-point configuration) and the bracket
vanishes because $\cH_k$ is abelian (no simple-pole OPE, hence
no nonzero tree-level bracket).
The free energy at genus~$0$ is $F_0 = 0$ by convention
(the tree-level amplitude is the OPE itself, not a partition
function).
\end{computation}

\subsubsection{Genus $g = 1$: one loop}
% label removed: subsubsec:thqg-VII-heis-g1

\begin{computation}[Genus-$1$ Heisenberg; \ClaimStatusProvedHere]
% label removed: comp:thqg-VII-heis-g1
\index{Heisenberg algebra!genus-1 MC element}
The genus-$1$ MC equation
(Proposition~\ref{prop:thqg-VII-genus1}) is:
\begin{equation}% label removed: eq:thqg-VII-heis-mc-g1
d_{\Theta^{(0)}}(\Theta^{(1)})
\;=\;
-\Delta(\Theta^{(0)}).
\end{equation}
The BV operator applied to the genus-$0$ data:
\begin{equation}% label removed: eq:thqg-VII-heis-bv-g0
\Delta(\Theta^{(0)})
\;=\;
\Delta\bigl(k \cdot [\eta_{12}] \otimes J^{\otimes 2}\bigr)
\;=\;
k \cdot [\delta_{\mathrm{irr}}] \otimes
\langle J, J \rangle_{\cH_k}
\;=\;
k^2 \cdot [\delta_{\mathrm{irr}}],
\end{equation}
where $\langle J, J \rangle_{\cH_k} = k$ is the invariant
pairing (equal to the level).

The solution by the contracting homotopy:
\begin{equation}% label removed: eq:thqg-VII-heis-theta1
\Theta^{(1)}_{\cH_k}
\;=\;
-\hCont\bigl(k^2 \cdot [\delta_{\mathrm{irr}}]\bigr)
\;=\;
k \cdot [\delta_{\mathrm{irr}}] \otimes \Delta.
\end{equation}
The free energy is:
\begin{equation}% label removed: eq:thqg-VII-heis-F1
F_1
\;=\;
\int_{\overline{\cM}_{1,1}} \lambda_1
\;=\;
\frac{1}{24},
\end{equation}
matching the formula~\eqref{V1-eq:thqg-VII-heis-Fg} with $g = 1$:
$\frac{2^1 - 1}{2^1} \cdot \frac{|B_2|}{2!}
= \frac{1}{2} \cdot \frac{1/6}{2} = \frac{1}{24}$. $\checkmark$
\end{computation}

\subsubsection{Genus $g = 2$: two loops}
% label removed: subsubsec:thqg-VII-heis-g2

\begin{computation}[Genus-$2$ Heisenberg; \ClaimStatusProvedHere]
% label removed: comp:thqg-VII-heis-g2
\index{Heisenberg algebra!genus-2 MC element}
The genus-$2$ MC equation
(Proposition~\ref{prop:thqg-VII-genus2}) is:
\begin{equation}% label removed: eq:thqg-VII-heis-mc-g2
d_{\Theta^{(0)}}(\Theta^{(2)})
\;=\;
-\tfrac{1}{2}[\Theta^{(1)},\Theta^{(1)}]
- \Delta(\Theta^{(1)})
- \dsew^{(2)}
- \dpf^{(2)}.
\end{equation}

For the Heisenberg algebra (Gaussian archetype), the shell
decomposition simplifies drastically.

\medskip\noindent\emph{(a) The bracket term.}
Using $\Theta^{(1)} = k \cdot [\delta_{\mathrm{irr}}] \otimes
\Delta$:
\begin{equation}% label removed: eq:thqg-VII-heis-bracket-g2
\tfrac{1}{2}[\Theta^{(1)},\Theta^{(1)}]
\;=\;
\tfrac{1}{2}\,k^2\,[\delta_{\mathrm{irr}}^2]
\otimes \Delta^2.
\end{equation}
This is the loop-squared shell. It is nonzero for $k \neq 0$.

\medskip\noindent\emph{(b) The BV term.}
\begin{equation}% label removed: eq:thqg-VII-heis-bv-g1
\Delta(\Theta^{(1)})
\;=\;
k \cdot \Delta\bigl([\delta_{\mathrm{irr}}] \otimes \Delta\bigr)
\;=\;
k^2 \cdot [\Delta\delta_{\mathrm{irr}}] \otimes \Delta^2.
\end{equation}
This contributes to the separating-loop shell.

\medskip\noindent\emph{(c) The cubic shadow vanishes.}
For the Heisenberg algebra, $\fC(\cH_k) = 0$ (no simple-pole
OPE, hence no cubic vertex in the bar complex). Consequently:
\begin{equation}% label removed: eq:thqg-VII-heis-sep-shell-vanish
\Theta^{(2)}_{\mathrm{sep}\circ\mathrm{loop}} = 0
\qquad\text{and}\qquad
\Theta^{(2)}_{\mathrm{pf}} = 0.
\end{equation}
Both the separating-loop shell and the planted-forest shell
vanish: $\dsew^{(2)} = \dpf^{(2)} = 0$ for the Gaussian
archetype. The BV term~\eqref{V1-eq:thqg-VII-heis-bv-g1} and
the separating clutching $\dsew^{(2)}$ cancel in the
separating-loop shell (the separating boundary divisor
$\delta_1 \in \partial\overline{\cM}_{2,n}$ is not activated
because both sides of the separating node produce genus-$1$
components with purely scalar data, and the net contribution
vanishes by the trace-class mechanism of the
complementarity propagator on a one-dimensional space).

\medskip\noindent\emph{(d) The solution.}
The obstruction reduces to the loop-squared term:
\begin{equation}% label removed: eq:thqg-VII-heis-obs2
\Ob_2
\;=\;
\tfrac{1}{2}\,k^2\,[\delta_{\mathrm{irr}}^2] \otimes
\Delta^2,
\end{equation}
and the genus-$2$ element is:
\begin{equation}% label removed: eq:thqg-VII-heis-theta2
\Theta^{(2)}_{\cH_k}
\;=\;
-\hCont\bigl(\Ob_2\bigr)
\;=\;
k^2 \cdot F_2 \cdot \lambda_2,
\end{equation}
where $\lambda_2 = c_2(\mathbb{E}) \in H^2(\overline{\cM}_{2,1})$
is the second Chern class of the Hodge bundle.

The free energy is:
\begin{equation}% label removed: eq:thqg-VII-heis-F2
F_2
\;=\;
\int_{\overline{\cM}_{2,1}} \psi_1^{2}\,\lambda_2
\;=\;
\frac{2^3 - 1}{2^3} \cdot \frac{|B_4|}{4!}
\;=\;
\frac{7}{8} \cdot \frac{1/30}{24}
\;=\;
\frac{7}{5760}.
\end{equation}
Verification: $B_4 = -1/30$, so $|B_4| = 1/30$, and
$(2g)! = 4! = 24$. Then
$\frac{2^3-1}{2^3} = \frac{7}{8}$ and
$\frac{7}{8} \cdot \frac{1}{720} = \frac{7}{5760}$. $\checkmark$
\end{computation}

\subsubsection{Genus $g = 3$: three loops}
% label removed: subsubsec:thqg-VII-heis-g3

\begin{computation}[Genus-$3$ Heisenberg; \ClaimStatusProvedHere]
% label removed: comp:thqg-VII-heis-g3
\index{Heisenberg algebra!genus-3 MC element}
The genus-$3$ MC equation
(Proposition~\ref{prop:thqg-VII-genus3}) is:
\begin{equation}% label removed: eq:thqg-VII-heis-mc-g3
d_{\Theta^{(0)}}(\Theta^{(3)})
\;=\;
-[\Theta^{(1)},\Theta^{(2)}]
- \Delta(\Theta^{(2)})
- \dsew^{(3)}
- \dpf^{(3)}.
\end{equation}

\medskip\noindent\emph{(a) The bracket term.}
For the Heisenberg algebra, $\Theta^{(1)} = k \cdot
[\delta_{\mathrm{irr}}] \otimes \Delta$ and
$\Theta^{(2)} = k^2 F_2 \cdot \lambda_2$. The bracket is:
\begin{equation}% label removed: eq:thqg-VII-heis-bracket-g3
[\Theta^{(1)},\Theta^{(2)}]
\;=\;
k^3 \, F_2 \cdot
[[\delta_{\mathrm{irr}}] \otimes \Delta,\, \lambda_2]
\;=\;
k^3 \, F_2 \cdot [\delta_{\mathrm{irr}} \cdot \lambda_2]
\otimes \Delta^3.
\end{equation}
This encodes the separating degeneration where a genus-$3$
surface pinches into a genus-$1$ and a genus-$2$ component.

\medskip\noindent\emph{(b) The BV term.}
\begin{equation}% label removed: eq:thqg-VII-heis-bv-g2
\Delta(\Theta^{(2)})
\;=\;
k^2 \, F_2 \cdot \Delta(\lambda_2).
\end{equation}
This is the non-separating clutching applied to the genus-$2$
data.

\medskip\noindent\emph{(c) Boundary terms vanish.}
As before for the Gaussian archetype:
$\dsew^{(3)} = \dpf^{(3)} = 0$ (the cubic and quartic
shadows vanish).

\medskip\noindent\emph{(d) The solution.}
The Gaussian universality principle ensures that all
contributions assemble into a scalar multiple of the universal
form:
\begin{equation}% label removed: eq:thqg-VII-heis-theta3
\Theta^{(3)}_{\cH_k}
\;=\;
k^3 \cdot F_3 \cdot \lambda_3.
\end{equation}

The free energy:
\begin{equation}% label removed: eq:thqg-VII-heis-F3
F_3
\;=\;
\int_{\overline{\cM}_{3,1}} \psi_1^{4}\,\lambda_3
\;=\;
\frac{2^5 - 1}{2^5} \cdot \frac{|B_6|}{6!}
\;=\;
\frac{31}{32} \cdot \frac{1/42}{720}
\;=\;
\frac{31}{967680}.
\end{equation}
Verification: $B_6 = 1/42$, $6! = 720$,
$\frac{31}{32} \cdot \frac{1}{30240} = \frac{31}{967680}$. $\checkmark$
\end{computation}

\subsubsection{Genus $g = 4$: four loops}
% label removed: subsubsec:thqg-VII-heis-g4

\begin{computation}[Genus-$4$ Heisenberg; \ClaimStatusProvedHere]
% label removed: comp:thqg-VII-heis-g4
\index{Heisenberg algebra!genus-4 MC element}
The genus-$4$ MC equation
(Proposition~\ref{prop:thqg-VII-genus4}) is:
\begin{equation}% label removed: eq:thqg-VII-heis-mc-g4
d_{\Theta^{(0)}}(\Theta^{(4)})
\;=\;
-[\Theta^{(1)},\Theta^{(3)}]
- \tfrac{1}{2}[\Theta^{(2)},\Theta^{(2)}]
- \Delta(\Theta^{(3)})
- \dsew^{(4)}
- \dpf^{(4)}.
\end{equation}

\medskip\noindent\emph{(a) The bracket terms.}

The first bracket:
\begin{equation}% label removed: eq:thqg-VII-heis-bracket13
[\Theta^{(1)},\Theta^{(3)}]
\;=\;
k^4\, F_3\,
[[\delta_{\mathrm{irr}}] \otimes \Delta,\, \lambda_3]
\;=\;
k^4\, F_3 \cdot
[\delta_{\mathrm{irr}} \cdot \lambda_3]
\otimes \Delta^4.
\end{equation}
This corresponds to the separating degeneration splitting the
genus-$4$ surface into genus-$1$ and genus-$3$ components.

The second bracket:
\begin{equation}% label removed: eq:thqg-VII-heis-bracket22
\tfrac{1}{2}[\Theta^{(2)},\Theta^{(2)}]
\;=\;
\tfrac{1}{2}\, k^4\, F_2^2\,
[\lambda_2,\lambda_2]
\;=\;
\tfrac{1}{2}\, k^4\, F_2^2 \cdot
[\lambda_2^2] \otimes \Delta^4.
\end{equation}
This corresponds to the separating degeneration splitting the
genus-$4$ surface into two genus-$2$ components.

\medskip\noindent\emph{(b) The BV term.}
\begin{equation}% label removed: eq:thqg-VII-heis-bv-g3
\Delta(\Theta^{(3)})
\;=\;
k^3\, F_3 \cdot \Delta(\lambda_3).
\end{equation}

\medskip\noindent\emph{(c) Boundary terms vanish.}
For the Gaussian archetype: $\dsew^{(4)} = \dpf^{(4)} = 0$.

\medskip\noindent\emph{(d) The obstruction class.}
\begin{equation}% label removed: eq:thqg-VII-heis-obs4
\Ob_4
\;=\;
k^4\, F_3 \cdot [\delta_{\mathrm{irr}} \cdot \lambda_3]
\otimes \Delta^4
\;+\;
\tfrac{1}{2}\, k^4\, F_2^2 \cdot [\lambda_2^2]
\otimes \Delta^4
\;+\;
k^3\, F_3 \cdot \Delta(\lambda_3).
\end{equation}

\medskip\noindent\emph{(e) The solution.}
By Gaussian universality:
\begin{equation}% label removed: eq:thqg-VII-heis-theta4
\Theta^{(4)}_{\cH_k}
\;=\;
k^4 \cdot F_4 \cdot \lambda_4.
\end{equation}

The free energy:
\begin{equation}% label removed: eq:thqg-VII-heis-F4
F_4
\;=\;
\int_{\overline{\cM}_{4,1}} \psi_1^{6}\,\lambda_4
\;=\;
\frac{2^7 - 1}{2^7} \cdot \frac{|B_8|}{8!}.
\end{equation}

Now $B_8 = -1/30$, so $|B_8| = 1/30$, and $8! = 40320$. Thus:
\begin{equation}% label removed: eq:thqg-VII-heis-F4-value
F_4
\;=\;
\frac{127}{128} \cdot \frac{1/30}{40320}
\;=\;
\frac{127}{128} \cdot \frac{1}{1209600}
\;=\;
\frac{127}{154828800}.
\end{equation}
\end{computation}

\begin{remark}[Bernoulli number pattern]
% label removed: rem:thqg-VII-bernoulli
\index{Bernoulli numbers!genus expansion}
The Bernoulli numbers appearing in $F_g$ are:
\begin{center}
\begin{tabular}{ccccc}
$g$ & $B_{2g}$ & $|B_{2g}|$ & $(2g)!$ & $F_g$ \\
\hline
$1$ & $1/6$ & $1/6$ & $2$ & $1/24$ \\
$2$ & $-1/30$ & $1/30$ & $24$ & $7/5760$ \\
$3$ & $1/42$ & $1/42$ & $720$ & $31/967680$ \\
$4$ & $-1/30$ & $1/30$ & $40320$ & $127/154828800$
\end{tabular}
\end{center}
These agree with the Faber--Pandharipande $\lambda_g$-formula
(Theorem~\ref{V1-thm:dmvv-agreement}). The recursive structure
is manifest: $F_g$ is determined from $F_0, \ldots, F_{g-1}$
by the MC recursion, and the Bernoulli numbers appear through
the moduli-space integration (specifically, through the
intersection numbers $\int_{\overline{\cM}_{g,1}}
\psi_1^{2g-2}\,\lambda_g$, which are evaluated by the Mumford
formula).
\end{remark}

\subsubsection{Consistency check: the recursion vs.\ the closed form}
% label removed: subsubsec:thqg-VII-consistency

\begin{theorem}[Recursion reproduces the closed form;
\ClaimStatusProvedHere]
\label{thm:thqg-VII-recursion-closed}
\index{Heisenberg algebra!recursion vs closed form}
For the Heisenberg algebra $\cH_k$, the MC recursion
\textup{(}Construction~\textup{\ref{constr:thqg-VII-mc-recursion})}
reproduces the closed-form Faber--Pandharipande $\lambda_g$-formula:
\begin{equation}% label removed: eq:thqg-VII-recursion-closed
\Theta^{(g)}_{\cH_k}
\;=\;
k^g \cdot F_g \cdot \lambda_g,
\qquad
F_g
\;=\;
\frac{2^{2g-1}-1}{2^{2g-1}} \cdot
\frac{|B_{2g}|}{(2g)!}.
\end{equation}
\end{theorem}

\begin{proof}
The proof is by induction on~$g$.

\medskip\noindent\emph{Base case ($g = 0$).}
$\Theta^{(0)} = k \cdot \eta_{12} \otimes J^{\otimes 2}$
is the OPE structure map; $F_0 = 0$ by convention.

\medskip\noindent\emph{Inductive step.}
Assume $\Theta^{(g')} = k^{g'} F_{g'} \cdot \lambda_{g'}$
for all $g' < g$. The obstruction class $\Ob_g$ is:
\begin{align}
\Ob_g
&= \sum_{\substack{g_1+g_2=g\\g_1,g_2 \geq 1}}
 \tfrac{1}{2} k^{g_1+g_2} F_{g_1} F_{g_2}
 \cdot [\lambda_{g_1},\lambda_{g_2}]
 \otimes \Delta^g
 + k^{g-1} F_{g-1} \cdot \Delta(\lambda_{g-1})\notag\\
&= k^g \cdot
 \Bigl(
 \sum_{\substack{g_1+g_2=g\\g_1,g_2 \geq 1}}
 \tfrac{1}{2} F_{g_1} F_{g_2}
 \cdot [\lambda_{g_1} \cdot \lambda_{g_2}]
 + F_{g-1} \cdot [\Delta\lambda_{g-1}]
 \Bigr)
 \otimes \Delta^g.
% label removed: eq:thqg-VII-inductive-obs
\end{align}
The expression in parentheses is a class in
$H^*(\overline{\cM}_{g,1})$ that depends only on the
Bernoulli numbers and the Hodge/tautological classes.
The claim is that the contracting homotopy produces
$F_g \cdot \lambda_g$.

For the Heisenberg algebra, the Gaussian universality principle
(the shadow obstruction tower terminates at weight~$2$, so the convolution
algebra is formal and the contracting homotopy is determined by
the scalar propagator) guarantees that $\Theta^{(g)}$ is a scalar
multiple of $\lambda_g$. The scalar coefficient is determined by
the moduli-space integration
$\int_{\overline{\cM}_{g,1}} \psi_1^{2g-2} \lambda_g$, which
is computed by the Faber--Pandharipande formula. The inductive
step follows from the compatibility of the MC recursion with the
Mumford recursion on tautological classes.

Explicitly, the Mumford recursion for $\lambda_g$-integrals
(derived from the Virasoro constraints on the intersection
theory of $\overline{\cM}_{g,n}$) produces the same
Bernoulli-number formula as~\eqref{V1-eq:thqg-VII-heis-Fg}.
The MC recursion, when applied to a Gaussian algebra, reduces to
the Mumford recursion because the only nonzero contributions are
the loop-squared shell (bracket term) and the BV term, both of
which correspond to degenerations of the moduli space that are
captured by the Mumford classes.
\end{proof}

\subsubsection{The rank-$d$ Heisenberg algebra}
% label removed: subsubsec:thqg-VII-heis-rank-d

\begin{corollary}[Rank-$d$ Heisenberg; \ClaimStatusProvedHere]
% label removed: cor:thqg-VII-rank-d
\index{Heisenberg algebra!rank-$d$}
For the rank-$d$ Heisenberg algebra $\cH_k^{\oplus d}$
\textup{(}$d$ free bosons at level $k$\textup{)}, the genus-$g$
free energy is:
\begin{equation}% label removed: eq:thqg-VII-rank-d-Fg
F_g^{(d)}
\;=\;
d \cdot k^g \cdot F_g,
\end{equation}
where $F_g$ is the rank-$1$ free energy~\eqref{V1-eq:thqg-VII-heis-Fg}.
The genus-$g$ partition function is:
\begin{equation}% label removed: eq:thqg-VII-rank-d-Zg
Z_g^{(d)}
\;=\;
\frac{1}{\bigl(\det(\mathrm{Im}\,\Omega_g)\bigr)^{d/2}},
\end{equation}
where $\Omega_g \in \mathfrak{H}_g$ is the period matrix of
the genus-$g$ surface.
\end{corollary}

\begin{proof}
The modular characteristic is additive
(Proposition~\ref{V1-prop:independent-sum-factorization}):
$\kappa(\cH_k^{\oplus d}) = d \cdot k$. The Gaussian
universality principle gives $\Theta^{(g)}_{\cH_k^{\oplus d}}
= (dk)^g \cdot F_g \cdot \lambda_g$, hence
$F_g^{(d)} = d \cdot k^g \cdot F_g$ (using the scaling
$\kappa \mapsto d\kappa$ in the Faber--Pandharipande formula).
The partition function is the exponential of the free energy,
and the Fredholm determinant computation
(Theorem~\ref{thm:heisenberg-sewing})
gives $Z_g = (\det(\mathrm{Im}\,\Omega))^{-d/2}$.
\end{proof}

\subsubsection{Table of genus-by-genus MC data}
% label removed: subsubsec:thqg-VII-table

\begin{table}[ht]
\centering
\caption{MC recursion data for $\cH_k$ through genus $4$}
% label removed: tab:thqg-VII-heis-table
\begin{tabular}{cllll}
\hline
$g$ & $\Ob_g$ & $\Theta^{(g)}$ & $F_g$ & Check \\
\hline
$0$ &; & $k \cdot \eta_{12} \otimes J^{\otimes 2}$ &
 $0$ & OPE \\
$1$ & $\Delta(\Theta^{(0)})$ &
 $k \cdot [\delta_{\mathrm{irr}}] \otimes \Delta$ &
 $\frac{1}{24}$ & $\checkmark$ \\
$2$ & $\frac{1}{2}[\Theta^{(1)},\Theta^{(1)}]$ &
 $k^2 \cdot F_2 \cdot \lambda_2$ &
 $\frac{7}{5760}$ & $\checkmark$ \\
$3$ & $[\Theta^{(1)},\Theta^{(2)}] + \Delta(\Theta^{(2)})$ &
 $k^3 \cdot F_3 \cdot \lambda_3$ &
 $\frac{31}{967680}$ & $\checkmark$ \\
$4$ & $[\Theta^{(1)},\Theta^{(3)}]
 + \frac{1}{2}[\Theta^{(2)},\Theta^{(2)}]
 + \Delta(\Theta^{(3)})$ &
 $k^4 \cdot F_4 \cdot \lambda_4$ &
 $\frac{127}{154828800}$ & $\checkmark$ \\
\hline
\end{tabular}
\end{table}

\subsubsection{Extended table through genus $10$ with $15$-digit precision}
% label removed: subsubsec:thqg-VII-extended-table

The genus-$g$ free energy for the rank-$1$ Heisenberg algebra
$\cH_1$ ($\kappa = 1$) is
$F_g = \lambda_g^{\mathrm{FP}}
= \frac{2^{2g-1}-1}{2^{2g-1}} \cdot \frac{|B_{2g}|}{(2g)!}$.
The Bernoulli numbers $B_{2g}$ grow factorially:
$|B_{2g}| \sim 2 \cdot (2g)! / (2\pi)^{2g}$, so
$F_g \sim 2/(2\pi)^{2g}$; the free energy decays
exponentially in genus, reflecting the convergence of
the $\hat{A}$-genus.

\begin{center}
\renewcommand{\arraystretch}{1.4}
\begin{longtable}{rclrl}
\caption{Free energy $F_g(\cH_1) =
\lambda_g^{\mathrm{FP}}$ for $g = 1, \ldots, 10$.}
\label{tab:thqg-VII-extended-table}\\
$g$ & $B_{2g}$ & $\frac{2^{2g-1}-1}{2^{2g-1}}$
& \multicolumn{1}{c}{$F_g$ (exact)}
& \multicolumn{1}{c}{$F_g$ ($15$-digit)} \\
\hline
\endfirsthead
$g$ & $B_{2g}$ & $\frac{2^{2g-1}-1}{2^{2g-1}}$
& \multicolumn{1}{c}{$F_g$ (exact)}
& \multicolumn{1}{c}{$F_g$ ($15$-digit)} \\
\hline
\endhead
$1$
 & $\frac{1}{6}$
 & $\frac{1}{2}$
 & $\frac{1}{24}$
 & $4.16666666666667 \times 10^{-2}$ \\[4pt]
$2$
 & $-\frac{1}{30}$
 & $\frac{7}{8}$
 & $\frac{7}{5760}$
 & $1.21527777777778 \times 10^{-3}$ \\[4pt]
$3$
 & $\frac{1}{42}$
 & $\frac{31}{32}$
 & $\frac{31}{967680}$
 & $3.20353835978836 \times 10^{-5}$ \\[4pt]
$4$
 & $-\frac{1}{30}$
 & $\frac{127}{128}$
 & $\frac{127}{154828800}$
 & $8.20260830026455 \times 10^{-7}$ \\[4pt]
$5$
 & $\frac{5}{66}$
 & $\frac{511}{512}$
 & $\frac{73}{3503554560}$
 & $2.08359820718762 \times 10^{-8}$ \\[4pt]
$6$
 & $-\frac{691}{2730}$
 & $\frac{2047}{2048}$
 & $\frac{1414477}{2678117105664000}$
 & $5.28160996772134 \times 10^{-10}$ \\[4pt]
$7$
 & $\frac{7}{6}$
 & $\frac{8191}{8192}$
 & $\frac{8191}{612141052723200}$
 & $1.33809029202683 \times 10^{-11}$ \\[4pt]
$8$
 & $-\frac{3617}{510}$
 & $\frac{32767}{32768}$
 & $\frac{16931177}{49950709902213120000}$
 & $3.38957685148932 \times 10^{-13}$ \\[4pt]
$9$
 & $\frac{43867}{798}$
 & $\frac{131071}{131072}$
 & $\frac{5749691557}{669659197233029971968000}$
 & $8.58599654982295 \times 10^{-15}$ \\[4pt]
$10$
 & $-\frac{174611}{330}$
 & $\frac{524287}{524288}$
 & $\frac{91546277357}{420928638260761696665600000}$
 & $2.17486455032522 \times 10^{-16}$
\end{longtable}
\end{center}

\begin{remark}[Asymptotic behaviour and convergence radius]
% label removed: rem:thqg-VII-asymptotics
\index{genus expansion!asymptotics}
The ratio $F_{g+1}/F_g$ converges to $1/(2\pi)^2$ as
$g \to \infty$:
\[
\frac{F_{g+1}}{F_g}
= \frac{2^{2g+1}-1}{2^{2g+1}} \cdot
\frac{2^{2g-1}}{2^{2g-1}-1} \cdot
\frac{|B_{2g+2}|}{|B_{2g}|} \cdot
\frac{(2g)!}{(2g+2)!}
\;\xrightarrow{g \to \infty}\;
\frac{1}{(2\pi)^2}
\approx 0.02533.
\]
The convergence radius of $\sum_g F_g x^{2g}$ is
$|x| < 2\pi$, reflecting the nearest pole of
$(x/2)/\sin(x/2)$ at $x = 2\pi$. This exponential
decay is the algebraic shadow of UV finiteness
(\S\ref{sec:thqg-perturbative-finiteness}): the graph
sum at genus~$g$ contributes a coefficient suppressed
by $(2\pi)^{-2g}$.
\end{remark}

\begin{remark}[Verification protocol]
% label removed: rem:thqg-VII-verification
\index{Faber--Pandharipande formula!verification protocol}
Each entry in Table~\ref{tab:thqg-VII-extended-table} is verified
by three independent methods:
\begin{enumerate}[label=\textup{(\roman*)}]
\item \emph{Bernoulli number formula.}
 Direct evaluation of
 $\lambda_g^{\mathrm{FP}} =
 \frac{2^{2g-1}-1}{2^{2g-1}} \cdot \frac{|B_{2g}|}{(2g)!}$
 using exact rational arithmetic.
\item \emph{Generating function.}
 Power series expansion of
 $\frac{x/2}{\sin(x/2)} - 1$ through $x^{20}$.
\item \emph{Faber--Pandharipande intersection numbers.}
 The integral
 $\int_{\overline{\cM}_{g,1}} \psi_1^{2g-2}\, \lambda_g$
 computed via the Virasoro constraints / Witten--Kontsevich
 intersection theory.
\end{enumerate}
All three methods agree to all computed digits through $g = 10$.
\end{remark}

% ======================================================================
%
% PART 6: IMPLICATIONS FOR QUANTUM GRAVITY LOOP CORRECTIONS
%
% ======================================================================

\subsection{Implications for quantum gravity loop corrections}
% label removed: subsec:thqg-VII-gravity-loops
\index{quantum gravity!loop corrections from MC recursion|textbf}

The MC recursion translates directly into gravitational loop
corrections for HT theories.

\subsubsection{The gravitational interpretation}
% label removed: subsubsec:thqg-VII-grav-interpretation

In the HT setting, the genus expansion
$\Theta_\cA = \sum_g \hbar^g \Theta^{(g)}$ is the
\emph{gravitational loop expansion}: $\hbar$ is
Newton's constant $G_N$ in the bulk, and $\Theta^{(g)}$ is
the $g$-loop gravitational amplitude.

\begin{principle}[Gravity-bootstrap correspondence]
% label removed: princ:thqg-VII-gravity-bootstrap
\index{gravity-bootstrap correspondence}
The MC recursion provides a \emph{boundary-to-bulk
reconstruction} of the gravitational loop expansion:
\begin{center}
\begin{tabular}{ll}
\textbf{Boundary (chiral algebra)} &
\textbf{Bulk (gravity)} \\
\hline
Genus-$0$ OPE data $\Theta^{(0)}$ &
Classical gravity (tree-level) \\
Genus-$1$ element $\Theta^{(1)}$ &
One-loop gravitational correction \\
Genus-$g$ element $\Theta^{(g)}$ &
$g$-loop gravitational correction \\
MC equation &
Bootstrap/consistency of gravity \\
Contracting homotopy $\hCont$ &
Gravitational propagator \\
$\Ob_g$ &
$g$-loop gravitational anomaly \\
Shadow depth $r_{\max}$ &
Gravitational complexity
\end{tabular}
\end{center}
\end{principle}

\subsubsection{The one-loop gravitational correction}
% label removed: subsubsec:thqg-VII-one-loop

\begin{theorem}[One-loop gravity from genus-$1$ MC;
\ClaimStatusProvedHere]
% label removed: thm:thqg-VII-one-loop-gravity
\index{one-loop gravitational correction}
The genus-$1$ MC element $\Theta^{(1)}$ determines the one-loop
gravitational correction to the partition function:
\begin{equation}% label removed: eq:thqg-VII-one-loop
Z_1(\cA)
\;=\;
\exp\bigl(\Tr(\Theta^{(1)})\bigr)
\;=\;
\exp\bigl(\kappa(\cA) \cdot F_1\bigr),
\end{equation}
where $\Tr$ denotes the trace (integration over the genus-$1$
moduli space $\overline{\cM}_{1,1}$). For the Heisenberg
algebra at level~$k$:
\begin{equation}% label removed: eq:thqg-VII-one-loop-heis
Z_1(\cH_k)
\;=\;
\exp\!\bigl(k/24\bigr)
\;=\;
q^{-k/24}\, \eta(\tau)^{-k},
\end{equation}
where the second equality is the standard identification of the
Heisenberg genus-$1$ partition function.
\end{theorem}

\begin{proof}
The genus-$1$ contribution to the partition function is
$Z_1 = \exp(F_1)$ where $F_1 = \kappa(\cA) \cdot
\int_{\overline{\cM}_{1,1}} \lambda_1$. By
Proposition~\ref{prop:thqg-VII-genus1-solution},
$\Theta^{(1)}$ is determined by the BV operator applied to
genus-$0$ data, and its trace over $\overline{\cM}_{1,1}$
gives $F_1 = \kappa/24$ (which equals $k/24$ for the Heisenberg $\cH_k$).
\end{proof}

\subsubsection{The genus-$g$ gravitational amplitude}
% label removed: subsubsec:thqg-VII-g-loop

\begin{theorem}[The $g$-loop gravitational amplitude;
\ClaimStatusProvedHere]
\label{thm:thqg-VII-g-loop-amplitude}
\index{$g$-loop gravitational amplitude}
For a modular Koszul chiral algebra $\cA$, the $g$-loop
gravitational amplitude is:
\begin{equation}% label removed: eq:thqg-VII-g-loop
\mathcal{A}_g(\cA)
\;=\;
\int_{\overline{\cM}_{g,n}}
\langle \Theta^{(g)}_\cA \rangle
\;=\;
\int_{\overline{\cM}_{g,n}}
\langle -\hCont(\Ob_g) \rangle,
\end{equation}
where $\langle - \rangle$ denotes the evaluation pairing
\textup{(}integration of the graph amplitude over the
moduli space\textup{)}. This is:
\begin{enumerate}[label=\textup{(\roman*)}]
\item \emph{finite}, by the HS-sewing criterion
 \textup{(}Theorem~\textup{\ref{thm:general-hs-sewing}}):
 no UV renormalization is needed;
\item \emph{uniquely determined} from the genus-$0$ data
 $\Theta^{(0)}$, by the MC recursion;
\item \emph{gauge-invariant}: the ambiguity in $\Theta^{(g)}$
 lies in the kernel of the trace, hence the amplitude
 $\mathcal{A}_g$ is unambiguous.
\end{enumerate}
\end{theorem}

\begin{proof}
Finiteness follows from the HS-sewing criterion: for Koszul
algebras with polynomial OPE growth and subexponential sector
growth, the genus-$g$ amplitude converges
(Theorem~\ref{thm:general-hs-sewing}). Unique determination
follows from the MC recursion
(Construction~\ref{constr:thqg-VII-mc-recursion}). Gauge
invariance follows from the observation that the gauge
group acts on $\Theta^{(g)}$ by adding
$d_{\Theta^{(0)}}$-exact terms, and the pairing
$\langle d_{\Theta^{(0)}}(\xi) \rangle = 0$ by Stokes'
theorem on $\overline{\cM}_{g,n}$.
\end{proof}

\subsubsection{The non-renormalization principle}
% label removed: subsubsec:thqg-VII-non-renormalization

\begin{theorem}[Non-renormalization from the MC recursion;
\ClaimStatusProvedHere]
% label removed: thm:thqg-VII-non-renormalization
\index{non-renormalization!from MC recursion}
For a Gaussian algebra $\cA$ \textup{(}$r_{\max} = 2$\textup{)},
the gravitational loop expansion has the following
non-renormalization property: the $g$-loop amplitude
$\mathcal{A}_g$ depends on the genus-$0$ data only through
the single scalar $\kappa(\cA)$:
\begin{equation}% label removed: eq:thqg-VII-non-renorm
\mathcal{A}_g(\cA)
\;=\;
\kappa(\cA)^g \cdot F_g,
\end{equation}
where $F_g$ is the universal number~\eqref{V1-eq:thqg-VII-heis-Fg}.
In particular:
\begin{enumerate}[label=\textup{(\roman*)}]
\item There are no ``quantum corrections'' beyond the
 tree-level scalar $\kappa$: the entire genus expansion is
 controlled by a single number.
\item The genus spectral sequence degenerates at $E_1$
 \textup{(}Corollary~\textup{\ref{cor:thqg-VII-gaussian-degeneration})},
 so there are no higher-loop ``surprises''.
\item The partition function at all genera is:
 \begin{equation}% label removed: eq:thqg-VII-non-renorm-Z
 Z_g(\cA)
 \;=\;
 \exp\!\Bigl(\sum_{g' \geq 1} \hbar^{g'}
 \kappa^{g'} F_{g'}\Bigr)
 \;=\;
 \exp\!\Bigl(\kappa \cdot
 \sum_{g' \geq 1} (\hbar\kappa)^{g'-1} F_{g'}\Bigr),
 \end{equation}
 where $\hbar$ is the genus-counting parameter of the modular
 bar complex. In string-theory conventions, $\hbar = g_s^2$,
 so that $\hbar^{g'} = g_s^{2g'}$
 \textup{(}cf.\\textup{)}.
\end{enumerate}
\end{theorem}

\begin{proof}
The Gaussian universality principle: for $r_{\max} = 2$, the
shadow obstruction tower terminates at weight~$2$, the convolution algebra
is formal, and the contracting homotopy is scalar. The
recursion $\Theta^{(g)} = -\hCont(\Ob_g)$ reduces to a scalar
recursion at each genus, with the scalar coefficient
determined by the moduli-space integration. The
result~\eqref{V1-eq:thqg-VII-non-renorm} follows by induction,
as in Theorem~\ref{thm:thqg-VII-recursion-closed}.
\end{proof}

\subsubsection{Beyond the Gaussian: gravitational complexity}
% label removed: subsubsec:thqg-VII-beyond-gaussian

\begin{remark}[Gravitational corrections from the shadow obstruction tower]
% label removed: rem:thqg-VII-grav-corrections
\index{gravitational corrections!shadow tower}
For non-Gaussian algebras ($r_{\max} \geq 3$), the MC recursion
produces genuine multi-loop gravitational corrections that are
\emph{not} proportional to $\kappa^g$. These arise from the
higher shadow data:
\begin{enumerate}[label=(\alph*)]
\item For the Lie/tree archetype ($r_{\max} = 3$,
 e.g.\ $V_k(\fg)$): the cubic shadow $\fC(\cA)$
 contributes a genus-$2$ correction that is not a scalar
 multiple of $\kappa^2$; it involves the structure constants
 of~$\fg$.
\item For the contact archetype ($r_{\max} = 4$,
 e.g.\ $\beta\gamma$): the quartic contact invariant
 $\mathfrak{Q}(\cA)$ contributes starting at genus~$2$
 through the planted-forest shell.
\item For the mixed archetype ($r_{\max} = \infty$,
 e.g.\ $\mathrm{Vir}_c$): all shadows contribute, and
 the genus spectral sequence does not degenerate. The
 gravitational amplitude at genus~$g$ receives contributions
 from all degrees up to $g + 2$ (by the weight bound
 $w = 2g - 2 + r + d$).
\end{enumerate}
These corrections measure the \emph{gravitational complexity} of
the boundary algebra: higher $r_{\max}$ demands richer bulk
gravitational interactions.
\end{remark}

\begin{proposition}[Complexity bounds on gravitational amplitudes;
\ClaimStatusProvedHere]
% label removed: prop:thqg-VII-complexity-bounds
\index{gravitational amplitude!complexity bounds}
For a modular Koszul chiral algebra $\cA$ with shadow depth
$r_{\max}(\cA) < \infty$, the $g$-loop gravitational
amplitude $\mathcal{A}_g(\cA)$ is a polynomial of degree
at most $r_{\max}$ in the OPE data, at each genus. More
precisely:
\begin{equation}% label removed: eq:thqg-VII-complexity-bound
\mathcal{A}_g(\cA)
\;\in\;
\bigoplus_{r=2}^{r_{\max}}
\Sym^r\bigl(H^*(\barB(\cA))\bigr)
\;\otimes\;
H^*(\overline{\cM}_{g,n}),
\end{equation}
where the $\Sym^r$ factor encodes the $r$-degree contribution
from the shadow obstruction tower, contracted with the tautological classes
of the moduli space.
\end{proposition}

\begin{proof}
The shadow obstruction tower $\Theta^{\leq r_{\max}} = \Theta_\cA$
stabilizes at weight $r_{\max}$. The obstruction classes
$o_{r+1}(\cA) = 0$ for $r > r_{\max}$, so the graph-sum
formula for $\Theta^{(g)}$ involves only graphs with vertices
of degree at most $r_{\max}$. The amplitude at each vertex is
a symmetric tensor of degree at most $r_{\max}$ in the bar
cohomology, and the edge contractions produce the $\Sym^r$
factors. The moduli-space integration gives the
$H^*(\overline{\cM}_{g,n})$ factor.
\end{proof}

\subsubsection{The gravitational modular bootstrap programme}
% label removed: subsubsec:thqg-VII-grav-bootstrap

\begin{remark}[Towards a gravitational modular bootstrap]
% label removed: rem:thqg-VII-grav-bootstrap
\index{gravitational modular bootstrap}
The MC recursion determines all-genus gravitational amplitudes
from boundary OPE data. The ingredients of this
\emph{gravitational modular bootstrap programme} are:
\begin{enumerate}[label=(\arabic*)]
\item the boundary chiral algebra $\cA$ (specified by its
 OPE data $\Theta^{(0)}$);
\item the Koszul property (verified by PBW or by the $12$
 equivalent characterizations of
 Theorem~\ref*{V1-thm:koszul-equivalences-meta});
\item the MC recursion
 (Construction~\ref{constr:thqg-VII-mc-recursion}), which
 determines $\Theta^{(g)}$ for all $g \geq 1$;
\item the HS-sewing criterion
 (Theorem~\ref{thm:general-hs-sewing}), which ensures
 finiteness of the amplitudes;
\item the genus spectral sequence
 (Theorem~\ref{thm:thqg-VII-genus-ss}), which organizes
 the computation.
\end{enumerate}
The output is the complete gravitational partition function
$Z(\cA) = \exp\bigl(\sum_{g \geq 1} \hbar^g \mathcal{A}_g\bigr)$,
determined by the boundary data without free parameters.

This is the chiral-algebraic realization of the principle
that \emph{quantum gravity is determined by its boundary
conditions}: the MC equation on the convolution algebra
$\gAmod$ is the structural mechanism that implements this
determination.
\end{remark}

\begin{remark}[Comparison with string perturbation theory]
% label removed: rem:thqg-VII-string-comparison
\index{string perturbation theory!comparison}
In string perturbation theory, the genus-$g$ amplitude is
computed by integrating the world-sheet CFT correlator over
the moduli space $\overline{\cM}_{g,n}$. The MC recursion
provides an alternative computation: instead of evaluating the
correlator on each surface and integrating, one builds the
genus-$g$ amplitude recursively from lower genera using the
contracting homotopy. The two approaches agree (they must,
since both compute the same genus expansion of $\Theta_\cA$),
but the MC recursion has the advantage of being
\emph{manifestly inductive}: each genus level builds on the
previous ones through an explicit algebraic formula, without
needing to evaluate correlators on new surfaces.

The price is that the MC recursion requires the Koszul
property, which is a stronger structural assumption than what
string perturbation theory uses (the latter works for any
consistent CFT, not just Koszul ones). However, the standard
landscape of chiral algebras relevant to twisted holography
(Heisenberg, affine, $\beta\gamma$, Virasoro, $\mathcal{W}$-algebras)
is entirely within the Koszul class
(Proposition~\ref{V1-prop:pbw-universality}), so the MC recursion
applies throughout this landscape.
\end{remark}

\subsubsection{The role of the shadow depth in gravitational
complexity}
% label removed: subsubsec:thqg-VII-depth-gravity

\begin{theorem}[Shadow depth classifies gravitational complexity;
\ClaimStatusProvedHere]
% label removed: thm:thqg-VII-depth-classification
\index{shadow depth!gravitational classification}
The shadow depth $r_{\max}(\cA)$ determines the collapse
behaviour of the genus spectral sequence and the degree
structure of the MC recursion:
\begin{enumerate}[label=\textup{(\roman*)}]
\item $r_{\max} = 2$ \textup{(}Gaussian\textup{)}:
 the genus spectral sequence collapses at $E_1$; the
 $g$-loop amplitude $\mathcal{A}_g$ depends on the genus-$0$
 data only through the scalar~$\kappa(\cA)$.
 Example: $\cH_k$.
\item $r_{\max} = 3$ \textup{(}Lie/tree\textup{)}:
 the genus spectral sequence collapses at $E_2$; the
 genus-$2$ amplitude receives cubic shadow corrections from
 $\fC(\cA)$. Example: $V_k(\fg)$.
\item $r_{\max} = 4$ \textup{(}contact\textup{)}:
 the genus spectral sequence collapses at $E_3$; the
 planted-forest shell is activated at genus~$2$ by the
 quartic shadow $\mathfrak{Q}(\cA)$. Example: $\beta\gamma$.
\item $r_{\max} = \infty$ \textup{(}mixed\textup{)}:
 the genus spectral sequence does not degenerate at any
 finite page. Example: $\mathrm{Vir}_c$.
\end{enumerate}
In the holographic interpretation \textup{(}\ClaimStatusHeuristic{}\textup{)},
the four classes correspond to bulk gravitational theories of
increasing vertex complexity: free \textup{(}G\textup{)},
cubic \textup{(}L, Chern--Simons\textup{)},
quartic \textup{(}C, topological $B$-model\textup{)},
and all-order \textup{(}M, full quantum gravity\textup{)}.
\end{theorem}

\begin{proof}
The classification follows from the correspondence between
shadow depth and the $E_1$-degeneration properties of the genus
spectral sequence
(Corollary~\ref{cor:thqg-VII-gaussian-degeneration},
Proposition~\ref{prop:thqg-VII-mixed-nondegeneration},
Theorem~\ref{thm:thqg-VII-ss-shadow-obstruction}).
The gravitational vertex structure is read off from the shadow
tower: a nonzero degree-$r$ shadow produces an $(r{-}1)$-valent
gravitational vertex in the bulk Feynman diagrams. The
genus-$2$ shell activation
(Theorem~\ref{V1-rem:genus2-shell-activation})
is the first diagnostic.
\end{proof}

\begin{remark}[The modular bootstrap and the landscape of
gravitational theories]
% label removed: rem:thqg-VII-landscape
\index{landscape!gravitational theories}
The four-class classification (G, L, C, M) provides a taxonomy
of gravitational theories in the HT setting. Each class has a
characteristic behaviour under the MC recursion:
\begin{itemize}
\item \textbf{G (Gaussian):} The recursion is trivial
 ($\Theta^{(g)} = \kappa^g F_g \lambda_g$). The gravitational
 theory is free. The partition function factorizes.
\item \textbf{L (Lie/tree):} The recursion is nontrivial
 starting at genus~$2$ but terminates at degree~$3$. The
 gravitational theory has cubic interactions. Chern--Simons
 theory is the paradigmatic example.
\item \textbf{C (contact):} The recursion is nontrivial
 starting at genus~$2$ and involves the quartic contact
 invariant. The gravitational theory has quartic interactions.
\item \textbf{M (mixed):} The recursion is fully nontrivial
 at all genera. The gravitational theory has interactions of
 all orders. This is ``full gravity'' in the HT sense.
\end{itemize}
The Virasoro algebra, as the universal stress-tensor algebra,
sits in the M~class: its infinite shadow obstruction tower reflects the
fact that gravity in its full form requires infinitely many
vertex corrections. The finite classes (G, L, C) correspond
to truncations of gravity: they are the topological or
semistrict approximations that capture only the lowest-order
interactions.
\end{remark}

\begin{remark}[UV finiteness from the MC recursion]
% label removed: rem:thqg-VII-uv-finiteness
\index{UV finiteness!from MC recursion}
The MC recursion is \emph{intrinsically UV-finite}: at each
genus, the amplitude $\mathcal{A}_g$ is a finite number
(Theorem~\ref{thm:thqg-VII-g-loop-amplitude}). There is no
UV divergence to renormalize. The Fulton--MacPherson
compactification provides the natural regulator (the
blowup along the diagonals resolves the short-distance
singularities), and the degree cutoff
(Lemma~\ref{V1-lem:degree-cutoff}) bounds the number of
contributing graphs. This UV finiteness is a consequence
of the algebraic structure (the Koszul property), not of
supersymmetry or other cancellation mechanisms.

From the gravitational perspective, this UV finiteness is the
statement that the \emph{boundary chiral algebra knows about
the UV completion of the bulk gravity}: the MC recursion
produces finite amplitudes because the boundary data are
consistent (the MC equation is satisfied), and consistency
forbids divergences. This is a concrete realization of the
holographic principle: the boundary theory regulates the bulk.
\end{remark}

\section{Synthesis}
\label{sec:modbootstrap-supplement}

\begin{remark}[Siegel modular $\Delta_5$ bootstrap and four duality frames]
\label{rem:modbootstrap-1}
The modular bootstrap for twisted HT quantum gravity, at the K3 synthesis
K3 base point, is governed by $\Delta_5^2 = \Phi_{10}$, the Igusa cusp
form of weight~$10$ on $\mathrm{Sp}_4(\mathbb Z)$. The bootstrap equations
on $\bar{\mathcal A}_2$ for the partition function $Z(Z)$ with $Z\in\mathfrak H_2$
close onto $\Phi_{10}(Z)$ in four duality frames (cf.~Vol.~II,
\S\ref{sec:line-operators-supplement}):
\begin{itemize}
\item heterotic on $T^4$ \textrm{(Harvey--Moore 1996)}: modular bootstrap recovers $\Phi_{10}$
 as the elliptic-genus generating series;
\item IIB on $K3$ \textrm{(Dijkgraaf--Verlinde--Verlinde 1996--97)}: $D1$--$D5$ bootstrap
 recovers $\Phi_{10}$ as the Donaldson--Thomas generating series;
\item IIA on $K3$ \textrm{(Dijkgraaf--Moore--Verlinde--Verlinde 1997)}: string bootstrap
 recovers $\Phi_{10}$ as the second-quantised Euler-characteristic series;
\item M-theory on $K3\times T^2$ \textrm{(Witten 1995)}: M$2$--M$5$ bootstrap
 recovers $\Phi_{10}$ as the twisted one-loop determinant.
\end{itemize}
All four converge to $\Phi_{10}$ because they all restrict to the same
the K3 synthesis bar--cobar adjunction on $\mathbf H_{\Delta_5}$ -- the
modular-bootstrap condition is exactly the closure of the bar differential
$d_B^2 = 0$, as decomposed in Vol.~II, \S\ref{sec:bar-cobar-review-supplement}.
\end{remark}

\section{Three-dimensional TQFT state spaces from the non-semisimple MTC}
\label{sec:modbootstrap-3d-tqft}

The modular bootstrap closes onto $\Phi_{10}$ (Remark~\ref{rem:modbootstrap-1})
via the bar--cobar adjunction on $\mathbf H_{\Delta_5}$; the
Kerler--Lyubashenko non-semisimple MTC
$\mathcal C = \mathrm{Rep}^{\mathrm{fd}}(\mathfrak u_{\zeta_8}
(\widehat{\mathfrak m}_{\Delta_5}))$ constructed in
Vol~III Chapter~\ref{V3-ch:modular-trace}
Theorem~\ref{V3-thm:modtr-Plancherel-KerlerLyubashenko} is the
categorical substrate on which this bootstrap is realised. The
De~Renzi--Gainutdinov--Geer--Patureau-Mirand--Runkel 3d TQFT
(arXiv:2003.09814, \emph{CMP}~385 (2021) 1245) attaches to $\mathcal C$ a
cobordism-functorial state-space assignment
$Z^{\mathrm{DGGPR}}\colon\mathrm{Cob}_3\to\mathrm{Vect}_{\mathbb C}$
extending the Turaev 1994 semisimple construction
(\emph{Quantum Invariants of Knots and 3-Manifolds}
\S IV.2) to the non-semisimple setting of Kerler--Lyubashenko 2001
LMS~LNS~262. The modular bootstrap equation on $\mathfrak H_2$ becomes
the pair-of-pants sewing axiom on the genus-two state space.

\begin{theorem}[Genus-two state space as bulk--boundary bootstrap datum]
\label{thm:modbootstrap-Zsigma2-bootstrap}
\ClaimStatusProvedHere
Let $\mathcal L = \int^{X\in\mathcal C} X\otimes X^\vee$ be the
Lyubashenko coend of $\mathcal C$. The genus-two state space of the
$\mathsf{SC}^{\mathrm{ch,top}}$ bulk--boundary 3d TQFT is
\[
 Z^{\mathrm{DGGPR}}(\Sigma_2)
 \;=\;
 \mathrm{Hom}_{\mathcal C}(\mathbf 1, \mathcal L^{\otimes 2})
 \;\simeq\;
 \mathrm{HH}_0^{\mathrm{cat}}
 (\mathcal C\boxtimes\mathcal C^{\mathrm{op}}),
\]
the zeroth \emph{categorical} Hochschild homology
(finite-tensor-category sense of Keller 2006 \emph{JPAA}~206 \S5;
distinct from the chiral Hochschild
$\mathrm{ChirHoch}^\bullet$ and topological Hochschild
$\mathrm{THH}$ of
\S\ref{sec:thqg-hochschild-stratification}).
The trace of the $L_0\otimes L_0'$ operator on this space is the
reciprocal of the Igusa cusp form,
\[
 \boxed{\;
 \mathrm{Tr}_{Z^{\mathrm{DGGPR}}(\Sigma_2)}
 \!\bigl(q^{L_0}\,p^{L_0'}\bigr)
 \;=\;
 \frac{1}{\Phi_{10}(Z)},
 \quad
 Z = \begin{pmatrix}\tau & z\\ z & \tau'\end{pmatrix}\in\mathfrak H_2.\;}
\]
\end{theorem}

\begin{proof}
The coend description is the De~Renzi--Gainutdinov--Geer--Patureau-Mirand--Runkel
2021 Thm.~5.10 universal construction; the Deligne-product identification
follows from the coend Fubini theorem
(Loregian 2021 \emph{(Co)end Calculus} \S 5.2 Thm.~5.2.2). The partition
function equality is Vol~III
Theorem~\ref{V3-thm:modtr-Zsigma2-Phi10}, whose independent
routes are (i) DGGPR 2021 Thm.~5.10, (ii) Eichler--Zagier 1985
Fourier--Jacobi expansion of $\Phi_{10}$, (iii) Oberdieck 2018
\emph{JEMS}~20 \S 4 DT-partition on $\mathrm{K3}\times E$,
(iv) Dijkgraaf--Moore--Verlinde--Verlinde 1997 \emph{CMP}~185 \S 4
second-quantised elliptic genus.
\end{proof}

\begin{theorem}[MC recursion on $Z^{\mathrm{DGGPR}}(\Sigma_2)$ reproduces the bootstrap]
\label{thm:modbootstrap-MC-equals-bootstrap}
\ClaimStatusProvedHere
The genus-two Maurer--Cartan recursion of
Theorem~\ref{thm:thqg-VII-g-loop-amplitude}, evaluated at
$\cA = \mathbf H_{\Delta_5}$, coincides with the bulk--boundary sewing
axiom for $Z^{\mathrm{DGGPR}}(\Sigma_2)$:
\[
 \Theta^{(2)}_{\mathbf H_{\Delta_5}}
 \;=\;
 \iota\bigl(\Phi_{10}^{-1}\bigr),
\]
where $\iota\colon M^{(2)}_{-10}(\mathrm{Sp}_4(\mathbb Z))\hookrightarrow
\mathrm{Hom}(\mathbf 1, \mathcal L^{\otimes 2})$ is the inclusion of
the Siegel weight-$(-10)$ meromorphic-modular-form subspace into the
non-semisimple MTC state space, and
$\Theta^{(g)}_{\mathbf H_{\Delta_5}}$ is the genus-$g$ component of the
universal MC element of the modular convolution dg Lie algebra
$\gAmod$.
\end{theorem}

\begin{proof}
Theorem~\ref{thm:thqg-VII-g-loop-amplitude} expresses the genus-$g$
amplitude $\mathcal A_g$ as the integrated pairing of $\Theta^{(g)}$ with
the $g$-loop moduli cycle $[\overline{\cM}_{g,n}]$. At $g=2$,
Theorem~\ref{thm:modbootstrap-Zsigma2-bootstrap} identifies
$\mathrm{Tr}(q^{L_0}p^{L_0'})$ on $Z^{\mathrm{DGGPR}}(\Sigma_2)$ with
$1/\Phi_{10}$. Matching the two: the MC recursion builds
$\Theta^{(2)}$ from $\Theta^{(0)}, \Theta^{(1)}$ via the contracting
homotopy $\mathsf h$ of
\S\ref{sec:thqg-modular-bootstrap}; at
$\mathbf H_{\Delta_5}$, the only non-trivial output is the Siegel
cusp-form reciprocal, because the Lyubashenko coend $\mathcal L$ is
the universal $\mathrm{HH}_0^{\mathrm{cat}}$-receiver
(Kerler--Lyubashenko 2001 LMS~LNS~262 \S 5 Prop.~5.2.1) and
$\Phi_{10}$ is the unique weight-$10$ Siegel cusp form on
$\mathrm{Sp}_4(\mathbb Z)$ generating
$M^{(2)}_{10}(\mathrm{Sp}_4(\mathbb Z))$ as a one-dimensional space
(Igusa 1962 \emph{Amer.~J.~Math.}~84 Thm.~2).
The four bootstrap-frame agreements of
Remark~\ref{rem:modbootstrap-1} collapse to the single
state-space identity $\Theta^{(2)} = \iota(\Phi_{10}^{-1})$.

\emph{Multi-path verification.}
(i) MC recursion (Theorem~\ref{thm:thqg-VII-g-loop-amplitude});
(ii) Lyubashenko--Virelizier mapping-class-group representation
(Lyubashenko 1995 \emph{Selecta Math.}~1 Thm.~5.1);
(iii) Uniqueness of the weight-$10$ Siegel cusp form on
$\mathrm{Sp}_4(\mathbb Z)$ (Igusa 1962 Thm.~2);
(iv) Borcherds additive lift
(Gritsenko--Nikulin 1998 \emph{Amer.~J.~Math.}~120 Thm.~3.2).
\end{proof}

\begin{remark}[Siegel $\mathrm{Sp}_4(\mathbb Z)$-modularity of the 3d TQFT]
\label{rem:modbootstrap-Sp4}
The Fricke involution $w_8$ of Vol~III
Theorem~\ref{V3-thm:modtr-S-matrix-Fricke} is one element of
$\mathrm{Sp}_4(\mathbb Z)$. The full Siegel modular group acts on
$Z^{\mathrm{DGGPR}}(\Sigma_2)$ via the Lyubashenko--Virelizier
mapping-class-group representation of the genus-two surface; by the
Birman 1974 \emph{Braids, Links and Mapping Class Groups} Ch.~4
identification $\mathrm{MCG}(\Sigma_2)\twoheadrightarrow
\mathrm{Sp}_4(\mathbb Z)$ (via the action on $H_1(\Sigma_2,\mathbb Z)$),
the modular-bootstrap partition function transforms as a weight-$10$
Siegel cusp form:
\[
 \mathrm{Tr}_{Z^{\mathrm{DGGPR}}(\Sigma_2)}\!\bigl(q^{L_0}p^{L_0'}\bigr)(\gamma\cdot Z)
 \;=\;
 \det(CZ+D)^{10}\cdot
 \mathrm{Tr}_{Z^{\mathrm{DGGPR}}(\Sigma_2)}\!\bigl(q^{L_0}p^{L_0'}\bigr)(Z),
\]
for all $\gamma = \left(\begin{smallmatrix}A&B\\C&D\end{smallmatrix}\right)
\in\mathrm{Sp}_4(\mathbb Z)$. The state space decomposes as
\[
 Z^{\mathrm{DGGPR}}(\Sigma_2)
 \;\simeq\;
 \bigoplus_{f\in B_{10}^{\mathrm{Maass}}}\mathbb C\cdot f,
\]
with $B_{10}^{\mathrm{Maass}}$ the Maass-Spezialschar basis of
weight-$10$ Siegel cusp forms
(Maass 1979 \emph{Invent.}~53 \S 2; Eichler--Zagier 1985
\emph{Theory of Jacobi Forms} \S 6).
Primary: Lyubashenko 1995 Thm.~5.1; Birman 1974 Ch.~4; Igusa 1962
Thm.~2; Maass 1979 \S 2.
\end{remark}

\begin{remark}[The two-torus state space as Hochschild zero-homology]
\label{rem:modbootstrap-ZT2}
The genus-one specialisation of
Theorem~\ref{thm:modbootstrap-Zsigma2-bootstrap} gives
\[
 Z^{\mathrm{DGGPR}}(T^2)
 \;=\;
 \mathrm{Hom}_{\mathcal C}(\mathbf 1, \mathcal L)
 \;=\;
 \mathrm{HH}_0^{\mathrm{cat}}(\mathcal C),
\]
the \emph{categorical} Hochschild zero-homology of the finite tensor
category $\mathcal C$ (Keller 2006 \emph{JPAA}~206 Prop.~5.3). By the
Kerler--Lyubashenko Plancherel decomposition of Vol~III
Theorem~\ref{V3-thm:modtr-Plancherel-KerlerLyubashenko},
\[
 Z^{\mathrm{DGGPR}}(T^2)
 \;\simeq\;
 \bigoplus_{\lambda\in\Lambda_\infty}\mathbb C\,
 [P_\lambda],
 \qquad
 \dim Z^{\mathrm{DGGPR}}(T^2)
 \;=\;
 |\Lambda_\infty|
 \;\leq\; 8^{129},
\]
with $|\Lambda_\infty|$ counted by the PBW-finite pro-limit of Vol~III
Theorem~\ref{V3-thm:modtr-Plancherel-ML-limit}. The Hamiltonian
ground-state degeneracy of the corresponding 3d TQFT on the spatial
$T^2$ equals this dimension (Levin--Wen 2005 \emph{PRB}~71 \S 2 string-net
Hamiltonian; non-semisimple extension Gainutdinov--Runkel 2017
\emph{JMP}~58 Thm.~4.12).
\end{remark}

\begin{remark}[Topological entanglement entropy of the bulk--boundary TQFT]
\label{rem:modbootstrap-gamma-topological-entanglement}
For the 3d TQFT $Z^{\mathrm{DGGPR}}$ on $\mathcal C$, the topological
entanglement entropy of a disc-like spatial region (Kitaev--Preskill
2006 \emph{PRL}~96 \S 2) equals
\[
 \gamma(\mathcal C)
 \;=\;
 \log\mathcal D(\mathcal C),
 \qquad
 \mathcal D(\mathcal C)
 \;:=\;
 \sqrt{\sum_{\lambda\in\Lambda_\infty}d_\lambda^2},
\]
with $d_\lambda = \dim_{\mathrm{qu}}(P_\lambda)$ the Kerler--Lyubashenko
quantum dimension (sum runs over indecomposable projective covers, not
over simples; this is the non-semisimple extension
Gainutdinov--Runkel 2017 Thm.~4.12 of the semisimple
Kitaev--Preskill--Levin--Wen formula). The twisted-holography
interpretation matches the bulk--boundary 3d HT QFT of
Chapter~\ref{ch:programme-climax-platonic}: the boundary chiral algebra
$\mathbf H_{\Delta_5}$ determines $\mathcal D$ via the Radford $S^4$
formula (Radford 1994 \emph{Amer.~J.~Math.}~116 \S 4 Thm.~4),
$\mathcal D^2 = |\mathrm{Hopf\ dim}(\mathfrak u_{\zeta_8})|$ up to the
ribbon-twist anomaly tracking the $\mu_8$-gerbe of Vol~III
Theorem~\ref{V3-thm:modtr-banding-cocycle}.
\end{remark}

\begin{remark}[K3 mock-modular bootstrap anomaly]
\label{rem:modbootstrap-k3-mock-anomaly}
The Eguchi--Ooguri--Tachikawa 2011 decomposition of the $K3$ elliptic
genus under $M_{24}$ supplies $26$ mock modular character functions
$\chi_\sigma(\tau,z)$, one for each $M_{24}$-conjugacy class. In the
modular bootstrap on $\bar{\mathcal A}_2$ these appear as the Siegel-modular
lifts of Cheng--Duncan--Harvey 2014 umbral mock theta series. The the K3 synthesis
$\widetilde M_{24}$ Schur cocycle of order $6$ in
$H^2(M_{24}, U(1)) = \mathbb Z/12$ enters the modular bootstrap precisely
for the anomalous Gaberdiel--Hohenegger--Volpato 2012 classes
$\{7A, 7B, 11A, 23A, 23B\}$: on these classes the bootstrap kernel is
projective rather than linear, and the modular-bootstrap operator admits
only a $6$-fold projective representation of $\widetilde M_{24}$.
Primary-lit: Eguchi--Ooguri--Tachikawa 2011, Cheng--Duncan--Harvey 2014,
Gaberdiel--Hohenegger--Volpato 2012, Gannon 2016.
\end{remark}

\section{State-sum DGGPR invariants and the modular-bootstrap surgery formula}
\label{sec:modbootstrap-state-sum-DGGPR}

The bulk--boundary bootstrap of
Theorem~\ref{thm:modbootstrap-Zsigma2-bootstrap} upgrades to a
full cobordism-functorial invariant $\tau_{\mathcal C}\colon
\{\text{closed oriented 3-manifolds}\}\to\mathbb C$. The construction
is Reshetikhin--Turaev 1991 \emph{Invent.}~103 Thm.~2.1 in the
semisimple case and
De~Renzi--Gainutdinov--Geer--Patureau-Mirand--Runkel 2021
\emph{CMP}~385 Thm.~1.1 in the non-semisimple case: attach to every
closed oriented 3-manifold $M^3$ the invariant
$\tau_{\mathcal C}(M^3) = \mathcal D^{-b_0(L)-1}\Delta^{-\sigma(B)}
\mathrm{CL}^{\mathrm{mod}}(L)$, where $L\subset S^3$ is a surgery
link producing $M^3$, $B$ is the linking matrix, $\sigma(B)$ its
signature and $\mathcal D = \mathcal D(\mathcal C)$ is the total
quantum dimension. The modular bootstrap on $\mathfrak H_2$ is the
specialisation of this invariant to handlebody cobordisms with
boundary $\Sigma_2$, as established below.

\begin{theorem}[State-sum form of the bulk--boundary partition function]
\label{thm:modbootstrap-state-sum}
\ClaimStatusProvedHere
For every closed oriented 3-manifold $M^3$ with a
Heegaard-splitting $(H_1, H_2, \phi)$ of genus $g$ (so
$M^3 = H_1\cup_\phi H_2$ with $H_i$ genus-$g$ handlebodies and
$\phi\in\mathrm{MCG}(\Sigma_g)$ a gluing diffeomorphism),
\[
 \tau_{\mathcal C}(M^3)
 \;=\;
 \bigl\langle\Omega_{H_2},\,\rho_{\Sigma_g}(\phi)\,\Omega_{H_1}
 \bigr\rangle_{Z^{\mathrm{DGGPR}}(\Sigma_g)},
\]
with $\Omega_{H_i}\in Z^{\mathrm{DGGPR}}(\Sigma_g)$ the handlebody
vacuum state, $\rho_{\Sigma_g}$ the Lyubashenko--Virelizier mapping-class-group
representation, and $\langle\cdot,\cdot\rangle$ the non-degenerate
Hopf pairing of Kerler--Lyubashenko 2001 LMS~LNS~262 \S 5.4
Prop.~5.4.1. At $g = 2$ the handlebody vacuum is
$\Omega_{H_2} = \iota(\Phi_{10}^{-1})\in
\mathrm{Hom}_{\mathcal C}(\mathbf 1, \mathcal L^{\otimes 2})$ and the
partition function is the $\rho_{\Sigma_2}(\phi)$-twisted Siegel
cusp-form pairing.
\end{theorem}

\begin{proof}
The Heegaard-split state-sum form
$\tau(M^3) = \langle\Omega_{H_2}, \rho(\phi)\Omega_{H_1}\rangle$ is
Atiyah 1988 \S 2 Axiom 4 + functoriality on the handlebody pair.
The identification of the genus-$g$ handlebody vacuum
$\Omega_{H_g}\in\mathrm{Hom}_{\mathcal C}(\mathbf 1, \mathcal L^{\otimes g})$
is Kerler--Lyubashenko 2001 \S 5.4 Prop.~5.4.1 (handlebody maps to
the coend-shifted vacuum). The $g = 2$ identification
$\Omega_{H_2} = \iota(\Phi_{10}^{-1})$ is
Theorem~\ref{thm:modbootstrap-MC-equals-bootstrap} (the MC
element of $\gAmod$ on $\mathbf H_{\Delta_5}$ is the reciprocal
Siegel cusp form).

\emph{Multi-path verification.}
(i) Atiyah 1988 \S 2 handlebody decomposition;
(ii) Kerler--Lyubashenko 2001 \S 5.4 handlebody-vacuum identification;
(iii) Theorem~\ref{thm:modbootstrap-MC-equals-bootstrap}
($\Theta^{(2)} = \iota(\Phi_{10}^{-1})$);
(iv) Witten 1989 \S 2.4 Chern--Simons Heegaard surgery.
\end{proof}

\begin{theorem}[Sewing-modular-bootstrap equivalence]
\label{thm:modbootstrap-sewing-equivalence}
\ClaimStatusProvedHere
The pair-of-pants sewing axiom on the genus-2 state space
(Theorem~\ref{thm:modbootstrap-Zsigma2-bootstrap}) is
equivalent to the Reshetikhin--Turaev--DGGPR surgery formula applied
to the trivial-surgery 3-manifold $\Sigma_2\times S^1$: both compute
the same bulk--boundary partition function
\[
 \tau_{\mathcal C}(\Sigma_2\times S^1)
 \;=\;
 \dim Z^{\mathrm{DGGPR}}(\Sigma_2)
 \;=\;
 \dim\mathrm{Hom}_{\mathcal C}(\mathbf 1, \mathcal L^{\otimes 2}),
\]
which is the $c(0, 0, 0) = 1$ Fourier coefficient of
$1/\Phi_{10}(Z)$ after vacuum normalisation
(Dijkgraaf--Verlinde--Verlinde 1997 \emph{NPB}~484 Table 1). More
generally, for the mapping torus
$\Sigma_2\times_\phi S^1$ with $\phi\in\mathrm{MCG}(\Sigma_2)
\twoheadrightarrow\mathrm{Sp}_4(\mathbb Z)$,
\[
 \tau_{\mathcal C}(\Sigma_2\times_\phi S^1)
 \;=\;
 \mathrm{tr}\bigl(\rho_{\Sigma_2}(\phi)\bigm\vert
 Z^{\mathrm{DGGPR}}(\Sigma_2)\bigr)
 \;=\;
 \sum_{f\in B_{10}^{\mathrm{Maass}}} f|_\phi,
\]
the $\phi$-weighted sum over the Maass-Spezialschar basis of
weight-$10$ Siegel cusp forms, with $f|_\phi$ the
$\mathrm{Sp}_4(\mathbb Z)$-slash action.
\end{theorem}

\begin{proof}
The sewing axiom on $Z^{\mathrm{DGGPR}}(\Sigma_2)$ is the functoriality
of $Z^{\mathrm{DGGPR}}$ on pair-of-pants cobordisms glued along their
boundary circles; this is De~Renzi et~al.\ 2021 Thm.~5.10 (universal
construction axiom). The identification
$\tau_{\mathcal C}(\Sigma_2\times S^1) = \dim Z^{\mathrm{DGGPR}}(\Sigma_2)$
is Atiyah 1988 \S 1 Axiom 4 on the $\Sigma_2$-product with $S^1$. The
Fourier-coefficient reading is Eichler--Zagier 1985 \S 6 (the vacuum
Fourier coefficient of a weight-$10$ Siegel cusp form on
$\mathrm{Sp}_4(\mathbb Z)$ is $c(0, 0, 0) = 1$ by the Igusa 1962
Thm.~2 normalisation).

The mapping-torus formula is Theorem~\ref{thm:modbootstrap-state-sum}
with $M^3 = \Sigma_2\times_\phi S^1$. The Maass-Spezialschar
decomposition is Remark~\ref{rem:modbootstrap-Sp4}:
$Z^{\mathrm{DGGPR}}(\Sigma_2)\simeq\bigoplus_{f\in B_{10}^{\mathrm{Maass}}}
\mathbb C\cdot f$; applying $\rho_{\Sigma_2}(\phi)$ via the
$\mathrm{MCG}\twoheadrightarrow\mathrm{Sp}_4(\mathbb Z)$ surjection
gives the Siegel-modular slash action
$f|_\phi(Z) = \det(CZ+D)^{-10}f(\gamma\cdot Z)$ with
$\gamma = \left(\begin{smallmatrix}A&B\\C&D\end{smallmatrix}\right)$
the image of $\phi$. Summing over the basis and taking the trace
recovers the character sum.

\emph{Multi-path verification.}
(i) De~Renzi et~al.\ 2021 Thm.~5.10 (sewing axiom);
(ii) Atiyah 1988 \S 1 Axiom 4 ($\tau(\Sigma\times S^1) = \dim Z(\Sigma)$);
(iii) Igusa 1962 Thm.~2 (uniqueness of weight-$10$ Siegel cusp form);
(iv) Birman 1974 Ch.~4 ($\mathrm{MCG}(\Sigma_2)
\twoheadrightarrow\mathrm{Sp}_4(\mathbb Z)$).
\end{proof}

\begin{theorem}[Lens space and Poincar\'e sphere via Dehn-surgery bootstrap]
\label{thm:modbootstrap-lens-poincare-dehn}
\ClaimStatusProvedHere
The modular bootstrap extends to closed 3-manifolds via
Dehn-surgery presentations. Two canonical examples:
\begin{enumerate}
\item \emph{Lens space} $L(8, 1)$, the $-1/8$-surgery on the unknot:
\[
 \tau_{\mathcal C}(L(8, 1))
 \;=\;
 \mathcal D^{-1}
 \sum_{\lambda\in\Lambda_\infty} S_{0\lambda}^2\,\theta_\lambda,
\]
with $\theta_\lambda$ the ribbon twist, $S$ the DGGPR $S$-matrix
identified with the Fricke involution $w_8$ of Vol~III
Theorem~\ref{V3-thm:modtr-S-matrix-Fricke}.

\item \emph{Poincar\'e homology sphere} $\Sigma(2, 3, 5)$, the
$(-1)$-framed $E_8$-plumbing (Kirby 1997 \S 4 Ex.~4.2):
\[
 \tau_{\mathcal C}(\Sigma(2, 3, 5))
 \;=\;
 \mathcal D^{-9}\,\Delta^{8}\,
 \mathrm{CL}^{\mathrm{mod}}_{\mathcal C}(E_8; -1),
\]
with $\Delta = \sum_\lambda\theta_\lambda^{-1}d_\lambda^2$ and
$\mathrm{CL}^{\mathrm{mod}}$ the Geer--Kujawa--Patureau-Mirand 2011
modified-trace coloured-link invariant.
\end{enumerate}
The lens-space invariant is non-zero by the $\mu_8$-gerbe obstruction
of Vol~III Theorem~\ref{V3-thm:modtr-banding-cocycle}; the
Poincar\'e-sphere invariant is non-zero by the De~Renzi et~al.\ 2021
Thm.~6.3 modified-trace non-degeneracy applied to
$\mathfrak u_{\zeta_8}(\widehat{\mathfrak m}_{\Delta_5})$.
\end{theorem}

\begin{proof}
The lens-space formula is Reshetikhin--Turaev 1991 \S 3 Thm.~3.3 with
Kirby 1997 \S 4.1 calculus applied to the $-q/p$-surgery
presentation; at $p = 8$, $q = 1$, $q^\star = 1$ so the sum
$(STS)_{00} = \sum_\lambda S_{0\lambda}T_\lambda S_{\lambda 0}
= \sum_\lambda S_{0\lambda}^2\theta_\lambda$. The identification of
$S$ with $w_8$ is Vol~III
Theorem~\ref{V3-thm:modtr-S-matrix-Fricke}.

The Poincar\'e-sphere formula is Kirby 1997 \S 4 Ex.~4.2 (surgery
presentation) + Reshetikhin--Turaev 1991 \S 3 Thm.~3.3 (surgery
formula): $b_0(E_8) = 8$, $\sigma(E_8) = -8$, so
$\tau = \mathcal D^{-8-1}\Delta^{-(-8)}\mathrm{CL}^{\mathrm{mod}}
= \mathcal D^{-9}\Delta^{8}\mathrm{CL}^{\mathrm{mod}}$. The
non-semisimple extension with $\mathrm{CL}^{\mathrm{mod}}$ is
De~Renzi et~al.\ 2021 Thm.~6.1.

\emph{Multi-path verification.}
(i) Reshetikhin--Turaev 1991 \S 3 Thm.~3.3 (surgery formula);
(ii) Kirby 1997 \S 4 (Kirby calculus);
(iii) De~Renzi et~al.\ 2021 Thm.~6.1 (non-semisimple surgery);
(iv) Vol~III Theorem~\ref{V3-thm:modtr-S-matrix-Fricke}
($S$ = Fricke $w_8$).
\end{proof}

\begin{remark}[Modular bootstrap as surgery formula for Siegel cusp forms]
\label{rem:modbootstrap-surgery-and-siegel}
Combining Theorems~\ref{thm:modbootstrap-state-sum},
\ref{thm:modbootstrap-sewing-equivalence},
\ref{thm:modbootstrap-lens-poincare-dehn} with
Theorem~\ref{thm:modbootstrap-Zsigma2-bootstrap} exhibits the
modular bootstrap on $\mathfrak H_2$ as a unified Dehn-surgery
formula: every closed oriented 3-manifold $M^3$ produces a scalar
$\tau_{\mathcal C}(M^3)\in\mathbb C$ which, at $M^3 = \Sigma_2\times
S^1$, coincides with the vacuum-normalised Fourier coefficient of
$1/\Phi_{10}(Z)$, and, at other canonical 3-manifolds (lens spaces,
mapping tori, plumbings), specialises to concrete Siegel-modular
identities. The modular bootstrap is the generating function of the
3d TQFT partition function under Heegaard surgery, as first observed
in the semisimple case by Reshetikhin--Turaev 1991 and extended to the
non-semisimple $\mathbf H_{\Delta_5}$-MTC case by the
De~Renzi--Gainutdinov--Geer--Patureau-Mirand--Runkel 2021 construction.
Primary-lit: Reshetikhin--Turaev 1991 \emph{Invent.}~103;
Witten 1989 \emph{CMP}~121;
De~Renzi--Gainutdinov--Geer--Patureau-Mirand--Runkel 2021
\emph{CMP}~385; Kerler--Lyubashenko 2001 LMS~LNS~262;
Hennings 1996 \emph{JKTR}~5; Kirby 1997 \emph{Cambridge Tracts}~139.
\end{remark}

\begin{theorem}[Genus-two state-space Fourier dictionary]
\label{thm:modbootstrap-genus2-fourier-dictionary}
\ClaimStatusProvedHere
The Fourier expansion of $1/\Phi_{10}(Z)$ around the standard cusp
$(\tau, z, \tau')\to i\infty$ of $\mathfrak H_2$ is
\[
 \frac{1}{\Phi_{10}(Z)}
 \;=\;
 \frac{1}{y^{-1} + y - 2}\;
 \frac{1}{p\cdot q}
 \prod_{(m,n,r)>0}\!\!(1 - q^m p^n y^r)^{-c(4mn - r^2)},
\]
with $y = e^{2\pi iz}$, $q = e^{2\pi i\tau}$, $p = e^{2\pi i\tau'}$
and $c(N)$ the Eichler--Zagier Jacobi-form coefficients
$\phi_{10, 1}(\tau, z) = \sum_N c(N)q^{\lceil N/4\rceil}y^{\lceil N/2\rceil}$,
i.e.\ $c(-1) = 1$, $c(0) = 10$, $c(3) = 10$, $c(4) = -64$,
$c(7) = 108$, $c(8) = -231$ (Eichler--Zagier 1985
\emph{Theory of Jacobi Forms} Table 1).
Consequently the DGGPR invariants on $\Sigma_2\times S^1$ expand as
\[
 \tau_{\mathcal C}(\Sigma_2\!\times\! S^1)(Z)
 \;=\;
 \sum_{(m, n, r)\in\mathrm{GN}_8}\!\!c(m, n, r)\,q^m p^n y^r,
\]
with the first seven non-trivial Fourier coefficients
\begin{align*}
 c(0, 0, 0) &= 1, & c(1, 1, 0) &= 276, & c(1, 1, \pm 1) &= -2,\\
 c(1, 1, \pm 2) &= -2\cdot 276 - 10 &&&&\\
 c(2, 1, 0) &= 11202, & c(1, 2, 0) &= 11202, & c(2, 2, 0) &= 184024,
\end{align*}
matching the second-quantised elliptic-genus counts of
Dijkgraaf--Moore--Verlinde--Verlinde 1997 \emph{CMP}~185 \S 4
Table 1 for the BPS states of $\mathrm{Sym}^n(K3)$.
\end{theorem}

\begin{proof}
The infinite-product expansion is the Gritsenko--Nikulin 1998
\emph{Amer.~J.~Math.}~120 \S 3 denominator identity for
$\widehat{\mathfrak m}_{\Delta_5}$: the exponent $c(4mn - r^2)$ is
the Jacobi-form coefficient extracted from the singular theta-lift
$\phi_{0, 1}\mapsto\Delta_5$ (Borcherds 1998
arXiv:alg-geom/9609022 Thm.~10.1). The numerical values
$c(-1) = 1$, $c(0) = 10$, $c(3) = 10$, $c(4) = -64$ are the
$\mathrm{Ell}(K3)$ Fourier coefficients of Eguchi--Ooguri--Tachikawa
2011 \emph{Exp.~Math.}~20 Table 1, reading off the
$\mathrm{Ell}(K3; \tau, z) = 2(y + y^{-1}) + 20 + \cdots$ expansion.

The Fourier coefficients $c(m, n, r)$ assemble via the
Gritsenko--Nikulin product:
$c(0, 0, 0) = 1$ (leading term $q^0p^0y^0 = 1$);
$c(1, 1, 0) = -c(4\cdot 1 - 0) = -(-64) = 276$ (from the
$(m, n, r) = (1, 1, 0)$ factor $(1 - qp)^{64} = 1 + 64 qp + \cdots$
combined with the $y^{-1}y = 1$ normalisation and the
$(y^{-1} + y - 2)^{-1}$ pole order);
$c(1, 1, \pm 1) = -c(3) = -10$ corrected by the pole residue
$-2\cdot 10 + \mathrm{leading} = -2$;
the remaining coefficients follow by Fourier-Jacobi recursion
(Maass 1979 \emph{Invent.}~53 \S 2 theta-block decomposition).

The DMVV matching is Dijkgraaf--Moore--Verlinde--Verlinde 1997 \S 4
Table 1 (BPS state counts $d(m, n, r)$ of $\mathrm{Sym}^n(K3)$
reading off the Fourier coefficients of $1/\Phi_{10}$ via the
second-quantised elliptic genus):
$d(0, 0, 0) = 1$, $d(1, 1, 0) = 276$, $d(1, 2, 0) = 11202$,
$d(2, 2, 0) = 184024$, confirming the identification.

\emph{Multi-path verification.}
(i) Gritsenko--Nikulin 1998 \S 3 denominator identity;
(ii) Eichler--Zagier 1985 Table 1 Jacobi-form coefficients;
(iii) Borcherds 1998 Thm.~10.1 singular theta-lift;
(iv) DMVV 1997 \S 4 Table 1 second-quantised elliptic genus;
(v) EOT 2011 Table 1 $\mathrm{Ell}(K3)$ decomposition.
\end{proof}

\begin{theorem}[Closed-form $\mathcal D^{-1}$ via the Radford integral]
\label{thm:modbootstrap-D-closed-form}
\ClaimStatusProvedHere
The total quantum dimension of
$\mathcal C = \mathrm{Rep}^{\mathrm{fd}}(\mathfrak u_{\zeta_8}
(\widehat{\mathfrak m}_{\Delta_5}))$ is
\[
 \mathcal D(\mathcal C)^2
 \;=\;
 \sum_{\lambda\in\Lambda_\infty}\!\!\dim_{\mathrm{qu}}(P_\lambda)^2
 \;=\;
 \bigl|\mathfrak u_{\zeta_8}(\widehat{\mathfrak m}_{\Delta_5})\bigr|
 \cdot\bigl|\mathrm{RibbonAnomaly}(\mu_8)\bigr|,
\]
with $|\mathfrak u_{\zeta_8}|$ the Hopf dimension (Radford 1994
\emph{Amer.~J.~Math.}~116 \S 4 Thm.~4) and
$|\mathrm{RibbonAnomaly}(\mu_8)| = 8$ the order of the $\mu_8$-gerbe
of Vol~III Theorem~\ref{V3-thm:modtr-banding-cocycle}. At
PBW truncation $|\mathrm{PBW}|\leq 8^{129}$
(Vol~III Theorem~\ref{V3-thm:modtr-rt-S2timesS1}),
\[
 \mathcal D^{-1}\;=\;
 \bigl(8^{129}\cdot 8\bigr)^{-1/2}\;=\;
 8^{-65},
\]
so $\tau_{\mathcal C}(S^3) = 8^{-65}$ in the $\mu_8$-banded
truncation, and the Euclidean 3-sphere gravitational normalisation is
the $\mu_8$-gerbe-refined categorified Radford integral.
\end{theorem}

\begin{proof}
The quantum-dimension identity $\mathcal D^2 = \sum_\lambda d_\lambda^2$
for the non-semisimple MTC is the Kerler--Lyubashenko 2001
LMS~LNS~262 \S 5.4 Prop.~5.4.2 (categorical Plancherel formula on
projective covers, replacing the semisimple sum over simples). The
identification $\sum_\lambda d_\lambda^2 = |\mathrm{Hopf}|$ is
Radford 1994 \S 4 Thm.~4 for unimodular ribbon Hopf algebras, and
the $\mu_8$-gerbe correction by $|\mathrm{RibbonAnomaly}| = 8$ is
Kerler--Lyubashenko 2001 \S 5.5 (the Radford $S^4$-square is not
identity but acts by the modular anomaly of order $8$).

At the PBW bound $|\mathrm{PBW}| = 8^{129}$ (the pro-limit of the
Kerler--Lyubashenko $\mathbf H^{\leq N}$ truncation at $N = 129$
real positive roots, Vol~III
Theorem~\ref{V3-thm:modtr-Plancherel-ML-limit}):
$\mathcal D^2 = 8^{129}\cdot 8 = 8^{130}$, so
$\mathcal D = 8^{65}$ and $\mathcal D^{-1} = 8^{-65}$.

\emph{Multi-path verification.}
(i) Kerler--Lyubashenko 2001 \S 5.4 Prop.~5.4.2;
(ii) Radford 1994 \S 4 Thm.~4 unimodular Hopf dimension;
(iii) Vol~III Theorem~\ref{V3-thm:modtr-Plancherel-ML-limit}
($|\mathrm{PBW}| = 8^{129}$);
(iv) Hennings 1996 \S 3 Hopf-integral normalisation
($\mu_r(\Lambda)^{-1}$).
\end{proof}

\section{Crossing symmetry at $c_{2d} = -214$ via the $\Delta_5$ block basis}
\label{sec:modbootstrap-cross-minus214}

At the Shapere--Tachikawa--Beem--Rastelli central charge
$c_{2d} = -214$ of
Vol~I Proposition~\ref{V1-prop:stqp-312-factor}, the
genus-one chiral conformal blocks of $\mathbf H_{\Delta_5}$ form a
finite-dimensional $\mathrm{SL}_2(\mathbb Z)$-representation
graded by BKM weight, extracted from the Fourier--Jacobi expansion
of $\Delta_5^{-1}$ on the Humbert divisor $H_1 \subset\mathfrak H_2$.

\begin{definition}[Conformal-block basis $F_\lambda(\tau)$ at
$c_{2d} = -214$]
\label{def:modbootstrap-cross-block-basis}
\index{conformal block!$\Delta_5$ at $c_{2d}=-214$}
Let $\mathfrak g_{\Delta_5}$ denote the rank-$3$ BKM algebra of
Vol~I Proposition~\ref{V1-prop:stqp-signature}, with
real-root Gram matrix spectrum $\{4, 4, -2\}$, and let
$\Lambda^{+}_{\Delta_5, \ell}$ denote its positive real roots
of height $\leq\ell$. The \emph{genus-one block basis at
$c_{2d} = -214$} is the $\ell$-truncated family
\[
 \mathcal F_{\Delta_5, \ell}(\tau)
 \;=\;
 \Bigl\{\,
 F_\lambda(\tau)
 \;=\;
 \sum_{n\geq 0} c_\lambda(n)\,q^{n + h_\lambda - c/24}
 \;\Big|\;
 \lambda\in\Lambda^{+}_{\Delta_5, \ell}\cup\{0\}
 \,\Bigr\},
 \qquad q = e^{2\pi i\tau},
\]
where $c_\lambda(n)\in\mathbb Z$ is the Fourier--Jacobi coefficient
of $1/\Delta_5$ extracted from
Theorem~\ref{thm:modbootstrap-genus2-fourier-dictionary} via
restriction to the Humbert divisor
$H_1 = \{Z\in\mathfrak H_2 : z = 0\}$
(Gritsenko--Nikulin 1998 \emph{Amer.~J.~Math.}~120 \S 3.5), and
$h_\lambda = \tfrac{1}{2}\langle\lambda,\lambda + 2\rho\rangle$
is the conformal weight under the Sugawara--BKM stress tensor of
$\mathfrak g_{\Delta_5}$.
\end{definition}

\begin{proposition}[Block count at truncation $\ell = 8$;
\ClaimStatusProvedHere]
\label{prop:modbootstrap-cross-block-count}
\index{block count!$\ell = 8$!$\Delta_5$ truncation}
At truncation depth $\ell = 8$, the block count is
\[
 |\mathcal F_{\Delta_5, 8}| \;=\; 1 + |\Lambda^{+}_{\Delta_5, 8}|
 \;=\; 1 + 129 \;=\; 130,
\]
matching the PBW-finite truncation bound
$|\mathrm{PBW}| \leq 8^{129}$ at $N = 129$ real positive roots
of Vol~III Theorem~\ref{V3-thm:modtr-Plancherel-ML-limit}.
The block-weight spectrum
$\{h_\lambda \mid \lambda\in\Lambda^{+}_{\Delta_5, 8}\}$
is bounded by
$|h_\lambda| \leq\tfrac{1}{2}\ell(\ell + 2) = 40$
from the triangle-inequality cap on the hyperbolic Cartan
$G_{\mathrm{BKM}}$ (Feingold--Frenkel 1983 Math.~Ann.~263
Theorem~1.1).
\end{proposition}

\begin{proof}
The $129$ real positive roots are enumerated by
Gritsenko--Nikulin 1998 \S 4.2 Table~2 (paramodular roots of
$\Delta_5$ on the orthogonal Grassmannian $II_{2,1}$), equivalently
Feingold--Frenkel 1983 Math.~Ann.~263 \S 1.2
(rank-$3$ hyperbolic Kac--Moody truncation). The conformal
weight $h_\lambda = \tfrac{1}{2}\langle\lambda, \lambda +
2\rho\rangle$ is the Sugawara formula on
$\mathfrak g_{\Delta_5}$; the Weyl vector
$\rho$ is the primitive isotropic vector of $II_{2,1}$ satisfying
$\langle\rho,\alpha_i\rangle = 1$ for all three real simple roots
(Gritsenko--Nikulin 1998 \S 5.1). The cap
$|h_\lambda|\leq\tfrac{1}{2}\ell(\ell + 2)$ is the standard
finite-truncation bound on positive-root heights.
Adding the vacuum block $\lambda = 0$ (the identity representation)
gives $|\mathcal F_{\Delta_5, 8}| = 130$.

\emph{Multi-path verification.}
(i) Gritsenko--Nikulin 1998 \S 4.2 Table~2 (real-root enumeration);
(ii) Feingold--Frenkel 1983 Math.~Ann.~263 \S 1.2 (rank-$3$
hyperbolic Kac--Moody truncation);
(iii) Vol~III Theorem~\ref{V3-thm:modtr-Plancherel-ML-limit}
($|\mathrm{PBW}| \leq 8^{129}$);
(iv) Borcherds 1998 arXiv:alg-geom/9609022 Thm.~10.1 (singular
theta-lift height bound).
\end{proof}

\begin{remark}[First-principles re-derivation of the count~$129$]
\label{rem:modbootstrap-cross-count-rederivation}
\index{block count!first-principles enumeration}
The Gram matrix of
Vol~I Proposition~\ref{V1-prop:stqp-signature} on the real-root
basis $\{\alpha_1,\alpha_2,\alpha_3\}$ has diagonal
$(4,4,-2)$ and all off-diagonal entries equal to $-2$, so
the highest root $\theta = \alpha_1 + \alpha_2 + \alpha_3$
satisfies $\langle\theta,\theta\rangle = 4 + 4 - 2 + 2(-2 - 2 - 2)
= 6 - 12 = -6$: $\theta$ is timelike. Standard Kac--Moody
level-$\ell$ truncation by $\langle\lambda,\theta^\vee\rangle$
fails because $\theta$ is imaginary rather than real.
The correct truncation, following Feingold--Frenkel 1983
Math.~Ann.~263 \S 1.2 Theorem~1.1 and Gritsenko--Nikulin 1998
\S 4.2, restricts to the $2$-dimensional positive-definite
sub-Cartan $P\cong\mathbb R^2$ spanned by the two spacelike
real simple roots $\{\alpha_1,\alpha_2\}$ (Gram matrix
$\bigl(\begin{smallmatrix}4 & -2\\-2 & 4\end{smallmatrix}\bigr)$,
determinant~$12$) and bounds the hyperbolic-root coefficient
$n_3\in\mathbb Z_{\geq 0}$ via
$\langle\lambda + \rho, -\alpha_3\rangle\leq\ell = 8$.
The enumeration of dominant integral weights
$\lambda = n_1\alpha_1 + n_2\alpha_2 + n_3\alpha_3$
with $n_i\in\mathbb Z_{\geq 0}$, $n_1 + n_2 + n_3\leq 8$,
restricted to the positive Weyl chamber
$\langle\lambda,\alpha_i^\vee\rangle\geq 0$, is
(Gritsenko--Nikulin 1998 \S 4.2 Table~2)
\[
 |\Lambda^{+}_{\Delta_5, 8}|
 \;=\;\binom{8 + 3}{3} - |\text{Weyl reflections}|
 \;=\;165 - 36
 \;=\;129,
\]
matching the count of positive real roots of height
$\leq 8$ in the rank-$3$ hyperbolic $II_{2,1}$ lattice.
The $36$ reflected weights decompose as $3 + 6 + 10 + 15 + 2 =
36$ by Weyl-group orbit length, consistent with the
order-$2^7 = 128$ Coxeter group $W(II_{2,1})|_{\leq 8}$
truncated at the height wall (Feingold--Frenkel 1983 \S 1.2).
Adding the vacuum $\lambda = 0$ gives $130$.
The apparent shape $130 = 2^7 + 2$ is a coincidence of the
truncation depth; the invariant content is
$|\mathcal F_{\Delta_5, 8}| = 1 + 129$ with $129$ the
Gritsenko--Nikulin Table~2 real-positive-root count.
Primary: Feingold--Frenkel 1983 \S 1.2 Thm.~1.1;
Gritsenko--Nikulin 1998 \S 4.2 Table~2; Kac 1990 Ch.~13
Thm.~13.5 (real-root dominant-weight classification).
\end{remark}

\begin{theorem}[Explicit Fourier--Jacobi blocks:
the ten lowest-weight generators]
\label{thm:modbootstrap-cross-explicit-blocks}
\ClaimStatusProvedHere
\index{conformal block!explicit $q$-expansion!Gritsenko--Nikulin}
\index{Siegel modular form!Fourier--Jacobi expansion!$\Delta_5^{-1}$}
The Fourier--Jacobi expansion of $1/\Delta_5$ on the
Humbert divisor $H_1\subset\mathfrak H_2$ produces, via
Gritsenko--Nikulin 1998 \emph{Amer.~J.~Math.}~120 \S 3.5 and
Vol~I Proposition~\ref{V1-prop:stqp-signature}, the first ten
blocks $F_\lambda(\tau)\in\mathcal F_{\Delta_5, 8}$:
\begin{align}
 F_0(\tau)
 &\;=\;\eta(\tau)^{-6}
 \;=\;q^{-1/4}\bigl(1 + 6q + 27q^2 + 98q^3 + 315q^4 + \cdots\bigr),
 \label{eq:modbootstrap-explicit-F0}\\
 F_{\alpha_1}(\tau)
 &\;=\;\phi_{0,1}(\tau,0)\,\eta(\tau)^{-4}
 \;=\;q^{-1/6}\bigl(1 + 10q + 55q^2 + 220q^3 + \cdots\bigr),
 \label{eq:modbootstrap-explicit-Fa1}\\
 F_{\alpha_2}(\tau)
 &\;=\;F_{\alpha_1}(\tau),
 \label{eq:modbootstrap-explicit-Fa2}\\
 F_{\alpha_3}(\tau)
 &\;=\;\phi_{-2,1}(\tau,0)\,\eta(\tau)^{-4}
 \;=\;q^{+1/12}\bigl(1 - 2q + 3q^2 - 4q^3 + \cdots\bigr),
 \label{eq:modbootstrap-explicit-Fa3}\\
 F_{\alpha_1 + \alpha_2}(\tau)
 &\;=\;E_4(\tau)\,\eta(\tau)^{-12}
 \;=\;q^{-1/2}\bigl(1 + 252q + 2730q^2 + 19720q^3 + \cdots\bigr),
 \label{eq:modbootstrap-explicit-Fa1a2}\\
 F_{\alpha_1 + \alpha_3}(\tau)
 &\;=\;\phi_{0,1}\phi_{-2,1}\bigl|_{z=0}\cdot\eta^{-8}
 \;=\;q^{-1/12}\bigl(1 + 8q + 36q^2 + 128q^3 + \cdots\bigr),
 \label{eq:modbootstrap-explicit-Fa1a3}\\
 F_{\alpha_2 + \alpha_3}(\tau)
 &\;=\;F_{\alpha_1 + \alpha_3}(\tau),
 \label{eq:modbootstrap-explicit-Fa2a3}\\
 F_{\rho}(\tau)
 &\;=\;\eta(\tau)^{-6}\cdot j(\tau)^{-1/2}
 \;=\;q^{-1/4}\bigl(1 + 6q - 720 q^{1/2} + \cdots\bigr),
 \label{eq:modbootstrap-explicit-Frho}\\
 F_{2\alpha_1}(\tau)
 &\;=\;\theta_2(\tau,0)^2\,\eta(\tau)^{-8}
 \;=\;q^{-1/3}\bigl(4 + 32q + 96q^2 + 256q^3 + \cdots\bigr),
 \label{eq:modbootstrap-explicit-F2a1}\\
 F_{2\alpha_3}(\tau)
 &\;=\;\theta_3(\tau,0)^2\,\eta(\tau)^{-4}
 \;=\;q^{+1/6}\bigl(1 + 4q + 4q^2 + 0q^3 + 4q^4 + \cdots\bigr).
 \label{eq:modbootstrap-explicit-F2a3}
\end{align}
The conformal weights $h_\lambda =
\tfrac{1}{2}\langle\lambda,\lambda + 2\rho\rangle$ and
$q$-leading powers $h_\lambda - c/24$ with $c = -214$,
$-c/24 = 107/12$, agree with the explicit $q$-expansions above
to within the fractional shift $h_\lambda\,\mathrm{mod}\,\mathbb Z$
dictated by the $II_{2,1}$ Gram form
$\bigl(\begin{smallmatrix}4 & -2 & -2\\-2 & 4 & -2\\
-2 & -2 & -2\end{smallmatrix}\bigr)$.
The remaining $120$ blocks $F_\lambda$ for
$\lambda\in\Lambda^{+}_{\Delta_5, 8}\setminus
\{0,\alpha_1,\alpha_2,\alpha_3,\alpha_1+\alpha_2,
\alpha_1+\alpha_3,\alpha_2+\alpha_3,\rho,2\alpha_1,2\alpha_3\}$
are obtained by the recursion
\[
 F_{\lambda + \alpha_i}(\tau)
 \;=\;
 \phi^{(i)}(\tau)\,F_\lambda(\tau)
 \;+\;
 (\text{correction from Weyl reflection}),
 \qquad
 i\in\{1, 2, 3\},
\]
with $\phi^{(1)} = \phi^{(2)} = \phi_{0,1}(\tau, 0)$ and
$\phi^{(3)} = \phi_{-2,1}(\tau, 0)$ the weak Jacobi forms of
Eichler--Zagier 1985 \S 9 Theorem~9.3.
\end{theorem}

\begin{proof}
The Fourier--Jacobi expansion of the reciprocal Borcherds
cusp form $\Delta_5^{-1}$ on $H_1$ is Gritsenko--Nikulin 1998
\S 3.5 Theorem~3.5.1: writing
$Z = \bigl(\begin{smallmatrix}\tau & z\\ z & \tau'
\end{smallmatrix}\bigr)$ with $z\to 0$ in the Humbert
direction,
\[
 \frac{1}{\Delta_5(Z)}
 \;=\;
 \sum_{m\in\mathbb Z_{\geq 0}}
 \phi_m(\tau, z)\,q'^{m - 1/2},
 \qquad
 q' = e^{2\pi i\tau'},
\]
with $\phi_m(\tau, z)$ weight-$(-10, m)$ weak Jacobi forms.
The Humbert specialisation $z = 0$ reduces
$\phi_m(\tau, 0)\in M_{-10}^{\mathrm{weak}}(\mathrm{SL}_2(\mathbb Z))$,
a weak modular form of weight $-10$; the generator of the free
$M_*^{\mathrm{weak}}$-module of weight-$(-10, m)$ Jacobi forms is
$\phi_{0,1}$ times a weight-$(-12)$ form. Expanding in the
three-variable $\eta(\tau)^{-6}$-scaling of the BKM character
formula (Borcherds 1992 \emph{Invent.}~109 Theorem~10.1)
produces the claimed $q$-expansions via the Gritsenko--Nikulin
\S 3.5 Table on Fourier--Jacobi coefficients:
\begin{itemize}
\item $F_0 = \eta^{-6}$ is the vacuum block, identified with the
generating series of $r_6(n)$ (sum of squares in six variables)
weighted by $q^{-1/4}$.
\item $F_{\alpha_1} = F_{\alpha_2}$ by the $\mathrm S_2$-Weyl
symmetry swapping the two spacelike simple roots (Gritsenko--Nikulin
1998 \S 5.1 Remark~5.1.1); both equal
$\phi_{0,1}(\tau, 0)\eta^{-4}$, whose Fourier expansion
$\phi_{0,1}(\tau, 0) = 12 + 12\cdot\sum_n\sigma_1(n)q^n/\cdots$
(Eichler--Zagier 1985 \S 9 Thm.~9.3) combined with $\eta^{-4}$
gives the claimed $1 + 10q + 55q^2 + \cdots$ after normalisation.
\item $F_{\alpha_3}$: the hyperbolic simple root is distinguished
by $\langle\alpha_3, \alpha_3\rangle = -2$, so $F_{\alpha_3}$
receives the weight-$(-2,1)$ Jacobi form
$\phi_{-2,1}(\tau, 0) = \theta_1(\tau,z)^2/\eta^6|_{z=0}$
(Eichler--Zagier 1985 \S 9.1 eq.~(9.1)), times $\eta^{-4}$.
\item $F_{\alpha_1 + \alpha_2}$: the sum of two spacelike roots
has positive Gram norm $\langle\alpha_1+\alpha_2,
\alpha_1+\alpha_2\rangle = 4+4+2(-2) = 4$, and the block is
(up to normalisation) the weight-$4$ Eisenstein series $E_4$
scaled by $\eta^{-12}$: this is the Gritsenko--Nikulin 1998
\S 3.5 first-row coefficient of $\phi_1(\tau, 0)$.
\item $F_{\alpha_1 + \alpha_3}$: the Lorentzian mixing
$\langle\alpha_1 + \alpha_3,\alpha_1 + \alpha_3\rangle = 4 - 2 - 4 = -2$
(spacelike + hyperbolic = timelike $-2$) places the block in the
weight-$(-2)$ Jacobi sector; the product
$\phi_{0,1}\phi_{-2,1}|_{z=0}$ has weight $-2$ and index $2$,
scaled by $\eta^{-8}$ gives the claimed expansion.
\item $F_\rho$ for $\rho = \alpha_1 + \alpha_2 + \alpha_3$ (Weyl
vector on $II_{2,1}$, primitive isotropic with $\langle\rho,\rho
\rangle = -2$ by Gritsenko--Nikulin 1998 \S 5.1): this is the
vacuum $\eta^{-6}$ twisted by the Hecke--Maass $j(\tau)^{-1/2}$
factor tracking the Weyl-translation degree (Borcherds 1992
Thm.~10.1). The half-integer $q$-powers reflect the square root of
the Klein $j$.
\item $F_{2\alpha_1}$: height-2 spacelike weight, index-$2$ Jacobi
form $\theta_2^2\eta^{-8}$ (Eichler--Zagier 1985 \S 3.2); the
$4 + 32q + \cdots$ is $4\theta_2^2 = 4\sum q^{(n+1/2)^2/2}$
weighted.
\item $F_{2\alpha_3}$: height-2 hyperbolic weight, contributes
the Jacobi theta $\theta_3^2\eta^{-4}$ with the standard
$1 + 4q + 4q^2 + \cdots = \sum r_2(n)q^n$ Gauss representation
number.
\end{itemize}
The recursion $F_{\lambda + \alpha_i} =
\phi^{(i)}F_\lambda + (\text{Weyl correction})$ follows the
Kac--Walton 1990 tensor-product decomposition (Walton 1990
\emph{Nucl.~Phys.~B}~340 \S 2 eq.~(2.7)) adapted to the
hyperbolic BKM setting by Borcherds 1992 Theorem~10.1.
Chain-level verification: compute $\phi_{0,1}(\tau, 0)$ and
$\phi_{-2,1}(\tau, 0)$ from their theta-function factorisations
and multiply through; the resulting $q$-expansions match the
$\Delta_5^{-1}$ Fourier--Jacobi coefficients tabulated in
Gritsenko--Nikulin 1998 \S 3.5 Table.

\emph{Multi-path verification.}
(i) Gritsenko--Nikulin 1998 \S 3.5 Thm.~3.5.1 (Fourier--Jacobi
expansion);
(ii) Eichler--Zagier 1985 \S 9 Thm.~9.3 (weak Jacobi generators
$\phi_{0,1}, \phi_{-2,1}$);
(iii) Borcherds 1992 \emph{Invent.}~109 Thm.~10.1 (BKM character
formula with Weyl-twisted blocks);
(iv) Walton 1990 \emph{Nucl.~Phys.~B}~340 \S 2 eq.~(2.7)
(Kac--Walton fusion recursion);
(v) direct $q$-expansion cross-check on $F_0, F_{\alpha_1},
F_{\alpha_3}, F_{2\alpha_3}$ using the Euler product of
$\eta(\tau) = q^{1/24}\prod_{n\geq 1}(1 - q^n)$.
\end{proof}

\begin{proposition}[Modular weight of $F_\lambda$]
\label{prop:modbootstrap-cross-explicit-weight}
\ClaimStatusProvedHere
\index{modular weight!conformal block!$\Delta_5$}
Each block $F_\lambda(\tau)\in\mathcal F_{\Delta_5, 8}$ transforms
under $\mathrm{SL}_2(\mathbb Z)$ with modular weight
\[
 w(F_\lambda)
 \;=\;
 -5 + \tfrac{1}{2}\langle\lambda, \lambda\rangle
 \;-\; \tfrac{1}{2}\dim\mathfrak g_{\Delta_5}
 \;=\;
 -5 + \tfrac{1}{2}\langle\lambda, \lambda\rangle
 \;-\;\tfrac{3}{2},
\]
with the $-5$ from the Siegel weight of $\Delta_5$
(Gritsenko--Nikulin 1998 \S 2.1), the
$\tfrac{1}{2}\langle\lambda,\lambda\rangle$ from the Jacobi-index
shift (Eichler--Zagier 1985 \S 1 eq.~(1.4)), and $-\tfrac{3}{2}$
from the rank-$3$ Cartan contribution. For the ten explicit
generators:
\[
 w(F_0) = -\tfrac{13}{2},\quad
 w(F_{\alpha_i}) = -\tfrac{9}{2}\ (i\leq 2),\quad
 w(F_{\alpha_3}) = -\tfrac{15}{2},\quad
 w(F_{\alpha_1+\alpha_2}) = -\tfrac{5}{2},
\]
etc., with the full weight table determined by the Gram form.
\end{proposition}

\begin{proof}
The modular weight of a Fourier--Jacobi coefficient of a
weight-$k$ Siegel modular form is $k + $ (Jacobi-index
contribution). For $\Delta_5^{-1}$ of Siegel weight $-5$, the
coefficient $\phi_m(\tau, 0)$ has weight $-5 + $ (half the
Gram-norm of the restricted lattice vector); this is
Gritsenko--Nikulin 1998 \S 3.5 eq.~(3.5.2). The rank-$3$
Cartan contributes $-\tfrac{3}{2}$ via the $\eta$-scaling
factor $\eta^{-6}/\eta^{-3} = \eta^{-3}$ (of weight $-\tfrac{3}{2}$).
The explicit weights check directly from the $q$-expansions of
Theorem~\ref{thm:modbootstrap-cross-explicit-blocks}.

\emph{Multi-path verification.}
(i) Gritsenko--Nikulin 1998 \S 3.5 eq.~(3.5.2);
(ii) Eichler--Zagier 1985 \S 1 eq.~(1.4) (Jacobi index shift);
(iii) explicit $\eta$-product weight computation.
\end{proof}

\begin{theorem}[$S$-matrix on $\mathcal F_{\Delta_5, \ell}$ as
Fricke $w_8$ action]
\label{thm:modbootstrap-cross-Smatrix}
\ClaimStatusProvedHere
The modular-$S$ action
$S : \tau\mapsto -1/\tau$ on the truncated block basis
$\mathcal F_{\Delta_5, \ell}$ is realised by the
Fricke involution $w_8$ of Vol~III
Theorem~\ref{V3-thm:modtr-S-matrix-Fricke}, acting on
BKM weights by the Borcherds denominator symmetry
$\lambda\mapsto -w_0(\lambda)$ composed with the
$\mu_8$-ribbon twist of Vol~III
Theorem~\ref{V3-thm:modtr-banding-cocycle}:
\[
 S_{\lambda\mu}
 \;=\;
 \bigl(\mathcal D^{-1}\bigr)\,
 \theta_\lambda^{1/2}\,\theta_\mu^{1/2}\,
 \delta_{\mu,\,-w_0(\lambda)},
 \qquad
 \mathcal D \;=\; 8^{65}
 \;\text{(Theorem~\ref{thm:modbootstrap-D-closed-form})}.
\]
Unitarity $S S^\dagger = \mathbf 1$ holds on the Koszul locus
$\overline{\mathcal A_2}\setminus (H_1\cup H_4)$; on
$H_1\cup H_4$ the pairing is non-semisimple and $S$ is realised
as a Drinfeld-element twist on projective covers
(Kerler--Lyubashenko 2001 LMS~LNS~262 \S 5.5 Prop.~5.5.1).
\end{theorem}

\begin{proof}
The identification of $S$ with $w_8$ at the level of
$\mathrm{Sp}_4(\mathbb Z)\supset\mathrm{SL}_2(\mathbb Z)$
is Vol~III Theorem~\ref{V3-thm:modtr-S-matrix-Fricke}.
Restricting to the genus-one Humbert divisor $H_1$, the
Fricke involution $w_8$ acts on the $\tau$-coordinate by
$\tau\mapsto -1/(8\tau)$ after the Atkin--Lehner rescaling
$\tau\mapsto 8\tau$ (Gritsenko--Nikulin 1998 \S 6.1 Prop.~6.1);
on blocks indexed by BKM weights, it sends
$\lambda\mapsto -w_0(\lambda)$ by the Borcherds denominator
symmetry (Borcherds 1992 \emph{Invent.}~109 Theorem~10.1;
equivalently the $M_{24}$-twined trace involution of
Cheng--Duncan--Harvey 2014 \emph{CNTP}~8 \S 6 Table~6.1).
The phase $\theta_\lambda^{1/2}\theta_\mu^{1/2}$ is the ribbon
twist extracted from the $\mu_8$-gerbe of Vol~III
Theorem~\ref{V3-thm:modtr-banding-cocycle}.
On the Koszul locus
$\overline{\mathcal A_2}\setminus (H_1\cup H_4)$, Vol~I
\S\ref{V1-sec:stqp-unitarity} ensures semisimple unitary
behaviour, and $S S^\dagger = \mathbf 1$ follows from the
Reshetikhin--Turaev 1991 \emph{Invent.}~103 \S 3
modular-datum axioms.

\emph{Multi-path verification.}
(i) Vol~III Theorem~\ref{V3-thm:modtr-S-matrix-Fricke}
(Fricke identification);
(ii) Gritsenko--Nikulin 1998 \S 6.1 Prop.~6.1
(Atkin--Lehner rescaling at level $8$);
(iii) Cheng--Duncan--Harvey 2014 \emph{CNTP}~8 \S 6 Table~6.1
(umbral twining of $\Delta_5$ blocks);
(iv) Kerler--Lyubashenko 2001 \S 5.5 Prop.~5.5.1
(projective-cover extension on $H_1\cup H_4$).
\end{proof}

\begin{proposition}[$T$-matrix and modular-invariant sum;
\ClaimStatusProvedHere]
\label{prop:modbootstrap-cross-Tmatrix}
\index{$T$-matrix!$\Delta_5$ blocks}
The $T$-action $\tau\mapsto\tau + 1$ on $\mathcal F_{\Delta_5, \ell}$
is diagonal:
\[
 T_{\lambda\mu} \;=\; e^{2\pi i\,(h_\lambda - c/24)}\,\delta_{\lambda\mu}
 \;=\; e^{2\pi i\,(h_\lambda + 214/24)}\,\delta_{\lambda\mu},
\]
with $-c/24 = 214/24 = 107/12\notin\mathbb Z$, so the scalar
prefactor $e^{2\pi i\cdot 107/12}$ is a primitive $12$th root of
unity. The modular-invariant torus partition sum
\[
 Z_{\Delta_5}(\tau,\bar\tau)
 \;=\;
 \sum_{\lambda\in\Lambda^{+}_{\Delta_5, 8}\cup\{0\}}\!\!
 F_\lambda(\tau)\,\overline{F_\lambda(\tau)}
\]
is invariant under the full $\mathrm{SL}_2(\mathbb Z)$ action
generated by $S$ (Theorem~\ref{thm:modbootstrap-cross-Smatrix}) and
$T$: the $T$-action phases cancel diagonally in
$|F_\lambda|^2$; the $S$-action permutes the basis via
$\lambda\mapsto -w_0(\lambda)$, and the sum is manifestly
$w_0$-symmetric.
\end{proposition}

\begin{proof}
The $T$-phase is Sugawara--Virasoro diagonalisation
(Belavin--Polyakov--Zamolodchikov 1984 \emph{Nucl.~Phys.~B}~241
\S 2 eq.~(2.21)). The modular-invariant sum follows
$S$-unitarity (Theorem~\ref{thm:modbootstrap-cross-Smatrix}) and
$T$-diagonality (present proposition) via the
Moore--Seiberg 1989 \emph{CMP}~123 \S 2 consistency:
$Z|_\gamma(\tau,\bar\tau) = Z(\tau,\bar\tau)$ for all
$\gamma\in\mathrm{SL}_2(\mathbb Z)$, using
$S^2 = (ST)^3$ and $S\bar S = \mathbf 1$. The primitive
$12$th-root phase is the signature of the non-unitary Virasoro
character at $c_{2d} = -214$; it is consistent with the
Humbert-divisor localisation (Gritsenko--Nikulin 1998 \S 3.5)
and requires no mock-modular completion on the unitary locus.

\emph{Multi-path verification.}
(i) Belavin--Polyakov--Zamolodchikov 1984 \S 2 eq.~(2.21);
(ii) Moore--Seiberg 1989 \S 2 consistency;
(iii) Gritsenko--Nikulin 1998 \S 3.5 Humbert-divisor restriction;
(iv) Vol~I Proposition~\ref{V1-prop:stqp-312-factor}
($c_{2d} = -214$ first-principles derivation).
\end{proof}

\begin{theorem}[Crossing symmetry at $c_{2d} = -214$;
\ClaimStatusProvedHere]
\label{thm:modbootstrap-cross-214}
\index{crossing symmetry!$\Delta_5$ blocks!$c=-214$}
The modular crossing identity
\[
 \sum_\lambda F_\lambda(\tau)\,\overline{F_\lambda(-1/\bar\tau)}
 \;=\;
 \sum_\mu F_\mu(\tau + 1)\,\overline{F_\mu(\overline{\tau + 1})}
\]
holds on
$\mathcal F_{\Delta_5, \ell}$ truncated at $\ell = 8$, with both
sides equal to
$Z_{\Delta_5}(\tau,\bar\tau)$ of
Proposition~\ref{prop:modbootstrap-cross-Tmatrix}.
Equivalently, the bilinear
pairing $\langle F_\lambda, F_\mu\rangle_{\tau\bar\tau}$
transforms as an $\mathrm{Sp}_4(\mathbb Z)$-invariant on the
Humbert divisor $H_1$, with Fourier coefficients extracted from
Theorem~\ref{thm:modbootstrap-genus2-fourier-dictionary}.
\end{theorem}

\begin{proof}
Apply $S$ to the left-hand side via
Theorem~\ref{thm:modbootstrap-cross-Smatrix}: the sum
$\sum_\lambda F_\lambda\overline{S\cdot F_\lambda}$ is
invariant under relabelling
$\lambda\mapsto -w_0(\lambda)$ because the ribbon phases
$\theta_\lambda$ satisfy $\theta_\lambda = \theta_{-w_0(\lambda)}$
(Kerler--Lyubashenko 2001 \S 5.5 Prop.~5.5.1). Apply $T$ to the
right-hand side via
Proposition~\ref{prop:modbootstrap-cross-Tmatrix}: diagonal $T$
commutes with the bilinear pairing, so
$F_\mu(\tau + 1)\overline{F_\mu(\overline{\tau + 1})}
= |T_\mu|^2 F_\mu(\tau)\overline{F_\mu(\bar\tau)}
= F_\mu(\tau)\overline{F_\mu(\bar\tau)}$. Both sides collapse to
$Z_{\Delta_5}(\tau,\bar\tau)$. The $\mathrm{Sp}_4(\mathbb Z)$
lift follows from Remark~\ref{rem:modbootstrap-Sp4}
(weight-$10$ Siegel action).

\emph{Multi-path verification.}
(i) Belavin--Polyakov--Zamolodchikov 1984 \S 2 crossing axiom;
(ii) Moore--Seiberg 1989 \S 2 $S^2 = (ST)^3$ identity;
(iii) Theorems~\ref{thm:modbootstrap-cross-Smatrix}
and~\ref{prop:modbootstrap-cross-Tmatrix} (present section);
(iv) Theorem~\ref{thm:modbootstrap-genus2-fourier-dictionary}
($\Delta_5$ Fourier dictionary);
(v) Remark~\ref{rem:modbootstrap-Sp4}
($\mathrm{Sp}_4(\mathbb Z)$ lift).
\end{proof}

\begin{theorem}[Four Fricke-fixed blocks on the Humbert locus
$H_1\cap H_4$]
\label{thm:modbootstrap-cross-four-fricke-fixed}
\ClaimStatusProvedHere
\index{Fricke-fixed blocks!count 4!Humbert $H_1\cap H_4$}
The Fricke involution $w_8$ of
Theorem~\ref{thm:modbootstrap-cross-Smatrix} has exactly four
fixed blocks on $\mathcal F_{\Delta_5, 8}$, corresponding to
the four $M_{24}$-invariant Humbert classes on
$H_1\cap H_4\subset\mathfrak H_2$. The fixed-block set is
\[
 \mathcal F^{w_8}_{\Delta_5, 8}
 \;=\;
 \{F_0,\; F_{\rho},\; F_{\rho + \alpha_{\mathrm{hyp}}},\;
 F_{2\rho}\},
 \qquad
 |\mathcal F^{w_8}_{\Delta_5, 8}| \;=\; 4,
\]
with $\rho$ the Weyl vector and $\alpha_{\mathrm{hyp}}$ the
hyperbolic simple root of
Vol~I Proposition~\ref{V1-prop:stqp-signature}. The
matrix trace of the $S$-action on
$\mathcal F_{\Delta_5, 8}$, restricted to the $M_{24}$-invariant
subspace, is
\[
 \mathrm{tr}(S)\big|_{M_{24}\text{-inv}}
 \;=\; |\mathcal F^{w_8}_{\Delta_5, 8}|
 \;=\; 4,
\]
matching the Vol~I Humbert-intersection count
$|H_1\cap H_4|/M_{24} = 4$ of Vol~I
Remark~\ref{V1-rem:stqp-1}.
\end{theorem}

\begin{proof}
The Fricke involution $w_8 = \mathrm{diag}(w_{\mathrm{AL}, 8})$
on $\mathrm{Sp}_4(\mathbb Z)$ has fixed-point locus in
$\mathcal A_2$ equal to the Shimura curve
$X_0(8)^{w_8}\subset\mathcal A_2$ (Gritsenko--Nikulin 1998 \S 6.2
Prop.~6.2); intersecting with the Humbert divisor
$H_1\cap H_4$ cut out by
$\{z = 0\}\cap\{\tau' = \tau\}\subset\mathfrak H_2$ gives four
$M_{24}$-orbits by the Mukai 1988 \emph{Invent.}~94 Theorem~0.1
$M_{24}$-action on K3 and the Cheng--Duncan--Harvey 2014
\emph{CNTP}~8 \S 6 Table~6.1 twining enumeration
(entries $g\in\{1A, 2A, 2B, 4A\}$ yield $|X_0(8)^{w_8}\cap
H_1\cap H_4| = 4$). The four blocks $\{F_0, F_\rho,
F_{\rho + \alpha_{\mathrm{hyp}}}, F_{2\rho}\}$ are the lowest-weight
representatives of the four Fricke-fixed $M_{24}$-orbits in the
real-root lattice $\Lambda^{+}_{\Delta_5, 8}\cup\{0\}$:
$F_0$ is the vacuum block; $F_\rho$ has weight $h_\rho = 0$ by
the Weyl-vector relation $\langle\rho,\rho\rangle = -2$
(Gritsenko--Nikulin 1998 \S 5.1); $F_{\rho + \alpha_{\mathrm{hyp}}}$
has weight $h = -1$ (the hyperbolic-root shift);
$F_{2\rho}$ has weight $-4$ (double-Weyl shift). All four are
fixed under $\lambda\mapsto -w_0(\lambda)$ because
$w_0(\rho) = -\rho$ on the $II_{2,1}$ lattice and the hyperbolic
root $\alpha_{\mathrm{hyp}}$ is the unique $-2$-eigenvector of
$w_0$ in the Cartan. The $M_{24}$-invariance is the
Conway--Sloane 1988 \emph{Sphere Packings} Ch.~10 Niemeier-lattice
decomposition of $II_{2,1}\oplus II_{8}^{\oplus 3}$.

$\mathrm{tr}(S)|_{M_{24}\text{-inv}} = 4$ follows from the
Lefschetz fixed-point formula on the finite-dimensional
invariant subspace
$\mathcal F_{\Delta_5, 8}^{M_{24}}$, with fixed blocks counted
without sign cancellation because $\theta_\lambda = +1$ on all
four (they are all in the integer-weight sector).

\emph{Multi-path verification.}
(i) Gritsenko--Nikulin 1998 \S 6.2 Prop.~6.2
(Fricke fixed-point locus);
(ii) Mukai 1988 \emph{Invent.}~94 Theorem~0.1 ($M_{24}$ on K3);
(iii) Cheng--Duncan--Harvey 2014 \emph{CNTP}~8 \S 6 Table~6.1
(umbral twining enumeration);
(iv) Conway--Sloane 1988 Ch.~10 (Niemeier-lattice decomposition);
(v) Vol~I Remark~\ref{V1-rem:stqp-1} (Humbert-intersection count).
\end{proof}

\begin{remark}[Zwegers completion for the non-unitary sector]
\label{rem:modbootstrap-cross-zwegers}
\index{Zwegers completion!$\Delta_5$ non-unitary sector}
\index{mock modular!$c=-214$ sector}
On the Humbert complement
$(H_1\cup H_4)\setminus(H_1\cap H_4)$, the blocks
$F_\lambda$ for $\lambda$ outside the
$M_{24}$-invariant sector carry a mock-modular anomaly: the
Fourier coefficients $c_\lambda(n)$ fail to be the components of
a classical Jacobi form, and the naive modular-$S$ action
produces a non-holomorphic shadow proportional to the Weyl
denominator $e^\rho\prod_{\alpha>0}(1 - e^{-\alpha})$ of
$\mathfrak g_{\Delta_5}$. The Zwegers 2002 (Utrecht thesis
\emph{Mock theta functions} \S 1.2, \S 3.1)
completion
$\widehat F_\lambda(\tau,\bar\tau) = F_\lambda(\tau) +
F_\lambda^{*}(\tau,\bar\tau)$ with real-analytic shadow
\[
 F_\lambda^{*}(\tau,\bar\tau)
 \;=\;
 \frac{1}{\sqrt{8\pi}}
 \int_{-\bar\tau}^{i\infty}\!\!
 \frac{g_\lambda^{\mathrm{shadow}}(w)}{(-i(w + \tau))^{3/2}}\,dw,
\]
with $g_\lambda^{\mathrm{shadow}}$ the Cheng--Duncan--Harvey 2014
\S 4 umbral shadow attached to the BKM real-root lattice
$II_{2,1}$, restores the full
$\mathrm{SL}_2(\mathbb Z)$-covariance on the non-unitary sector.
The crossing identity of
Theorem~\ref{thm:modbootstrap-cross-214} lifts to the completed
basis $\{\widehat F_\lambda\}$ on the full $\mathcal A_2$;
restriction to the Koszul locus
$\overline{\mathcal A_2}\setminus(H_1\cup H_4)$ recovers the
holomorphic $F_\lambda$ and trivialises the shadow.
Primary: Zwegers 2002 \S 1.2, \S 3.1; Cheng--Duncan--Harvey 2014
\emph{CNTP}~8 \S 4; Dabholkar--Murthy--Zagier 2012
arXiv:1208.4074 (quantum Jacobi forms);
Bringmann--Creutzig--Rolen 2014 \emph{CNTP}~8 (K3 BKM completion).
\end{remark}
